% Les facteurs de la mondialisation — Géographie Tle E — Cours signé OnBuch+
% Programme officiel MINESEC. Compiler avec : tectonic N11-facteurs-mondialisation.tex
\newcommand{\DOCMATIERE}{Géographie}
\newcommand{\DOCNIVEAU}{Tle E}
\newcommand{\DOCMODULE}{Module 2 : La libéralisation des échanges (la mondialisation)}
\newcommand{\DOCLECON}{N11}
\newcommand{\DOCTITRE}{Les facteurs de la mondialisation}
% OnBuch+ — préambule « cours signés » ENRICHI (compilé avec tectonic / XeLaTeX)
% Système visuel « Pop » d'OnBuch+ : fond crème (jamais blanc), bordures encre
% épaisses, ombres « dures » décalées, titres Archivo Black en capitales.
% Version MATHÉMATIQUES : graphes (pgfplots/tikz), boîtes pédagogiques
% (définition, propriété, méthode, attention, le savais-tu, exercice résolu,
% exercice, TP, bilan), page de garde et signature OnBuch+ ancrée en bas.
%
% Macros attendues dans main.tex AVANT \input{preamble} :
%   \DOCMATIERE  \DOCNIVEAU  \DOCMODULE  \DOCLECON (numéro)  \DOCTITRE
\documentclass[11pt]{article}

\usepackage{fontspec}
\usepackage[french]{babel}
\usepackage[a4paper,margin=2cm,top=3.4cm,bottom=2.8cm,headheight=1.4cm,footskip=1.3cm]{geometry}
\usepackage{amsmath,amssymb,amsfonts}
\usepackage{mathtools} % \xmapsto, \coloneqq... souvent utilisés par les modèles
\usepackage{cancel} % simplifications barrées (bilans de forces)
% \abs{z} (module, valeur absolue) et \norm{u} (norme), utilisés par les modèles
\providecommand{\abs}{}\let\abs\relax\DeclarePairedDelimiter{\abs}{\lvert}{\rvert}
\providecommand{\norm}{}\let\norm\relax\DeclarePairedDelimiter{\norm}{\lVert}{\rVert}
\usepackage[table]{xcolor} % [table] : fournit aussi \rowcolors (alternance de lignes)
\usepackage{pagecolor}
\usepackage{tikz}
\usetikzlibrary{arrows.meta,positioning,calc,shapes.geometric,shapes.misc,shapes.arrows,decorations.pathmorphing,decorations.pathreplacing,decorations.markings,patterns,shadows,mindmap,trees,fit,backgrounds,matrix,intersections,angles,quotes}
\usepackage{pgfplots}
\usepackage[version=4]{mhchem} % équations nucléaires \ce{^{235}_{92}U}
\usepackage[european,siunitx]{circuitikz} % schémas électriques
\pgfplotsset{compat=1.18}
\usepgfplotslibrary{fillbetween,groupplots} % aires entre courbes (calcul intégral)
\usepackage{enumitem}
\usepackage{fancyhdr}
% [most] charge toutes les bibliothèques tcolorbox (listings, minted, raster,
% documentation...) et gonfle inutilement la mémoire TeX sur un document long
% (~300 boîtes) : on ne charge que ce qui est réellement utilisé.
\usepackage[skins,breakable]{tcolorbox}
\usepackage{graphicx}
\graphicspath{{/home/user/t-t/geo-onbuch/images/}}
\usepackage{colortbl}
\usepackage{booktabs}
\usepackage{tabularx}
\usepackage{multirow}
\usepackage{diagbox} % case d'en-tête diagonale des tableaux à double entrée
% Colonnes extensibles centrée / gauche / droite (C, L, R), souvent utilisées par les modèles
\newcolumntype{C}{>{\centering\arraybackslash}X}
\newcolumntype{L}{>{\raggedright\arraybackslash}X}
\newcolumntype{R}{>{\raggedleft\arraybackslash}X}
\usepackage{array}
\usepackage{multicol}
\usepackage{siunitx}
% parse-numbers=false : accepte aussi \SI{2,5 \times 10^{-3}}{...}
\sisetup{locale=FR,output-decimal-marker={,},group-minimum-digits=4,parse-numbers=false}
\DeclareSIUnit\ohm{\ensuremath{\symup{\Omega}}} % U+2126 absent de la police
\DeclareSIUnit\division{div} % réglages d'oscilloscope (ms/div, V/div)
\DeclareSIUnit\tr{tr} % tours (vitesse de rotation, ex. \SI{1500}{\tr\per\minute})
\DeclareSIUnit\molar{M} % concentration molaire (ex. \SI{3}{\molar}, \SI{1}{\micro\molar})
\DeclareSIUnit\an{an} % année (le modèle confond parfois avec \year, un compteur TeX) — ex. \SI{5}{\kilo\meter\per\an}
\DeclareSIUnit\FCFA{FCFA} % franc CFA (ex. \SI{500}{\FCFA\per\kilo\gram})
\DeclareSIUnit\degreeBrix{\textdegree Brix} % teneur en sucre d'un sirop/jus (réfractométrie)
\DeclareSIUnit\month{mois} % le modèle utilise parfois \month, un compteur TeX, comme unité de durée
\DeclareSIUnit\microliter{\micro\liter} % le modèle écrit parfois \microliter en un seul token
\DeclareSIUnit\million{millions} % dénombrement (ex. \SI{15}{\million\per\mL} spermatozoïdes)
\usepackage{qrcode}
% \degree : commande fréquemment utilisée par les modèles pour un angle ou une
% température en dehors de \SI{}{} (non fournie par les paquets chargés).
\providecommand{\degree}{^{\circ}}
% \qed (fin de démonstration), utilisé par les modèles sans amsthm
\providecommand{\qed}{\hfill$\square$}
% \barre : double barre des valeurs interdites dans un tableau de variations (tkz-tab)
\providecommand{\barre}{\ensuremath{\|}}
% environnement proof (amsthm non chargé) utilisé par les modèles
\ifdefined\proof\else\newenvironment{proof}[1][Démonstration]{\par\noindent\textit{#1.}\ }{\qed\par}\fi
\usepackage{lastpage}
\usepackage{hyperref}
\hypersetup{hidelinks}

% ============================================================
% Palette « Pop » OnBuch+ — mêmes hex que la classe P (plus_ui.dart)
% ============================================================
\definecolor{poporange}{HTML}{F59321}
\definecolor{popdark}{HTML}{E07A0C}
\definecolor{popcream}{HTML}{FFF6E8}
\definecolor{popink}{HTML}{1C1714}
\definecolor{poppurple}{HTML}{7A5AE0}
\definecolor{popgreen}{HTML}{2E9E6A}
\definecolor{poppink}{HTML}{E0457B}
\definecolor{popblue}{HTML}{2F7DE1}
\definecolor{popgold}{HTML}{C98A12}
\definecolor{popmuted}{HTML}{8A7F76}
\definecolor{popline}{HTML}{EADFD2}
\definecolor{popgray}{HTML}{8A7F76}
\colorlet{popgrey}{popgray}
% Alias tolérants (le modèle nomme parfois les couleurs différemment)
\colorlet{popwhite}{white}
\colorlet{popblack}{popink}
\colorlet{popred}{poppink}
\colorlet{popyellow}{popgold}
% Teintes claires (fonds de tableaux/encadrés) — le modèle nomme parfois
% les couleurs avec un suffixe L (ex. popblueL) sans les avoir définies.
\colorlet{poporangeL}{poporange!20}
\colorlet{popdarkL}{popdark!20}
\colorlet{poppurpleL}{poppurple!20}
\colorlet{popgreenL}{popgreen!20}
\colorlet{poppinkL}{poppink!20}
\colorlet{popblueL}{popblue!20}
\colorlet{popgoldL}{popgold!20}
\colorlet{popredL}{popred!20}
\colorlet{popgrayL}{popgray!20}
% Teintes claires pour fonds de tableaux / zones
\colorlet{popblueL}{popblue!10!white}
\colorlet{poporangeL}{poporange!14!white}
\colorlet{popgreenL}{popgreen!12!white}
\colorlet{poppurpleL}{poppurple!12!white}
\colorlet{poppinkL}{poppink!12!white}
\colorlet{popgoldL}{popgold!14!white}

\pagecolor{popcream}

% ============================================================
% Typographie
% ============================================================
\defaultfontfeatures{Ligatures=TeX}
\setmainfont{PlusJakartaSans-Medium.ttf}[Path=../../../fonts/,BoldFont=PlusJakartaSans-Bold.ttf]
% Maths en sans-serif (Fira Math) avec chiffres et lettres droites de Plus
% Jakarta, pour que les formules se fondent dans le texte.
\usepackage[math-style=ISO,bold-style=ISO]{unicode-math}
\setmathfont{FiraMath-Regular.otf}[Path=../../../fonts/]
\setmathfont{PlusJakartaSans-Medium.ttf}[Path=../../../fonts/,range={up/num,up/{Latin,latin},\mathup}]
\usepackage{icomma} % virgule décimale sans espace en mode math
% Caractères mathématiques Unicode absents des polices de texte (Plus Jakarta,
% Archivo Black) : rendus par leur équivalent mathématique, en texte comme en
% formule — sinon ils s'affichent en carré vide (ex. « Arithmétique dans ℤ »).
\usepackage{newunicodechar}
\newunicodechar{ℝ}{\ensuremath{\mathbb{R}}}
\newunicodechar{ℤ}{\ensuremath{\mathbb{Z}}}
\newunicodechar{ℂ}{\ensuremath{\mathbb{C}}}
\newunicodechar{ℕ}{\ensuremath{\mathbb{N}}}
\newunicodechar{ℚ}{\ensuremath{\mathbb{Q}}}
\newunicodechar{𝔻}{\ensuremath{\mathbb{D}}}
\newunicodechar{α}{\ensuremath{\alpha}}
\newunicodechar{β}{\ensuremath{\beta}}
\newunicodechar{γ}{\ensuremath{\gamma}}
\newunicodechar{δ}{\ensuremath{\delta}}
\newunicodechar{ε}{\ensuremath{\varepsilon}}
\newunicodechar{θ}{\ensuremath{\theta}}
\newunicodechar{λ}{\ensuremath{\lambda}}
\newunicodechar{μ}{\ensuremath{\mu}}
\newunicodechar{σ}{\ensuremath{\sigma}}
\newunicodechar{τ}{\ensuremath{\tau}}
\newunicodechar{φ}{\ensuremath{\varphi}}
\newunicodechar{ψ}{\ensuremath{\psi}}
\newunicodechar{ω}{\ensuremath{\omega}}
\newunicodechar{Σ}{\ensuremath{\Sigma}}
% majuscules produites par \MakeUppercase dans les titres (α-aminés -> Α)
\newunicodechar{Α}{\ensuremath{\alpha}}
\newunicodechar{Β}{\ensuremath{\beta}}
\newunicodechar{Γ}{\ensuremath{\Gamma}}
\newunicodechar{Δ}{\ensuremath{\Delta}}
\newunicodechar{Ε}{\ensuremath{\varepsilon}}
\newunicodechar{Θ}{\ensuremath{\Theta}}
\newunicodechar{Λ}{\ensuremath{\Lambda}}
\newunicodechar{Μ}{\ensuremath{\mu}}
\newunicodechar{Τ}{\ensuremath{\tau}}
\newunicodechar{Φ}{\ensuremath{\Phi}}
\newunicodechar{Ψ}{\ensuremath{\Psi}}
\newunicodechar{Ω}{\ensuremath{\Omega}}
\newunicodechar{↦}{\ensuremath{\mapsto}}
\newunicodechar{∈}{\ensuremath{\in}}
\newunicodechar{∉}{\ensuremath{\notin}}
\newunicodechar{∧}{\ensuremath{\wedge}}
\newunicodechar{⇔}{\ensuremath{\Leftrightarrow}}
\newunicodechar{⟺}{\ensuremath{\Longleftrightarrow}}
\newunicodechar{⇒}{\ensuremath{\Rightarrow}}
\newunicodechar{⟹}{\ensuremath{\Longrightarrow}}
\newunicodechar{∘}{\ensuremath{\circ}}
\newunicodechar{∪}{\ensuremath{\cup}}
\newunicodechar{∩}{\ensuremath{\cap}}
\newunicodechar{≡}{\ensuremath{\equiv}}
\newunicodechar{⊂}{\ensuremath{\subset}}
\newunicodechar{⊥}{\ensuremath{\perp}}
\newunicodechar{∃}{\ensuremath{\exists}}
\newunicodechar{∀}{\ensuremath{\forall}}
\newunicodechar{∛}{\ensuremath{\sqrt[3]{\,}}}
\newunicodechar{∜}{\ensuremath{\sqrt[4]{\,}}}
\newunicodechar{⃗}{\ensuremath{{}^{\to}}}
\newunicodechar{ⁿ}{\ifmmode^{n}\else\textsuperscript{\ensuremath{n}}\fi}
\newunicodechar{ˣ}{\ifmmode^{x}\else\textsuperscript{\ensuremath{x}}\fi}
\newunicodechar{ᵇ}{\ifmmode^{b}\else\textsuperscript{\ensuremath{b}}\fi}
\newunicodechar{ʳ}{\ifmmode^{r}\else\textsuperscript{\ensuremath{r}}\fi}
\newunicodechar{ᵘ}{\ifmmode^{u}\else\textsuperscript{\ensuremath{u}}\fi}
\newunicodechar{ʸ}{\ifmmode^{y}\else\textsuperscript{\ensuremath{y}}\fi}
\newunicodechar{ᵃ}{\ifmmode^{a}\else\textsuperscript{\ensuremath{a}}\fi}
\newunicodechar{ᵉ}{\ifmmode^{e}\else\textsuperscript{\ensuremath{e}}\fi}
\newunicodechar{ᵏ}{\ifmmode^{k}\else\textsuperscript{\ensuremath{k}}\fi}
\newunicodechar{ᵖ}{\ifmmode^{p}\else\textsuperscript{\ensuremath{p}}\fi}
\newunicodechar{ᴺ}{\ifmmode^{N}\else\textsuperscript{\ensuremath{N}}\fi}
\newunicodechar{ᶜ}{\ifmmode^{c}\else\textsuperscript{\ensuremath{c}}\fi}
\newunicodechar{ʰ}{\ifmmode^{h}\else\textsuperscript{\ensuremath{h}}\fi}
\newunicodechar{ᵠ}{\ifmmode^{q}\else\textsuperscript{\ensuremath{q}}\fi}
\newunicodechar{ₙ}{\ifmmode_{n}\else\textsubscript{\ensuremath{n}}\fi}
\newunicodechar{ₐ}{\ifmmode_{a}\else\textsubscript{\ensuremath{a}}\fi}
\newunicodechar{ᵢ}{\ifmmode_{i}\else\textsubscript{\ensuremath{i}}\fi}
\newunicodechar{ᵣ}{\ifmmode_{r}\else\textsubscript{\ensuremath{r}}\fi}
\newunicodechar{ₖ}{\ifmmode_{k}\else\textsubscript{\ensuremath{k}}\fi}
\newunicodechar{ᵦ}{\ifmmode_{\beta}\else\textsubscript{\ensuremath{\beta}}\fi}
\newunicodechar{ₓ}{\ifmmode_{x}\else\textsubscript{\ensuremath{x}}\fi}
\newfontfamily\popdisplay{ArchivoBlack-Regular.ttf}[Path=../../../fonts/]
\newfontfamily\poplogo{SpaceGrotesk-Bold.ttf}[Path=../../../fonts/]
\newfontfamily\popsemi{PlusJakartaSans-SemiBold.ttf}[Path=../../../fonts/]
\newfontfamily\popxbold{PlusJakartaSans-ExtraBold.ttf}[Path=../../../fonts/]
\newfontfamily\popxboldls{PlusJakartaSans-ExtraBold.ttf}[Path=../../../fonts/,LetterSpace=3.0]

\setlist[enumerate]{leftmargin=1.7em,itemsep=5pt,topsep=4pt}
\setlist[itemize]{leftmargin=1.7em,itemsep=4pt,topsep=4pt,label={\color{poporange}\textbullet}}
\setlength{\parindent}{0pt}
\setlength{\parskip}{5pt}
\renewcommand{\arraystretch}{1.25}

% pgfplots : style Pop par défaut
\pgfplotsset{
  every axis/.append style={axis line style={popink,line width=1pt},tick style={popink},
    grid style={popline,line width=0.5pt},label style={font=\small},tick label style={font=\footnotesize}},
  popcurve/.style={poporange,line width=1.8pt,no markers,smooth},
}
% Même style disponible dans un \draw TikZ simple (hors pgfplots)
\tikzset{popcurve/.style={poporange,line width=1.8pt}}
\tikzset{popgrid/.style={popline,very thin}}
\colorlet{popcurve}{poporange} % le nom du style est parfois utilisé comme couleur

% ============================================================
% Pill / squiggle
% ============================================================
\newcommand{\poppill}[3]{%
  \tikz[baseline=-0.55ex]\node[draw=popink,line width=1.2pt,rounded corners=2.6pt,%
    fill=#2,inner xsep=6.5pt,inner ysep=3pt]{%
    \popxboldls\fontsize{7}{7}\selectfont\color{#3}\MakeUppercase{#1}};%
}
\newcommand{\popsquiggle}[1]{%
  \tikz[baseline=-0.4ex]{
    \draw[#1,line width=2.6pt,line cap=round]
      plot [smooth] coordinates {(0,0.09) (0.34,0.01) (0.68,0.21) (1.02,0.01) (1.36,0.10)};
  }%
}
% Mot-clé surligné (marqueur orange sous le texte)
\newcommand{\cle}[1]{{\popxbold\color{popdark}#1}}

% ============================================================
% En-tête / pied
% ============================================================
\pagestyle{fancy}
\fancyhf{}
\renewcommand{\headrulewidth}{0pt}
\renewcommand{\footrulewidth}{0pt}
\lhead{\raisebox{-0.15ex}{\poplogo\fontsize{13}{13}\selectfont\color{popink}OnBuch\color{poporange}+}}
\rhead{\poppill{\DOCMATIERE\ \textbullet\ \DOCNIVEAU\ \textbullet\ Leçon \DOCLECON}{white}{popink}}
\lfoot{\popxbold\fontsize{7.6}{7.6}\selectfont\color{popmuted}\MakeUppercase{Cours signé OnBuch+}\ \textbullet\ {\popsemi\DOCTITRE}}
\rfoot{\tikz[baseline=-0.6ex]\node[rounded corners=8.5pt,fill=popink,minimum height=17pt,minimum width=17pt,inner xsep=5pt,inner sep=0]{\popxbold\fontsize{8}{8}\selectfont\color{popcream}\thepage\,/\,\pageref*{LastPage}};}

% ============================================================
% Titres
% ============================================================
\newcommand{\chaptitre}[2]{%
  \begin{tcolorbox}[enhanced,colback=poporange,colframe=popink,boxrule=2.6pt,arc=11pt,
      drop shadow=popink,left=15pt,right=15pt,top=12pt,bottom=11pt,before skip=3pt,after skip=18pt]
    \poppill{#2}{popink}{popcream}\par\vspace{9pt}
    {\raggedright\hyphenpenalty=10000\exhyphenpenalty=10000\popdisplay\fontsize{22}{24}\selectfont\color{popcream}\MakeUppercase{#1}\par}
  \end{tcolorbox}
}
% Section numérotée : pastille orange avec le numéro + titre Archivo
\newcounter{popsec}
\newcommand{\coursec}[1]{%
  \stepcounter{popsec}\setcounter{popsub}{0}%
  \par\vspace{16pt}\noindent
  \begin{minipage}{\linewidth}
  \tikz[baseline=-0.7ex]\node[draw=popink,line width=1.6pt,fill=poporange,rounded corners=4pt,minimum size=22pt,inner sep=0]{\popdisplay\fontsize{12}{12}\selectfont\color{popcream}\thepopsec};%
  \hspace{8pt}{\popdisplay\fontsize{15}{17}\selectfont\color{popink}\MakeUppercase{#1}}\par
  \vspace{4pt}\noindent{\color{poporange}\rule{36pt}{3.6pt}}
  \end{minipage}\par\nopagebreak\vspace{8pt}
}
\newcounter{popsub}
\newcommand{\courssub}[1]{%
  \stepcounter{popsub}%
  \par\vspace{9pt}\noindent{\popxbold\fontsize{11.5}{13}\selectfont\color{popink}{\color{poporange}\thepopsec.\thepopsub}\hspace{6pt}#1}\par\nopagebreak\vspace{4pt}
}

% ============================================================
% Boîtes pédagogiques — même recette : fond blanc, bordure épaisse colorée,
% ombre dure encre, badge plein. Titre optionnel [#1].
% ============================================================
\tcbset{popbase/.style={breakable,enhanced,colback=white,boxrule=2.3pt,arc=9pt,
  shadow={3pt}{-3pt}{0pt}{popink},left=14pt,right=14pt,top=11pt,bottom=11pt,before skip=11pt,after skip=15pt,
  fontupper=\popsemi\fontsize{10.6}{14.5}\selectfont\color{popink},
  fontlower=\popsemi\fontsize{10.6}{14.5}\selectfont\color{popink}}}
% Coche dessinée (le glyphe ✓ n'existe pas dans les polices du document)
\newcommand{\popcheck}[1]{\tikz[baseline=-0.5ex]\draw[#1,line width=1.6pt,line cap=round,line join=round](-0.13,0.0)--(-0.04,-0.09)--(0.13,0.1);}
\providecommand{\coche}{\popcheck{popgreen}}
\providecommand{\resultat}[1]{\par\smallskip\noindent\fcolorbox{popgreen}{popgreenL}{\parbox{\dimexpr\linewidth-2\fboxsep-2\fboxrule\relax}{\textbf{Résultat.} #1}}\par\smallskip}
\newcommand{\popboxhead}[3]{\poppill{#1}{#2}{white}\ifx\relax#3\relax\else\hspace{6pt}{\popxbold\fontsize{10.6}{12}\selectfont\color{popink}#3}\fi\par\vspace{7pt}}

\newtcolorbox{aretenir}[1][]{popbase,colframe=popgreen,before upper={\popboxhead{À retenir}{popgreen}{#1}}}
\newtcolorbox{exemplebox}[1][]{popbase,colframe=poporange,before upper={\popboxhead{Exemple}{poporange}{#1}}}
\newtcolorbox{definition}[1][]{popbase,colframe=popblue,before upper={\popboxhead{Définition}{popblue}{#1}}}
\newtcolorbox{propriete}[1][]{popbase,colframe=poppurple,before upper={\popboxhead{Propriété}{poppurple}{#1}}}
\newtcolorbox{methode}[1][]{popbase,colframe=popgold,before upper={\popboxhead{Méthode}{popgold}{#1}}}
\newtcolorbox{attention}[1][]{popbase,colframe=poppink,before upper={\popboxhead{Attention}{poppink}{#1}}}
\newtcolorbox{experience}[1][]{popbase,colframe=popblue,colback=popblueL,before upper={\popboxhead{Expérience}{popblue}{#1}}}
% « Le savais-tu ? » : carte violette pleine (contexte, culture, Cameroun)
\newtcolorbox{savaistu}[1][]{popbase,colback=poppurple,colframe=popink,
  fontupper=\popsemi\fontsize{10.4}{14.2}\selectfont\color{white},
  before upper={\poppill{Le savais-tu\,?}{white}{poppurple}\ifx\relax#1\relax\else\hspace{6pt}{\popxbold\color{white}#1}\fi\par\vspace{7pt}}}
% Exercice résolu : énoncé en haut, \tcblower puis la solution sur fond crème
\newcounter{popexo}\newcounter{popexr}
\newtcolorbox{exoresolu}[1][]{popbase,colframe=poporange,colbacklower=popcream,
  segmentation style={popink,line width=1.2pt,dashed},
  before upper={\stepcounter{popexr}\popboxhead{Exercice résolu \thepopexr}{poporange}{#1}},
  before lower={\poppill{Solution}{popink}{popcream}\par\vspace{6pt}}}
% Exercice (non résolu en place) — la correction est dans la partie Corrigés
\newtcolorbox{exercice}[1][]{popbase,colframe=popink,
  before upper={\stepcounter{popexo}\popboxhead{Exercice \thepopexo}{popink}{#1}}}
% Corrigé
\newtcolorbox{corrige}[1][]{popbase,colframe=popgreen,colback=popgreenL,
  before upper={\popboxhead{Corrigé}{popgreen}{#1}}}
% Objectifs / prérequis / bilan
\newtcolorbox{objectifs}{popbase,colframe=popink,colback=poporangeL,
  before upper={\popboxhead{Objectifs de la leçon}{popink}{}}}
\newtcolorbox{prerequis}{popbase,colframe=popink,colback=popblueL,
  before upper={\popboxhead{Prérequis}{popblue}{}}}
\newtcolorbox{bilan}{popbase,colframe=popink,colback=white,boxrule=2.8pt,
  before upper={\popboxhead{Fiche bilan}{poporange}{L'essentiel à connaître pour l'examen}}}
% Figure encadrée avec légende
\newtcolorbox{popfigure}[1][]{enhanced,breakable=false,colback=white,colframe=popline,boxrule=1.4pt,arc=8pt,
  left=10pt,right=10pt,top=10pt,bottom=8pt,before skip=10pt,after skip=12pt,halign=center,
  lower separated=false,halign lower=center,
  fontlower=\popsemi\fontsize{9}{11.5}\selectfont\color{popmuted},
  #1}
\newcommand{\legende}[1]{\par\vspace{6pt}{\popsemi\fontsize{9}{11.5}\selectfont\color{popmuted}#1}}

% ============================================================
% Signature OnBuch+
% ============================================================
\newcommand{\poplogoinline}[1]{{\poplogo\fontsize{#1}{#1}\selectfont\color{popink}OnBuch\color{poporange}+}}
\newcommand{\signatureonbuch}{%
  \begin{tcolorbox}[enhanced,colback=popink,colframe=popink,boxrule=2.2pt,arc=10pt,
      drop shadow=poporange,left=14pt,right=14pt,top=10pt,bottom=10pt,before skip=0pt,after skip=0pt]
    \tikz[baseline=-0.7ex]\node[circle,fill=popgreen,draw=popcream,line width=1.3pt,minimum size=22pt,inner sep=0]{\popcheck{white}};%
    \hspace{9pt}\begin{minipage}[c]{0.62\linewidth}
      {\popxbold\fontsize{12}{13}\selectfont\color{popcream}Signé\ \textbullet\ {\poplogo OnBuch\color{poporange}+}}\par\vspace{2pt}
      {\raggedright\popsemi\fontsize{8}{10.5}\selectfont\color{popline}\MakeUppercase{Équipe pédagogique OnBuch+}\\\MakeUppercase{Conforme au programme officiel MINESEC}\par}
    \end{minipage}\hfill
    {\poplogo\fontsize{22}{22}\selectfont\color{popcream}OnBuch\color{poporange}+}
  \end{tcolorbox}%
}

% ============================================================
% Page de garde
% #1 titre · #2 matière · #3 niveau · #4 module · #5 n° leçon · #6 durée ·
% #7 description / objectifs (texte ou itemize)
% ============================================================
\newcommand{\pagegarde}[7]{%
  \thispagestyle{empty}
  \newgeometry{a4paper,margin=1.7cm,top=1.6cm,bottom=1.4cm}%
  \begin{tikzpicture}[remember picture,overlay]
    \fill[poporange!18] ([xshift=-3.2cm,yshift=-3.6cm]current page.north east) circle (4.6cm);
    \fill[poppurple!14] ([xshift=2.4cm,yshift=7.6cm]current page.south west) circle (3.8cm);
  \end{tikzpicture}%
  {\poplogo\fontsize{30}{30}\selectfont\color{popink}OnBuch\color{poporange}+}\hfill
  \raisebox{0.6ex}{\poppill{Cours signé \textbullet\ Premium}{popink}{popcream}}\par
  \vspace{1pt}\popsquiggle{poppurple}\par
  \vspace{8pt}
  \begin{tcolorbox}[enhanced,colback=poporange,colframe=popink,boxrule=2.8pt,arc=13pt,
      drop shadow=popink,left=18pt,right=18pt,top=12pt,bottom=13pt,after skip=10pt]
    \poppill{#2 \textbullet\ #3}{popink}{popcream}\hspace{5pt}\poppill{#4}{white}{popink}\par\vspace{8pt}
    {\popdisplay\fontsize{14}{14}\selectfont\color{popink}LEÇON #5}\par\vspace{4pt}
    {\raggedright\hyphenpenalty=10000\exhyphenpenalty=10000\popdisplay\fontsize{25}{27}\selectfont\color{popcream}\MakeUppercase{#1}\par}
    \vspace{8pt}
    {\popxbold\fontsize{9.5}{9.5}\selectfont\color{popink}\MakeUppercase{Durée conseillée : #6}}
  \end{tcolorbox}
  \begin{tcolorbox}[enhanced,colback=white,colframe=popink,boxrule=2.2pt,arc=10pt,
      drop shadow=popink,left=16pt,right=16pt,top=10pt,bottom=10pt,after skip=8pt]
    \poppill{Dans cette leçon}{poppurple}{white}\par\vspace{5pt}
    \popsemi\fontsize{9.3}{12.6}\selectfont\color{popink}#7
  \end{tcolorbox}
  \noindent\begin{minipage}[t]{0.487\linewidth}
    \begin{tcolorbox}[enhanced,colback=popgreen,colframe=popink,boxrule=2.2pt,arc=10pt,
        drop shadow=popink,left=11pt,right=11pt,top=10pt,bottom=10pt,height=3.1cm,valign=center]
      \tcbox[colback=white,colframe=popink,boxrule=1.3pt,arc=5pt,boxsep=0pt,nobeforeafter,
        left=3pt,right=3pt,top=3pt,bottom=3pt,valign=center]{\qrcode[height=1.9cm]{https://play.google.com/store/apps/details?id=com.luvvix.onbuch&hl=fr}}%
      \hspace{8pt}\begin{minipage}[c]{0.5\linewidth}
        {\popdisplay\fontsize{11}{12.5}\selectfont\color{white}\MakeUppercase{Télécharge OnBuch}}\par\vspace{4pt}
        {\raggedright\popsemi\fontsize{8.4}{11}\selectfont\color{white}Gratuit \textbullet\ Google Play\\Tuteur IA, annales, concours, résultats\par}
      \end{minipage}
    \end{tcolorbox}
  \end{minipage}\hfill
  \begin{minipage}[t]{0.487\linewidth}
    \begin{tcolorbox}[enhanced,colback=poppurple,colframe=popink,boxrule=2.2pt,arc=10pt,
        drop shadow=popink,left=11pt,right=11pt,top=10pt,bottom=10pt,height=3.1cm,valign=center]
      \tikz[baseline=-0.7ex]\node[circle,fill=white,draw=popink,line width=1.3pt,minimum size=1.6cm,inner sep=0]{\popdisplay\fontsize{24}{24}\selectfont\color{poppurple}+};%
      \hspace{8pt}\begin{minipage}[c]{0.6\linewidth}
        {\popdisplay\fontsize{11}{12.5}\selectfont\color{white}\MakeUppercase{Passe à OnBuch+}}\par\vspace{4pt}
        {\raggedright\popsemi\fontsize{8.4}{11}\selectfont\color{white}Tous les cours signés, Tuteur IA illimité, fiches PDF — onglet OnBuch+ de l'app\par}
      \end{minipage}
    \end{tcolorbox}
  \end{minipage}
  \vspace{8pt}
  \popcontenu
  \vfill
  \signatureonbuch
  \restoregeometry
  \newpage
}

% Rangée « Ce cours contient » (page de garde)
\newcommand{\poptile}[3]{% #1 couleur #2 gros texte #3 légende
  \begin{tcolorbox}[enhanced,colback=#1,colframe=popink,boxrule=2pt,arc=9pt,shadow={2.5pt}{-2.5pt}{0pt}{popink},
      left=6pt,right=6pt,top=9pt,bottom=9pt,halign=center,valign=center,height=1.75cm,width=\linewidth,nobeforeafter]
    {\popdisplay\fontsize{12.5}{13}\selectfont\color{white}\MakeUppercase{#2}}\par\vspace{3pt}
    {\centering\popsemi\fontsize{7.8}{9.5}\selectfont\color{white}#3\par}
  \end{tcolorbox}}
\newcommand{\popcontenu}{%
  \par\noindent{\popxboldls\fontsize{7.5}{7.5}\selectfont\color{popmuted}CE COURS SIGNÉ CONTIENT}\par\vspace{7pt}
  \noindent\begin{minipage}[t]{0.235\linewidth}\poptile{popblue}{Cours}{complet, illustré, pas à pas}\end{minipage}\hfill
  \begin{minipage}[t]{0.235\linewidth}\poptile{poporange}{Méthodes}{et exercices résolus}\end{minipage}\hfill
  \begin{minipage}[t]{0.235\linewidth}\poptile{poppink}{Exercices}{type examen + corrigés}\end{minipage}\hfill
  \begin{minipage}[t]{0.235\linewidth}\poptile{popgreen}{Fiche bilan}{l'essentiel à retenir}\end{minipage}\par}

% Sommaire
\newtcolorbox{sommaire}{enhanced,colback=white,colframe=popink,boxrule=2.2pt,arc=10pt,
  drop shadow=popink,left=18pt,right=18pt,top=13pt,bottom=13pt,before skip=6pt,after skip=20pt,
  before upper={\poppill{Sommaire}{poppurple}{white}\par\vspace{9pt}\popsemi\fontsize{10.6}{14.5}\selectfont\color{popink}}}

% Signature de fin de cours — poussée tout en bas de la dernière page
\newcommand{\findecours}{%
  \par\vfill\nopagebreak
  \begin{center}\popsquiggle{poporange}\end{center}\vspace{4pt}
  \signatureonbuch
}
\AtBeginDocument{\def\llbracket{\mathopen{[\mkern-3.5mu[}}\def\rrbracket{\mathclose{]\mkern-3.5mu]}}} % intervalles d'entiers (glyphe ⟦ absent de la police)
% Glyphes absents de Fira Math : redéfinis à partir de glyphes disponibles
\AtBeginDocument{%
  \def\setminus{\mathbin{\backslash}}\let\smallsetminus\setminus
  \def\vdots{\mathord{\vbox{\baselineskip3.2pt\lineskiplimit0pt\kern4pt\hbox{.}\hbox{.}\hbox{.}}}}%
  \def\ddots{\mathinner{\mkern1mu\raise6.5pt\hbox{.}\mkern2mu\raise3.6pt\hbox{.}\mkern2mu\raise0.7pt\hbox{.}\mkern1mu}}%
  \def\triangle{\mathord{\scalebox{0.9}{$\symup{\Delta}$}}}%
  \def\nabla{\mathord{\raisebox{\depth}{\scalebox{1}[-1]{$\symup{\Delta}$}}}}%
  \def\checkmark{\popcheck{popdark}}%
  \def\star{\mathbin{\ast}}\def\bigstar{\mathord{\ast}}%
  \def\diamond{\mathbin{\scalebox{0.6}{$\Diamond$}}}%
  \def\bigcup{\mathop{\mathchoice{\raisebox{-0.3em}{\scalebox{1.7}{$\cup$}}}{\scalebox{1.25}{$\cup$}}{\cup}{\cup}}\displaylimits}%
  \def\bigcap{\mathop{\mathchoice{\raisebox{-0.3em}{\scalebox{1.7}{$\cap$}}}{\scalebox{1.25}{$\cap$}}{\cap}{\cap}}\displaylimits}%
  \def\lesseqqgtr{\mathrel{\lessgtr}}\def\bigoplus{\mathop{\oplus}\displaylimits}%
}
\providecommand{\ssi}{si et seulement si} % abréviation fréquente
\AtBeginDocument{\providecommand{\wideparen}{\overparen}} % arc de cercle

%% FIGFIX-BEGIN v5 — correctifs globaux des figures (docs/onbuch-generation/tools/figfix)
\usepackage{adjustbox}\usepackage{etoolbox}\tcbuselibrary{raster}
% toute figure TikZ/pgfplots plus large que la ligne est réduite à la largeur disponible
\BeforeBeginEnvironment{tikzpicture}{\begin{adjustbox}{max width=\linewidth}}
\AfterEndEnvironment{tikzpicture}{\end{adjustbox}}
% légendes pgfplots lisibles par-dessus les courbes
\pgfplotsset{every axis/.append style={legend style={fill=white,fill opacity=0.92,text opacity=1,draw=black!15,inner sep=3pt}}}
% étiquettes d'arêtes (syntaxe edge["..."]) avec fond blanc
\tikzset{every edge quotes/.append style={fill=white,inner sep=1.2pt}}
%% FIGFIX-END
\begin{document}

\pagegarde{Les facteurs de la mondialisation}{Géographie}{Tle E}{Module 2 : La libéralisation des échanges (la mondialisation)}{N11}{2 h}{Cette leçon analyse les facteurs qui expliquent l'accélération de la mondialisation : révolution des transports et des télécommunications, puissance des firmes multinationales, rôle des États et des diasporas, action des associations et des ONG, et explosion du commerce international. À partir de cartes, de photographies et de données chiffrées, l'élève apprend à observer, cartographier et interpréter les flux qui structurent le monde contemporain, avec des illustrations camerounaises.
\vspace{4pt}
\begin{multicols}{2}\small
\begin{itemize}
\item À la fin de cette leçon, je sais définir la mondialisation et citer ses principaux facteurs
\item À la fin de cette leçon, je sais décrire la révolution des transports (conteneurisation, aviation, autoroutes de la mer) à partir de documents
\item À la fin de cette leçon, je sais expliquer le rôle des télécommunications et d'Internet dans l'accélération des échanges
\item À la fin de cette leçon, je sais définir une firme multinationale et décrire ses stratégies (DIT, IDE) à partir d'exemples présents au Cameroun
\item À la fin de cette leçon, je sais analyser le rôle des États (libéralisation, accords, organisations) et des diasporas dans la mondialisation
\item À la fin de cette leçon, je sais présenter le rôle des associations et des ONG comme acteurs de la mondialisation
\end{itemize}\end{multicols}}

\begin{sommaire}
\begin{multicols}{2}
\begin{enumerate}
\item Introduction : la mondialisation et ses moteurs
\item Les transports et les télécommunications
\item Les firmes multinationales (FMN)
\item L'État et la diaspora
\item Les associations et les ONG
\item Le commerce international et synthèse
\item Activité d'intégration
\item Méthodes et exercices résolus
\item Exercices
\item Corrigés des exercices
\item Fiche bilan
\end{enumerate}
\end{multicols}
\end{sommaire}

\begin{objectifs}
\begin{itemize}
\item À la fin de cette leçon, je sais définir la mondialisation et citer ses principaux facteurs
\item À la fin de cette leçon, je sais décrire la révolution des transports (conteneurisation, aviation, autoroutes de la mer) à partir de documents
\item À la fin de cette leçon, je sais expliquer le rôle des télécommunications et d'Internet dans l'accélération des échanges
\item À la fin de cette leçon, je sais définir une firme multinationale et décrire ses stratégies (DIT, IDE) à partir d'exemples présents au Cameroun
\item À la fin de cette leçon, je sais analyser le rôle des États (libéralisation, accords, organisations) et des diasporas dans la mondialisation
\item À la fin de cette leçon, je sais présenter le rôle des associations et des ONG comme acteurs de la mondialisation
\item À la fin de cette leçon, je sais interpréter des données chiffrées sur le commerce international et ses déséquilibres
\item À la fin de cette leçon, je sais légender une carte, construire un croquis et un diagramme représentant les flux mondiaux
\end{itemize}
\end{objectifs}

\begin{prerequis}
\begin{itemize}
\item La notion de mondialisation vue en introduction du Module 2 (acteurs, flux, espaces)
\item Les grandes aires de puissance mondiale (Triade) étudiées en Première
\item Les notions de flux, de réseau et d'acteur géographique
\item La lecture et la légendation d'une carte (méthodes acquises depuis la Seconde)
\item Les échanges Nord-Sud et la notion d'inégalités de développement (Première)
\end{itemize}
\end{prerequis}

\newpage

\coursec{Introduction : la mondialisation et ses moteurs}

En octobre 2023, un jeune entrepreneur de Douala, appelons-le \textbf{Franck}, commande sur la plateforme \textbf{Alibaba} un lot de smartphones fabriqués à \textbf{Shenzhen} (Chine). Il paie par virement bancaire international via \textbf{SWIFT} en quelques secondes, suit son conteneur en temps réel grâce à une application de \textbf{tracking GPS} jusqu'au \textbf{Port Autonome de Douala (PAD)}, dédouane sa marchandise via le guichet unique \textbf{GUCE} (Guichet Unique du Commerce Extérieur), puis revend ses produits via \textbf{WhatsApp Business} à des clients de Yaoundé, Bafoussam, Ngaoundéré et même à sa cousine installée dans la diaspora camerounaise à \textbf{Paris}, qui paie par \textbf{Mobile Money} (Orange Money / MTN MoMo) depuis la France. Comment un tel circuit, impensable il y a trente ans, est-il devenu possible ? Quels acteurs — États, entreprises, réseaux humains, infrastructures — ont rendu le monde si interconnecté ? C'est à ces questions que répond la présente leçon, en identifiant les \cle{facteurs} qui ont fait naître et qui accélèrent la \cle{mondialisation}.

\courssub{Définition et compréhension intuitive}

Au sens large, la mondialisation n'est pas un phénomène nouveau : les routes de la soie, les grandes découvertes des XVe–XVIe siècles ou la première révolution industrielle ont déjà connecté des continents. Mais depuis les \textbf{années 1980}, on observe une \textbf{mutation qualitative} : la vitesse, le volume et la variété des échanges ont explosé, créant une \textbf{interdépendance systémique} entre les territoires. Il ne s'agit plus seulement d'échanger des marchandises, mais de fragmenter la production (\cle{décomposition verticale}), de faire circuler des capitaux instantanément, de diffuser des normes culturelles et de mobiliser des populations à l'échelle planétaire.

\begin{definition}[La mondialisation]
La \textbf{mondialisation} est le processus d'intensification et d'accélération des échanges (marchandises, services, capitaux, informations, personnes) et d'interdépendance croissante entre les différentes parties du monde, aboutissant à la formation d'un \emph{espace-monde} relativement intégré, tout en creusant de nouvelles inégalités entre territoires et acteurs.
\end{definition}

Cette définition souligne trois dimensions indissociables :
\begin{itemize}
    \item \textbf{L'extension spatiale :} la planète entière est concernée, même les zones en apparence isolées (ex. : l'arrivée de la 4G dans les villages de l'Adamaoua).
    \item \textbf{L'intensification des flux :} volume des échanges mondiaux multiplié par \num{7} entre 1980 et 2022 (OMC), flux financiers quotidiens supérieurs à \SI{6000}{milliards USD} (BRI).
    \item \textbf{L'interdépendance :} une crise locale (blocage du canal de Suez en 2021, pandémie de Covid-19, guerre en Ukraine) produit des effets immédiats sur les chaînes de valeur mondiales et le prix du pain à Douala.
\end{itemize}

\begin{savaistu}[Le mot « mondialisation » : une histoire récente]
Si le phénomène est ancien, le \textbf{terme} s'impose tardivement. En français, « mondialisation » (calque de l'anglais \emph{globalization}) devient courant \textbf{à la fin des années 1980}, concomitamment à la chute du mur de Berlin (1989) et à la fin de la Guerre froide. Auparavant, on parlait d'« internationalisation » (relations entre États-nations) ou de « planétarisation ». L'adoption du mot « mondialisation » marque la prise de conscience d'un changement d'échelle : le monde devient un \emph{seul} espace d'action pour les firmes et la finance, au-delà des frontières étatiques.
\end{savaistu}

\courssub{Observation de documents : le port de Douala et les flux mondiaux}

Pour saisir concrètement les moteurs de cette dynamique, observons deux documents « phares » du programme.

\textbf{Document 1 : Photographie aérienne du Port Autonome de Douala (PAD), prise en 2023.}
\emph{Description guidée :} On distingue le \textbf{terminal à conteneurs} (quais 14 à 18) avec ses \textbf{grues portiques (ship-to-shore)} de dernière génération (portique à 22 rangées), les \textbf{zones de stockage} où s'empilent des milliers de conteneurs \textbf{normalisés} (norme ISO 668 : 20 pieds, 40 pieds, 40 pieds High Cube) aux couleurs des armateurs mondiaux (Maersk \emph{bleu}, MSC \emph{rouge}, CMA CGM \emph{beige}, Cosco \emph{bleu foncé}). Au premier plan, le \textbf{chenal d'accès} dragué à \SI{-13,5}{m} (projet d'approfondissement à \SI{-15}{m}) permet le passage de porte-conteneurs \textbf{Post-Panamax} (capacité \num{8000} à \num{12000} EVP). En arrière-plan, le \textbf{pont sur le Wouri} (route et rail) et la ville de Douala, poumon économique du Cameroun.

\textbf{Document 2 : Carte des principaux flux mondiaux de marchandises (année 2022, source OMC / UNCTAD).}
\emph{Description guidée :} Planisphère en projection de Robinson. Les \textbf{flèches d'épaisseur proportionnelle} révèlent la structure tripolaire des échanges :
\begin{itemize}
    \item \textbf{Flux Transpacifique (Asie – Amérique du Nord) :} le plus dense (\textasciitilde \SI{30}{\%} du commerce maritime mondial), reliant la \emph{Usine du monde} (Chine, ASEAN) au premier marché de consommation (USA).
    \item \textbf{Flux Transatlantique (Europe – Amérique du Nord) :} flux intenses de produits finis, de composants high-tech et d'énergie (GNL).
    \item \textbf{Flux Europe – Asie (via Suez) :} artère vitale pour l'Europe (énergie, biens de consommation).
    \item \textbf{Flux Sud–Sud en forte croissance :} Afrique – Asie (minerais, pétrole, bois brut contre biens manufacturés), Afrique – Amérique latine, Asie – Amérique latine.
    \item \textbf{Position du Cameroun :} nœud secondaire sur le golfe de Guinée, entrée/sortie du corridor Douala–Ndjaména–Bangui (pays enclavés).
\end{itemize}

\begin{methode}[Démarche d'analyse d'un document géographique (photo, carte, croquis)]
Face à tout document en épreuve (Bac) ou en TP, suis \textbf{systématiquement} ces quatre étapes :
\begin{enumerate}
    \item \textbf{Observer :} Qu'est-ce que je vois ? (couleurs, formes, symboles, légende, échelle, date, source, auteur). Ne pas interpréter, seulement inventorier le visible.
    \item \textbf{Localiser :} Où cela se passe-t-il ? (coordonnées, pays, région, ville, repères physiques et humains). Situer par rapport à des axes, des frontières, des pôles.
    \item \textbf{Décrire :} Quels sont les éléments principaux et leur organisation ? (structures, contrastes, hiérarchie, flux, réseaux). Utiliser le vocabulaire précis (quai, grue, EVP, chenal, flux dominant, tripolaire, hub, feeder).
    \item \textbf{Interpréter / Expliquer :} Pourquoi cela est-il ainsi ? Quels processus, acteurs, politiques, contraintes physiques ou historiques produisent ce paysage ? Lier à la problématique du cours (facteurs de la mondialisation).
\end{enumerate}
\emph{Astuce Bac :} En introduction de commentaire, cite toujours la \textbf{nature}, le \textbf{titre}, la \textbf{source}, la \textbf{date} et l'\textbf{échelle} du document. En conclusion, ouvre sur une autre échelle ou un avenir possible.
\end{methode}

\begin{exoresolu}[Commentaire guidé : Le Port de Douala, porte de la mondialisation camerounaise]
\textbf{Énoncé :} À partir du Document 1 (photo du PAD) et de vos connaissances, montrez que le port de Douala est un acteur et un révélateur de la mondialisation au Cameroun.

\textbf{Solution détaillée :}
\tcblower
\textbf{1. Introduction}
\begin{itemize}
    \item \textbf{Nature :} Photographie aérienne verticale (vue zénithale).
    \item \textbf{Titre :} « Terminal à conteneurs du Port Autonome de Douala, 2023 ».
    \item \textbf{Source :} Port Autonome de Douala / Satellite (ex. Airbus Defence / Google Earth).
    \item \textbf{Localisation :} Douala, capitale économique du Cameroun, sur le fleuve Wouri, golfe de Guinée, façade atlantique de l'Afrique centrale.
    \item \textbf{Problématique :} En quoi cette infrastructure portuaire illustre-t-elle l'insertion du Cameroun dans la mondialisation ?
    \item \textbf{Plan :} I. Une infrastructure technique adaptée aux standards mondiaux. II. Un nœud de flux connectant l'hinterland africain au monde. III. Un révélateur des limites et défis de l'intégration.
\end{itemize}

\textbf{2. Développement}
\textbf{I. Une infrastructure technique adaptée aux standards mondiaux (Facteur technique : transports)}
\begin{itemize}
    \item \textbf{Normalisation :} Conteneurs ISO (20/40 pieds) = langage universel du transport multimodal (navire – camion – train). Permet la \textbf{rupture de charge} rapide.
    \item \textbf{Moyens de manutention :} Grues \emph{ship-to-shore} (STS) et portiques sur pneus (RTG) = productivité visée \num{30}–\num{35} mouvements/heure/grue. Investissements récents (projet d'extension \num{150} milliards FCFA financé par \textbf{AFD}, \textbf{BEI}, \textbf{BAD}).
    \item \textbf{Tirant d'eau :} \SI{-13,5}{m} (objectif \SI{-15}{m}) = accès aux navires \textbf{Post-Panamax} (\num{8000}–\num{12000} EVP), évitant le transbordement coûteux (ex. via Tanger Med ou Algeciras).
\end{itemize}

\textbf{II. Un nœud de flux connectant l'hinterland au monde (Facteurs économiques : commerce, FMN)}
\begin{itemize}
    \item \textbf{Flux sortants (Export) :} Bois (grumes, sciages), cacao, coton, café, aluminium (Alucam), pétrole brut (Kribi mais transit Douala), banane. Destination : UE (Pays-Bas, France, Belgique), Chine, USA.
    \item \textbf{Flux entrants (Import) :} Produits pétroliers raffinés, riz, blé, médicaments, véhicules, machines, biens de consommation (téléphones, textiles). Origine : Chine, France, Belgique, Inde, Nigeria.
    \item \textbf{Acteurs privés (FMN) :} \textbf{Bolloré Africa Logistics} (concessionnaire terminal conteneurs), \textbf{Maersk}, \textbf{MSC}, \textbf{CMA CGM}, \textbf{SDV} (logistique). Présence de \textbf{Zones Franches} (usines de transformation pour réexport).
    \item \textbf{Hinterland :} Corridor routier/ferroviaire Douala–Yaoundé–Ngaoundéré–Ndjaména (Tchad)–Bangui (RCA). Le port dessert \textbf{deux pays enclavés} (Tchad, RCA) $\rightarrow$ enjeu géopolitique régional.
\end{itemize}

\textbf{III. Un révélateur des limites et défis (Facteurs politiques, institutionnels)}
\begin{itemize}
    \item \textbf{Concurrence portuaire :} Montée en puissance de \textbf{Kribi} (port en eaux profondes, \SI{-16}{m}, terminal minéralier Mbalam, terminal conteneurs futur) et des ports voisins (Lagos/Lekki au Nigeria, Libreville/Owendo au Gabon, Pointe-Noire au Congo).
    \item \textbf{Gouvernance :} Délais de dédouanement encore longs (moyenne \num{5}–\num{7} jours malgré GUCE), coûts logistiques élevés (Douala dans le top 20 des ports les plus chers d'Afrique \emph{Doing Business}), corruption résiduelle.
    \item \textbf{Urbanisation contrainte :} Ville-port saturée, conflits d'usage (pêche, habitat, industrie), pollution, vulnérabilité climatique (montée des eaux, érosion côtière).
\end{itemize}

\textbf{3. Conclusion}
Le port de Douala est une \textbf{interface majeure} de la mondialisation camerounaise : il matérialise la \textbf{technicisation des transports} (conteneurisation, navires géants), l'\textbf{insertion dans le commerce international} (spécialisation primaire, dépendance aux importations manufacturées) et le \textbf{rôle des FMN} (concession Bolloré, armateurs). Mais sa position de \textbf{hub régional} est menacée par la concurrence (Kribi, Lagos) et les goulets d'étranglement institutionnels. L'avenir passe par le \textbf{corridor ferroviaire} (Camrail rénové, projet Douala–Ngaoundéré standard gauge) et l'\textbf{amélioration du climat des affaires} (facteur État).
\end{exoresolu}

\courssub{Problématique et annonce du plan}

L'observation croisée de la photographie du port de Douala (nœud local) et de la carte des flux mondiaux (structure globale) nous conduit à formuler la \textbf{problématique centrale} de ce module :

\begin{center}
\textbf{\large « Quels sont les facteurs techniques, économiques, politiques et sociaux qui expliquent l'accélération de la mondialisation depuis les années 1980, et comment s'articulent-ils pour façonner les territoires, notamment au Cameroun ? »}
\end{center}

Pour y répondre, nous identifierons \textbf{cinq familles de facteurs} qui agissent en \textbf{synergie} (et non de façon isolée). Le schéma ci-dessous résume l'architecture de la leçon :

\begin{popfigure}
\begin{tikzpicture}[
    scale=1.1,
    every node/.style={transform shape},
    center node/.style={circle, draw=popblue, thick, fill=popblueL, minimum width=3.2cm, text width=3cm, align=center, font=\bfseries\large},
    factor node/.style={rectangle, rounded corners=8pt, draw=popdark, thick, fill=popcream, minimum width=2.8cm, minimum height=1.6cm, text width=2.6cm, align=center, font=\small, drop shadow={shadow xshift=2pt, shadow yshift=-2pt, opacity=0.3}},
    arrow/.style={-Stealth, thick, popline, shorten >=3pt, shorten <=3pt}
]
% Centre
\node[center node, align=center] (center) at (0,0){Mondialisation\\(accélération\\depuis 1980)};

% Facteurs - positions en étoile
\node[factor node, fill=popblueL, align=center] (trans) at (0,4.2){1. Transports\\\emph{(conteneurisation,\\navires géants,\\corridors)}};
\node[factor node, fill=poporangeL, align=center] (telecom) at (3.6,1.2){2. Télécommunications\\\emph{(Internet, fibre,\\satellite, mobile,\\paiement numérique)}};
\node[factor node, fill=popgreenL, align=center] (fmn) at (2.1,-3.3){3. Firmes\\multinationales\\\emph{(IDE, chaînes\\de valeur,\\R\&D, marques)}};
\node[factor node, fill=poppurpleL, align=center] (etat) at (-2.1,-3.3){4. État \& Diaspora\\\emph{(libéralisation,\\accords, ZLECAf,\\transferts, compétences)}};
\node[factor node, fill=poppinkL, align=center] (ong) at (-3.6,1.2){5. ONG \& Commerce\\\emph{(normes, plaidoyer,\\humanitaire, OMC,\\règles du jeu)}};

% Flèches
\draw[arrow] (center) -- (trans);
\draw[arrow] (center) -- (telecom);
\draw[arrow] (center) -- (fmn);
\draw[arrow] (center) -- (etat);
\draw[arrow] (center) -- (ong);

% Légende interne discrète
\node[font=\tiny\itshape, text=popmuted] at (0,-4.8) {Schéma de synthèse : les 5 moteurs de la mondialisation (Module 2, Leçon 11)};
\end{tikzpicture}
\legende{Schéma en étoile : la mondialisation au centre, structurée par cinq facteurs interdépendants. Chaque facteur fera l'objet d'une section détaillée (Sections 2 à 6).}
\end{popfigure}

\textbf{Annonce du plan détaillé :}
\begin{enumerate}
    \item \textbf{Section 2 : Les transports et les télécommunications} (facteurs techniques) — La révolution du conteneur, le gigantisme naval, les corridors multimodaux (Douala–Ndjaména), la fibre optique (WACS, SAT-3, 2Africa), le mobile money.
    \item \textbf{Section 3 : Les firmes multinationales (FMN)} (facteur économique principal) — Stratégies d'IDE, fragmentation de la production (chaines de valeur mondiales), cas des firmes au Cameroun (Bolloré, Maersk, Heineken, SOCAPALM, Alucam, groupes télécoms).
    \item \textbf{Section 4 : L'État et la diaspora} (facteurs politiques et sociaux) — Libéralisation (PAS, PPTE), accords commerciaux (APE UE–Afrique, ZLECAf), rôle des transferts de la diaspora (estimation \textasciitilde \SI{300}{milliards FCFA/an}), transfert de compétences, lobbying.
    \item \textbf{Section 5 : Les associations et les ONG} (facteurs sociétaux et normatifs) — ONG humanitaires (MSF, Croix-Rouge), ONG de développement (CARE, Plan International), ONG environnementales (WWF, Greenpeace), syndicats internationaux, rôle dans la définition de normes (RSE, ODD).
    \item \textbf{Section 6 : Le commerce international et synthèse} — Évolution des flux (OMC, UNCTAD), termes de l'échange, balance commerciale du Cameroun, synthèse intégrée des cinq facteurs.
\end{enumerate}

\begin{attention}[Pièges à éviter dès l'introduction]
\begin{itemize}
    \item \textbf{Confondre « mondialisation » et « internationalisation » :} L'internationalisation = relations entre États-nations (commerce inter-industriel). La mondialisation = stratégies globales des firmes (commerce intra-industriel, IDE, chaînes de valeur), affaiblissement relatif de la frontière étatique.
    \item \textbf{Oublier le facteur « État » :} On croit souvent que la mondialisation = « retrait de l'État ». \textbf{Faux :} l'État \textbf{construit} la mondialisation (libéralisation, infrastructures, traités, ZLECAf, guichets uniques). Au Cameroun, l'État est actionnaire majeur (SNPC, Camrail, Camtel, PAD) et régulateur (ARMP, ART, MINCOMMERCE).
    \item \textbf{Négliger l'échelle locale :} Le Bac attend que tu ancre tes connaissances dans le \textbf{territoire camerounais} (Douala, Kribi, corridors, diaspora, FMN locales). Un cours sans exemple camerounais précis est un cours incomplet.
    \item \textbf{Lister les facteurs sans montrer leurs interactions :} Le conteneur (transport) ne fonctionne que grâce aux TIC (tracking, douane dématérialisée), aux normes (ISO, OMC), aux investissements (FMN Bolloré) et à la régulation (État). C'est le \textbf{système technique} qui compte.
\end{itemize}
\end{attention}

\coursec{Les transports et les télécommunications}

Nous venons de voir, dans l'introduction, que la mondialisation est un processus d'intensification des échanges de toutes natures (marchandises, capitaux, personnes, informations) à l'échelle de la planète. Mais un tel processus ne se décrète pas : il repose sur des \cle{facteurs}, c'est-à-dire des moteurs concrets qui rendent ces échanges possibles, rapides et bon marché. Le premier de ces moteurs, et sans doute le plus fondamental, est la \cle{révolution des transports et des télécommunications}. Sans navires capables de franchir les océans à bas coût, sans avions, sans câbles et sans satellites, les autres acteurs que nous étudierons ensuite (firmes multinationales, États, diasporas, ONG) ne pourraient tout simplement pas agir à l'échelle mondiale. Cette section répond à une question simple : comment les progrès des transports et des communications ont-ils \og rétréci \fg{} la planète, et avec quelles limites, notamment pour un pays comme le Cameroun ?

\courssub{La révolution des transports maritimes : la conteneurisation}

\textbf{Une observation pour commencer.} Au port autonome de Douala, sur le Wouri, on observe chaque semaine d'immenses navires déchargeant des boîtes métalliques standardisées, que des portiques soulèvent et déposent directement sur des camions. Ces boîtes, ce sont des \cle{conteneurs}. Leur apparition dans les années 1960 a bouleversé le commerce mondial bien plus que n'importe quel traité de libre-échange.

\begin{definition}[La conteneurisation]
La \cle{conteneurisation} est le transport de marchandises dans des \cle{conteneurs normalisés} (boîtes métalliques aux dimensions universelles, mesurées en \cle{EVP}, équivalent vingt pieds, soit environ 6,1 mètres de long). En supprimant la manutention marchandise par marchandise, elle réduit fortement les \textbf{coûts} et les \textbf{délais} de chargement et de déchargement.
\end{definition}

Avant le conteneur, le déchargement d'un navire exigeait des centaines de dockers pendant plusieurs jours : chaque sac, chaque caisse, chaque fût était manipulé individuellement. Avec le conteneur, une grue déplace une boîte entière en quelques minutes, et cette boîte passe du navire au camion ou au train \textbf{sans être ouverte}. C'est ce qu'on appelle la \emph{rupture de charge} sans rupture de contenu : la marchandise voyage de l'usine au consommateur dans la même boîte scellée.

\begin{popfigure}
\begin{center}
\begin{tikzpicture}[font=\small, >=Stealth]
% Usine
\draw[fill=popblueL] (0,0) rectangle (1.6,1.2);
\draw[fill=popblueL] (0,1.2) -- (0,1.6) -- (0.5,1.6) -- (0.5,1.2);
\node at (0.8,0.5) {\textbf{Usine}};
\node[below] at (0.8,0) {production};
% Camion 1
\draw[->, very thick, poporange] (1.7,0.6) -- (3.1,0.6);
\node[above, popdark] at (2.4,0.6) {camion};
% Port de départ
\draw[fill=popgreenL] (3.2,0) rectangle (4.8,1.2);
\node at (4.0,0.6) {\textbf{Port A}};
\node[below] at (4.0,0) {hub de départ};
% Navire
\draw[->, very thick, popblue] (4.9,0.6) -- (6.9,0.6);
\node[above, popdark] at (5.9,0.6) {porte-conteneurs};
% Port d'arrivée
\draw[fill=popgreenL] (7.0,0) rectangle (8.6,1.2);
\node at (7.8,0.6) {\textbf{Port B}};
\node[below] at (7.8,0) {hub d'arrivée};
% Camion 2
\draw[->, very thick, poporange] (8.7,0.6) -- (10.1,0.6);
\node[above, popdark] at (9.4,0.6) {camion};
% Consommateur
\draw[fill=poppinkL] (10.2,0) rectangle (11.8,1.2);
\node at (11.0,0.6) {\textbf{Marché}};
\node[below] at (11.0,0) {consommateur};
% Conteneur (le même tout du long)
\draw[fill=popgoldL, thick] (5.2,1.4) rectangle (6.6,2.0);
\node at (5.9,1.7) {\footnotesize conteneur scellé};
\draw[dashed, popmuted] (5.9,1.4) -- (5.9,0.8);
\end{tikzpicture}
\end{center}
\legende{Schéma 1 : la chaîne de la conteneurisation. Le même conteneur scellé circule de l'usine au consommateur, en changeant de mode de transport (route, mer, route) sans que la marchandise soit manipulée : c'est le principe du \og porte-à-porte \fg{}.}
\end{popfigure}

\textbf{Le gigantisme naval.} Pour rentabiliser le conteneur, les armateurs ont construit des navires de plus en plus grands : c'est la course au \cle{gigantisme}.

\begin{exemplebox}[Des navires toujours plus grands]
Dans les années 1960, un porte-conteneurs transportait quelques centaines d'EVP. Aujourd'hui, les plus grands porte-conteneurs du monde (par exemple ceux des compagnies MSC, Maersk ou CMA CGM) peuvent emporter \textbf{plus de 24 000 EVP}, soit l'équivalent d'un train de plus de 100 km de long. Un tel navire transporte des marchandises pour un coût unitaire dérisoire : acheminer un téléviseur de Chine en Europe revient à quelques dollars seulement.
\end{exemplebox}

\textbf{Les autoroutes de la mer et les hubs.} Ces navires géants ne desservent pas tous les ports : ils suivent des routes régulières entre quelques très grands ports, les \cle{autoroutes de la mer}, et s'arrêtent dans des plateformes de concentration et de redistribution appelées \cle{hubs}.

\begin{definition}[Le hub]
Un \cle{hub} (plateforme de correspondance) est un lieu de \textbf{concentration et de redistribution des flux} : portuaire (Singapour, Rotterdam, Shanghai), aéroportuaire (Dubaï, Addis-Abeba) ou numérique (centres de données). Les grands flux y arrivent, y sont triés, puis redispatchés vers des destinations plus petites.
\end{definition}

\begin{popfigure}
\includegraphics[width=\linewidth,height=9.5cm,keepaspectratio]{N11/shipping_routes.jpg}
\legende{Le trafic maritime commercial mondial (plus le bleu est foncé, plus le trafic est dense ; carte sans noms de lieux). Repère les grands corridors : Atlantique Nord entre l'Amérique du Nord et l'Europe, Pacifique Nord entre l'Asie et l'Amérique, axe Méditerranée--Suez--océan Indien--détroit de Malacca vers l'Asie de l'Est, et route du cap de Bonne-Espérance. Le golfe de Guinée, où se trouve Douala, est traversé par un trafic nettement plus faible : le Cameroun est en marge des grandes autoroutes de la mer, ce qui renchérit ses coûts de transport. \\ Source : Wikimedia Commons, « Shipping routes » (densité du trafic maritime commercial), T. Hengl, CC BY-SA 3.0.}
\end{popfigure}

\textbf{Et le Cameroun ?} Le port autonome de \cle{Douala} concentre l'essentiel du commerce extérieur camerounais (environ 90 \% des volumes, ordre de grandeur) et sert aussi le Tchad et la Centrafrique. Mais c'est un port sur fleuve, à l'accès nautique difficile (ensablement du chenal), qui ne peut recevoir les très grands navires. C'est pourquoi l'État a fait construire le \cle{port en eau profonde de Kribi} (au sud, sur l'océan), inauguré en 2014 (mise en service progressive) : il accueille des navires de plus fort tonnage et ambitionne de devenir un hub régional pour l'Afrique centrale. Ce projet illustre bien la logique des hubs : capter les grands flux mondiaux pour les redistribuer vers l'hinterland continental.

\courssub{Le transport aérien : la vitesse pour la valeur et les personnes}

Le transport maritime est roi pour les marchandises lourdes et peu périssables, mais il est lent : trois à cinq semaines entre l'Asie et l'Europe. Le \cle{transport aérien} répond à un autre besoin : la \textbf{rapidité}.

\begin{itemize}
\item Il transporte les \textbf{produits à forte valeur ajoutée} et de faible masse : composants électroniques, médicaments, pièces détachées, produits de luxe, mais aussi des denrées périssables (fleurs coupées, poissons, fruits tropicaux exportés par avion).
\item Il transporte surtout les \textbf{personnes} : hommes d'affaires, touristes, étudiants, diasporas. Le nombre de passagers aériens dans le monde a été multiplié par près de 50 depuis 1960 (ordre de grandeur), avant même la crise sanitaire de 2020.
\end{itemize}

Au Cameroun, les aéroports internationaux de \cle{Douala} et de \cle{Yaoundé-Nsimalen} assurent les liaisons avec l'Europe (Paris, Bruxelles, Istanbul), l'Afrique (Addis-Abeba, Nairobi, Casablanca) et le reste du monde. Ils jouent un rôle essentiel pour la mobilité des élites, de la diaspora et pour l'exportation de produits périssables. Mais le fret aérien reste marginal en volume : moins de 1 \% du commerce mondial en tonnage, contre environ 35 \% en valeur (ordres de grandeur). C'est la preuve que l'avion sert la \emph{valeur}, pas le \emph{poids}.

\courssub{Les transports terrestres : corridors et pays enclavés}

Sur les continents, ce sont les \textbf{routes} et les \textbf{chemins de fer} qui prolongent les ports vers l'intérieur des terres : on parle de \cle{corridors}. Un corridor est un axe de transport aménagé reliant un port à son arrière-pays (\emph{hinterland}).

Pour l'Afrique centrale, l'exemple est frappant. Le \cle{Tchad} et la \cle{Centrafrique} sont des \textbf{pays enclavés} (sans accès à la mer) : tout leur commerce extérieur transite par le corridor \cle{Douala--N'Djamena} (environ 1 800 km) ou Douala--Bangui. Or cet axe est un véritable \textbf{goulot d'étranglement} : routes en mauvais état par tronçons, ponts limités en tonnage, nombreux contrôles et tracasseries qui allongent les délais. Le chemin de fer camerounais (ligne \cle{Transcam} Douala--Yaoundé--Ngaoundéré, environ 1 000 km) dessert le nord du Cameroun mais ne rejoint ni le Tchad ni la Centrafrique. Résultat : un conteneur peut mettre plus de temps à aller de Douala à N'Djamena que de Shanghai à Douala. La mondialisation des transports reste donc très inégale selon les territoires.

\courssub{Une baisse spectaculaire des coûts et des durées depuis 1950}

La révolution des transports se mesure par un indicateur simple : l'évolution des coûts. Le graphique ci-dessous présente des \textbf{indices} (base 100 en 1950) du coût moyen du transport maritime et aérien. Il s'agit d'ordres de grandeur pédagogiques, cohérents avec les grandes tendances établies par les économistes (forte baisse du fret maritime grâce à la conteneurisation, baisse continue du fret aérien grâce aux avions à réaction puis gros-porteurs).

\begin{popfigure}
\begin{center}
\begin{tikzpicture}
\begin{axis}[
  ybar, bar width=7pt,
  width=0.95\linewidth, height=6.2cm,
  symbolic x coords={1950,1970,1990,2010,2020},
  xtick=data,
  ymin=0, ymax=110,
  ylabel={Indice du coût (base 100 en 1950)},
  legend style={at={(0.5,-0.18)}, anchor=north, legend columns=2, font=\small},
  nodes near coords, every node near coord/.append style={font=\scriptsize},
  ymajorgrids, grid style={popline!40},
]
\addplot[fill=popblue!70, draw=popblue] coordinates {(1950,100) (1970,60) (1990,35) (2010,25) (2020,20)};
\addplot[fill=poporange!70, draw=poporange] coordinates {(1950,100) (1970,45) (1990,25) (2010,15) (2020,12)};
\legend{Fret maritime, Fret aérien}
\end{axis}
\end{tikzpicture}
\end{center}
\legende{Graphique 3 : évolution indicative des coûts moyens du transport maritime et aérien, en indice base 100 en 1950 (ordres de grandeur, sources : estimations synthétiques d'après les travaux sur l'histoire des coûts de transport). Lecture : en 2020, transporter une tonne de marchandise par mer coûte environ cinq fois moins cher qu'en 1950, à monnaie constante.}
\end{popfigure}

\begin{exoresolu}[Lire et interpréter un graphique de coûts de transport]
\textbf{Énoncé.} À partir du graphique 3 : 1) Calculez le taux de baisse du coût du fret maritime entre 1950 et 2020. 2) Dégagez une conséquence de cette évolution pour le commerce mondial. 3) Expliquez pourquoi cette baisse ne profite pas également à tous les territoires.\par
\tcblower
\textbf{Solution détaillée.}\par
1) Le taux de variation se calcule ainsi :
\[
t = \frac{V_{2020}-V_{1950}}{V_{1950}}\times 100 = \frac{20-100}{100}\times 100 = -80\,\%.
\]
Le coût du fret maritime a donc \textbf{baissé d'environ 80 \%} entre 1950 et 2020 (indice passé de 100 à 20).\par
2) Conséquence : le transport ne représente plus un obstacle majeur au commerce. Il devient rentable de produire chaque étape d'un bien dans un pays différent (conception en Europe, assemblage en Chine, commercialisation au Cameroun) : c'est la base matérielle de la \textbf{mondialisation de la production} et de l'explosion des échanges (le commerce mondial a été multiplié par plus de 30 en valeur depuis 1950, ordre de grandeur).\par
3) Limite : la baisse des coûts profite d'abord aux espaces \textbf{connectés aux grands réseaux} (ports hub, corridors modernes). Les arrière-pays mal desservis, comme le corridor Douala--N'Djamena, cumulent mauvaises infrastructures et délais : leurs coûts de transport réels restent élevés, ce qui marginalise les pays enclavés malgré la baisse mondiale.
\end{exoresolu}

\courssub{La révolution des télécommunications : le monde en temps réel}

Transporter des \textbf{informations} est encore plus décisif que transporter des marchandises. Trois innovations ont transformé les communications depuis les années 1960.

\begin{itemize}
\item \textbf{Les câbles sous-marins de fibre optique} : ils acheminent aujourd'hui l'immense majorité (plus de 95 \%) du trafic Internet intercontinental, sous forme d'impulsions lumineuses, à des débits considérables.
\item \textbf{Les satellites} : ils assurent la télévision, la téléphonie et l'Internet des zones isolées, ainsi que le positionnement (GPS).
\item \textbf{Internet et la téléphonie mobile} : le réseau mondial connecte aujourd'hui plus de 5 milliards de personnes (ordre de grandeur) ; en Afrique, le téléphone mobile a permis de \og sauter \fg{} l'étape du téléphone fixe.
\end{itemize}

\begin{savaistu}[Le Cameroun branché sur la fibre mondiale]
Le Cameroun est connecté à l'Internet haut débit international depuis \textbf{2002} grâce au câble sous-marin \cle{SAT-3/WASC}, qui longe la côte occidentale africaine du Portugal à l'Afrique du Sud, avec un point d'atterrissement à Douala. D'autres câbles sont venus ensuite renforcer et sécuriser cette connexion : \textbf{WACS} (2012), \textbf{ACE} (2012) et \textbf{SAIL} (2018, reliant le Cameroun au Brésil). Un câble national de fibre optique relie en outre les grandes villes, mais son extension vers les zones rurales reste inachevée.
\end{savaistu}

\textbf{Les conséquences sont immenses.}

\begin{itemize}
\item \textbf{La compression de l'espace-temps} : une information fait le tour du monde en une fraction de seconde. La distance n'a quasiment plus d'effet sur la communication, alors qu'elle pèse encore sur le transport physique.
\item \textbf{La financiarisation en temps réel} : les bourses de New York, Londres, Tokyo ou Johannesburg sont interconnectées ; des milliards de francs circulent par voie électronique en continu, 24 heures sur 24.
\item \textbf{Le commerce électronique} : on achète depuis Yaoundé un produit expédié de Chine (plateformes comme Alibaba, Jumia au niveau africain).
\item \textbf{Les réseaux sociaux} : WhatsApp, Facebook ou TikTok relient instantanément la diaspora camerounaise au pays, transforment l'information, la politique et le commerce informel.
\end{itemize}

\begin{aretenir}[L'essentiel de la section]
La mondialisation repose d'abord sur une \textbf{révolution matérielle} : la conteneurisation et le gigantisme naval ont effondré les coûts du transport maritime (environ $-80\,\%$ depuis 1950), l'avion a conquis le transport des personnes et des produits de valeur, et les télécommunications (câbles, satellites, Internet, mobile) ont aboli le temps de transmission de l'information. Il en résulte une \textbf{compression de l'espace-temps} : le monde est devenu un \og village global \fg{}... mais seulement pour les territoires bien connectés.
\end{aretenir}

\courssub{Une limite majeure : la fracture numérique et les espaces déconnectés}

Il serait faux de conclure que la révolution des transports et des communications a uniformément \og aplati \fg{} le monde. Deux inégalités persistent.

\begin{itemize}
\item \textbf{La fracture numérique} : l'accès à Internet reste très inégal. En Afrique subsaharienne, environ 40 \% de la population utilise Internet (ordre de grandeur, en progression rapide), contre plus de 85 \% en Europe. Au sein du Cameroun, le fossé est net entre les villes connectées (Douala, Yaoundé) et les campagnes où le réseau mobile et l'électricité manquent.
\item \textbf{Les espaces mal desservis} : les arrière-pays enclavés ou dépourvus d'infrastructures (nord du Cameroun, Tchad, Centrafrique) supportent des coûts de transport élevés et restent à l'écart des flux mondiaux.
\end{itemize}

\begin{attention}[Ne pas confondre vitesse et accessibilité]
Piège classique du Bac : affirmer que \og la mondialisation a supprimé les distances \fg{}. C'est faux. La \textbf{vitesse} des transports et des communications a prodigieusement augmenté, mais l'\textbf{accessibilité} reste inégale : un espace ne profite de la mondialisation que s'il est \emph{connecté} aux réseaux (ports, corridors, câbles). L'arrière-pays camerounais, mal desservi par les routes et le numérique, reste marginalisé malgré la mondialisation. Dans une dissertation, opposez toujours les \textbf{espaces intégrés} (hubs, métropoles, littoraux) aux \textbf{espaces périphériques ou déconnectés}.
\end{attention}

\begin{exercice}[Pour s'entraîner : croquis du corridor Douala--N'Djamena]
Construisez un croquis schématique mettant en évidence le rôle du corridor Douala--N'Djamena pour les pays enclavés d'Afrique centrale. On attend : un titre, une flèche du nord, le port de Douala, l'axe routier (trait épais), la ligne ferroviaire Transcam (trait différent), les villes de Yaoundé, Ngaoundéré et N'Djamena, et une légende organisée. Analysez ensuite en cinq lignes pourquoi ce corridor est un goulot d'étranglement.
\end{exercice}

\begin{corrige}[Exercice : corridor Douala--N'Djamena]
Le croquis doit montrer le port de Douala sur la côte, l'axe routier remontant vers le nord-est via Yaoundé puis Ngaoundéré jusqu'à N'Djamena (Tchad), et la voie ferrée Transcam s'arrêtant à Ngaoundéré. Analyse attendue : le Tchad (et la Centrafrique) n'ayant pas de façade maritime, tout leur commerce extérieur transite par ce corridor unique ; or la route est longue (environ 1 800 km), dégradée par tronçons, le rail ne rejoint pas le Tchad, et les contrôles multiplient les délais. Les coûts et durées de transport y restent donc très élevés, ce qui pénalise les économies enclavées et illustre l'inégale accessibilité des territoires à la mondialisation.
\end{corrige}

Ainsi, transports et télécommunications forment le \textbf{premier moteur} de la mondialisation : ils ont réduit les coûts, les délais et les distances. Mais ces réseaux ont besoin d'acteurs pour les exploiter et les organiser à l'échelle planétaire. Ce sont précisément les \textbf{firmes multinationales}, que nous étudierons dans la section suivante.

\coursec{Les firmes multinationales (FMN)}

Nous venons de voir que les transports et les télécommunications ont « rétréci » la planète : un conteneur peut quitter Douala et arriver à Rotterdam en quelques semaines, un message peut circuler de Yaoundé à Shanghai en une fraction de seconde. Mais qui utilise ces réseaux ? Qui organise ces flux ? Derrière les infrastructures, il y a des acteurs. Et parmi eux, les plus puissants sont les \cle{firmes multinationales} (FMN). Ce sont elles qui décident où produire, où investir, où vendre. Comprendre les FMN, c'est comprendre le moteur économique principal de la mondialisation.

\courssub{Définition et caractéristiques d'une firme multinationale}

\textbf{Une situation pour commencer.} Tu es au marché Mokolo à Yaoundé. Tu achètes un téléphone de marque Tecno, conçu par une entreprise chinoise, assemblé en Chine avec des composants coréens, distribué au Cameroun par des grossistes. Au coin de la rue, une station-service TotalEnergies (entreprise française), une boutique Orange (entreprise française), un kiosque qui vend des produits Nestlé (entreprise suisse). En quelques mètres, tu as croisé quatre firmes multinationales. Qu'est-ce qui les définit ?

\begin{definition}[Firme multinationale (FMN)]
Une \cle{firme multinationale} (FMN), appelée aussi \emph{société transnationale}, est une entreprise implantée dans plusieurs pays grâce à ses \cle{filiales}, et qui organise sa production, sa recherche et sa vente à l'échelle mondiale. Elle est dirigée par une \cle{maison-mère} (ou société-mère), siège de la décision stratégique, généralement située dans le pays d'origine de l'entreprise.
\end{definition}

\begin{aretenir}[Les quatre caractéristiques d'une FMN]
\begin{itemize}
\item Une \textbf{maison-mère} qui concentre la stratégie, la recherche-développement (R\&D) et la gestion financière ;
\item Des \textbf{filiales} implantées dans plusieurs pays (parfois plus d'une centaine), qui produisent ou vendent localement ;
\item Un \textbf{chiffre d'affaires} considérable, parfois supérieur au PIB de certains États ;
\item Une \textbf{stratégie mondiale} : la FMN pense sa production à l'échelle de la planète, et non pays par pays.
\end{itemize}
\end{aretenir}

\textbf{Observation de documents.} Le programme te demande d'observer des photographies et des logos. Prenons trois exemples visibles au Cameroun : le logo de \textbf{TotalEnergies} sur les stations-service (maison-mère à Paris, présence dans environ 130 pays) ; le logo jaune de \textbf{MTN} (maison-mère en Afrique du Sud, présente dans une vingtaine de pays d'Afrique et du Moyen-Orient) ; le logo d'\textbf{Orange} (maison-mère française). Dans chaque cas, la marque est identique d'un pays à l'autre, mais l'entité locale (MTN Cameroon, Orange Cameroun) est une filiale de droit camerounais contrôlée par la maison-mère.

\begin{attention}[Exporter ne suffit pas !]
Une entreprise qui \emph{exporte} ses produits n'est pas forcément une FMN. Pour être qualifiée de multinationale, il faut une \textbf{implantation productive à l'étranger}, c'est-à-dire des filiales, des usines, des succursales contrôlées par la maison-mère. Une PME camerounaise qui exporte son café vers la France sans y posséder d'usine ni de filiale est une entreprise exportatrice, pas une FMN. Cette distinction est très souvent exigée dans les sujets de Bac.
\end{attention}

\courssub{La structure d'une FMN : maison-mère et filiales}

Pour visualiser cette organisation, construisons un croquis : la maison-mère au sommet, les filiales réparties sur les continents, et les flux qui les relient (capitaux, marchandises, informations).

\begin{popfigure}
\begin{center}
\begin{tikzpicture}[scale=1, every node/.style={transform shape}]
% Maison-mère
\node[draw=popdark, fill=popblueL, rounded corners, thick, minimum width=4.6cm, minimum height=1.1cm, align=center, font=\bfseries] (mere) at (0,0) {Maison-mère\\ \small (stratégie, R\&D, finances)};
% Filiales
\node[draw=popdark, fill=popgreenL, rounded corners, minimum width=2.6cm, minimum height=0.9cm, align=center] (eur) at (-5.4,-2.6) {Filiale Europe\\ \small production, vente};
\node[draw=popdark, fill=poporangeL, rounded corners, minimum width=2.6cm, minimum height=0.9cm, align=center] (asie) at (-1.8,-2.6) {Filiale Asie\\ \small assemblage};
\node[draw=popdark, fill=poppinkL, rounded corners, minimum width=2.6cm, minimum height=0.9cm, align=center] (afr) at (1.8,-2.6) {Filiale Afrique\\ \small matières premières};
\node[draw=popdark, fill=popgoldL, rounded corners, minimum width=2.6cm, minimum height=0.9cm, align=center] (ame) at (5.4,-2.6) {Filiale Amérique\\ \small vente, services};
% Flèches capitaux (descendantes)
\draw[-{Stealth[length=3mm]}, very thick, popblue] (mere.south west) -- (eur.north);
\draw[-{Stealth[length=3mm]}, very thick, popblue] (mere.south) -- (asie.north);
\draw[-{Stealth[length=3mm]}, very thick, popblue] (mere.south) -- (afr.north);
\draw[-{Stealth[length=3mm]}, very thick, popblue] (mere.south east) -- (ame.north);
% Flèches profits (remontantes, pointillés)
\draw[-{Stealth[length=3mm]}, thick, dashed, poporange] (eur.east) .. controls (-4,-1.2) .. (mere.west);
\draw[-{Stealth[length=3mm]}, thick, dashed, poporange] (ame.west) .. controls (4,-1.2) .. (mere.east);
% Flux entre filiales
\draw[-{Stealth[length=2.5mm]}, thick, popgreen] (afr.west) -- (asie.east);
\draw[-{Stealth[length=2.5mm]}, thick, popgreen] (asie.west) -- (eur.east);
% Légende des flèches
\node[align=left, font=\small] at (0,-4.6) {%
\begin{tabular}{@{}l@{}}
\textcolor{popblue}{$\longrightarrow$} Flux de capitaux et d'instructions (maison-mère $\rightarrow$ filiales)\\
\textcolor{poporange}{$\dashrightarrow$} Rapatriement des profits (filiales $\rightarrow$ maison-mère)\\
\textcolor{popgreen}{$\longrightarrow$} Flux de marchandises entre filiales (matières, composants)\\
\end{tabular}};
\end{tikzpicture}
\end{center}
\legende{Croquis de la structure d'une firme multinationale. La maison-mère concentre les décisions et la recherche ; les filiales exécutent la production et la vente sur chaque continent. Les profits remontent vers la maison-mère, ce qui explique la faible redistribution locale dans les pays d'accueil.}
\end{popfigure}

Ce croquis fait apparaître un point essentiel : la FMN fonctionne comme un \emph{réseau hiérarchisé}. La décision naît au sommet (la maison-mère), la production est dispersée (les filiales), et les profits sont rapatriés vers le siège. C'est exactement ce fonctionnement qui fait des FMN des acteurs à la fois puissants et contestés.

\courssub{Les stratégies des FMN : DIT, délocalisations et IDE}

\textbf{Pourquoi une entreprise s'implante-t-elle à l'étranger ?} L'intuition est simple : pour produire \emph{là où c'est le plus avantageux} et vendre \emph{là où c'est le plus rentable}. Cette logique se décline en trois stratégies complémentaires.

\begin{definition}[Investissements directs à l'étranger (IDE)]
Les \cle{IDE} sont des capitaux investis par une entreprise pour créer ou acquérir des activités durables dans un autre pays : construction d'une usine, achat d'une entreprise locale, création d'une filiale. Les IDE sont le principal instrument d'expansion des FMN. À l'échelle mondiale, ils représentent des flux de l'ordre de \num{1300} à \num{1500} milliards de dollars par an (ordre de grandeur, CNUCED).
\end{definition}

\textbf{1. La division internationale du travail (DIT).} La FMN découpe son processus de production en étapes et répartit ces étapes dans le monde selon les avantages de chaque territoire : la conception et la R\&D restent dans les pays riches (main-d'œuvre très qualifiée), l'assemblage va dans les pays à bas salaires (Asie du Sud-Est), l'extraction des matières premières dans les pays en développement (Afrique). Chaque territoire est spécialisé dans un maillon de la chaîne.

\textbf{2. Les délocalisations.} Une \cle{délocalisation} consiste à transférer tout ou partie de la production vers un pays où les coûts sont plus faibles : salaires modestes, fiscalité allégée, normes sociales et environnementales moins contraignantes. C'est ainsi que l'industrie textile européenne s'est largement déplacée vers le Bangladesh ou le Vietnam.

\textbf{3. La recherche des marchés et des ressources.} Les FMN s'implantent aussi pour être au plus près des consommateurs (MTN et Orange visent le marché camerounais des télécommunications) ou des matières premières (TotalEnergies pour les hydrocarbures, les négociants pour le cacao).

\begin{exemplebox}[La chaîne de valeur mondiale du smartphone]
Un smartphone vendu à Douala illustre parfaitement la DIT : le \emph{design} et les logiciels viennent des États-Unis ou de Chine ; les \emph{puces} sont fabriquées à Taïwan ou en Corée du Sud ; certains \emph{minerais} (cobalt, coltan) proviennent d'Afrique centrale ; l'\emph{assemblage} est réalisé en Chine ou au Vietnam ; la \emph{distribution} passe par des filiales et des opérateurs locaux. Aucun pays ne fabrique seul un smartphone : c'est le produit d'une chaîne de valeur mondiale organisée par des FMN.
\end{exemplebox}

\begin{savaistu}[Le cacao camerounais dans la chaîne de valeur mondiale]
Le Cameroun produit environ \num{280000} à \num{300000} tonnes de cacao par an (ordre de grandeur), principalement dans les régions du Centre, du Sud et du Sud-Ouest. Mais l'essentiel de cette fève est acheté et exporté par de grandes FMN du négoce et de la transformation, comme \textbf{Barry Callebaut} (Suisse) ou \textbf{Cargill} (États-Unis), présentes au Cameroun. La transformation en chocolat — l'étape la plus rentable — se fait surtout en Europe. Résultat : le planteur camerounais ne perçoit qu'une faible part (souvent moins de 10 \%) du prix final de la tablette de chocolat. Voilà la chaîne de valeur mondiale en action... et ses inégalités.
\end{savaistu}

\courssub{Le poids économique des FMN : des puissances comparables aux États}

Pour mesurer la puissance des FMN, comparons leur chiffre d'affaires annuel au produit intérieur brut (PIB) de quelques États africains. Le résultat est frappant.

\begin{popfigure}
\begin{center}
\begin{tikzpicture}
\begin{axis}[
    ybar, bar width=0.55cm,
    width=0.98\linewidth, height=7.2cm,
    ymin=0, ymax=700,
    ylabel={Milliards de dollars (ordres de grandeur)},
    symbolic x coords={Walmart, Amazon, Apple, Toyota, TotalEnergies, Côte d'Ivoire, Cameroun, Sénégal},
    xtick=data,
    x tick label style={rotate=30, anchor=east, font=\small},
    ytick={0,100,200,300,400,500,600,700},
    nodes near coords, nodes near coords style={font=\scriptsize},
    every axis plot/.append style={fill=popblue, draw=popdark}
]
\addplot+[ybar] coordinates {(Walmart,610) (Amazon,575) (Apple,385) (Toyota,275) (TotalEnergies,240) (Côte d'Ivoire,79) (Cameroun,45) (Sénégal,31)};
\end{axis}
\end{tikzpicture}
\end{center}
\legende{Comparaison entre le chiffre d'affaires annuel de cinq FMN et le PIB de trois États africains (Côte d'Ivoire, Cameroun, Sénégal), en milliards de dollars, ordres de grandeur arrondis (données 2022-2023, sources : classement Fortune Global 500, Banque mondiale). Le chiffre d'affaires de Walmart dépasse à lui seul le PIB cumulé du Cameroun, de la Côte d'Ivoire et du Sénégal.}
\end{popfigure}

\begin{exemplebox}[Un chiffre qui donne à réfléchir]
Le chiffre d'affaires de \textbf{Walmart} (environ 610 milliards de dollars) ou d'\textbf{Apple} (environ 385 milliards de dollars) dépasse le PIB cumulé de nombreux pays africains. Le PIB du Cameroun est de l'ordre de 45 milliards de dollars : une seule grande FMN « pèse » donc économiquement plus lourd que dix États comme le Cameroun. Attention toutefois à l'interprétation : un chiffre d'affaires (valeur des ventes) et un PIB (valeur ajoutée produite) ne mesurent pas exactement la même chose ; la comparaison donne un ordre de grandeur de la puissance économique, pas une équivalence stricte.
\end{exemplebox}

\courssub{Les FMN au Cameroun : observation et exemples}

Le territoire camerounais accueille de nombreuses filiales de FMN, que tu peux observer chaque jour à travers leurs logos, leurs enseignes et leurs produits :

\begin{center}
\begin{tabularx}{\linewidth}{|l|l|X|}
\hline
\rowcolor{popblueL} \textbf{FMN} & \textbf{Pays d'origine} & \textbf{Secteur d'activité au Cameroun} \\
\hline
TotalEnergies & France & Distribution de carburants, lubrifiants \\
\hline
MTN & Afrique du Sud & Téléphonie mobile, internet, mobile money \\
\hline
Orange & France & Télécommunications \\
\hline
Nestlé & Suisse & Agroalimentaire (bouillons, lait, boissons) \\
\hline
Castel (SABC) & France & Boissons (brasseries du Cameroun) \\
\hline
Bolloré (avec APM Terminals) & France & Gestion du terminal à conteneurs du port de Douala, chemin de fer (Camrail) \\
\hline
\end{tabularx}
\end{center}

\textbf{Commentaire de document guidé.} Face à une photographie du port de Douala montrant les portiques du terminal à conteneurs, voici la démarche attendue : (1) \emph{observer} — grues, conteneurs, navires ; (2) \emph{identifier} — le terminal est exploité par un consortium mené par Bolloré et APM Terminals (filiale du danois Maersk) ; (3) \emph{interpréter} — un maillon stratégique du commerce extérieur camerounais est géré par des FMN étrangères, ce qui montre l'insertion du pays dans les réseaux mondiaux, mais aussi sa dépendance.

\courssub{Poids et limites : l'influence des FMN et les critiques}

\begin{propriete}[Ce que les FMN apportent aux pays d'accueil]
Les FMN peuvent contribuer positivement : créations d'emplois, transferts de technologies et de savoir-faire, recettes fiscales, développement d'infrastructures (routes, ports, réseaux télécoms), intégration du pays dans les échanges mondiaux. Au Cameroun, les opérateurs télécoms ont par exemple permis l'essor du mobile money.
\end{propriete}

Mais cette puissance a un revers, que les sujets de Bac te demandent souvent d'analyser :

\begin{itemize}
\item \textbf{Influence sur les États} : une FMN dont le chiffre d'affaires dépasse le PIB du pays d'accueil dispose d'un rapport de force favorable lors des négociations (contrats miniers, concessions, fiscalité) ;
\item \textbf{Évasion et optimisation fiscales} : par des montages entre filiales (prix de transfert), certaines FMN déclarent leurs profits dans des pays à faible fiscalité, privant les États africains de recettes ;
\item \textbf{Faible redistribution locale} : les profits sont rapatriés vers la maison-mère ; les emplois créés sont souvent peu qualifiés et les activités à forte valeur ajoutée (R\&D, direction) restent au siège ;
\item \textbf{Dépendance économique} : quand des secteurs stratégiques (port, rail, télécoms, agroalimentaire) sont contrôlés par des capitaux étrangers, la souveraineté économique de l'État est réduite.
\end{itemize}

\begin{exoresolu}[Analyse d'un cas : les FMN au Cameroun, chance ou dépendance ?]
\textbf{Énoncé.} À partir de l'exemple de la gestion du terminal à conteneurs du port de Douala par un consortium étranger, dégagez en quelques lignes un avantage et une limite de la présence des FMN au Cameroun.
\tcblower
\textbf{Solution détaillée.} \emph{Avantage} : le consortium (Bolloré -- APM Terminals) a investi dans des équipements modernes (portiques, logiciels de gestion), ce qui a amélioré la fluidité du trafic portuaire, première porte d'entrée du commerce extérieur du Cameroun et de l'hinterland (Tchad, Centrafrique). La FMN apporte donc capitaux et technologie que l'État peinait à mobiliser seul. \emph{Limite} : les bénéfices de cette activité stratégique sont en partie rapatriés vers les maisons-mères (France, Danemark), et l'État camerounais ne contrôle pas pleinement un maillon vital de son économie. \emph{Conclusion} : la présence des FMN est ambivalente — elle intègre le Cameroun à la mondialisation tout en entretenant une dépendance ; l'enjeu pour l'État est de négocier des contrats qui maximisent les retombées locales (emplois, fiscalité, formation).
\end{exoresolu}

\courssub{Complément : la montée en puissance des FMN chinoises et des pays émergents}

Longtemps, les FMN étaient presque toutes originaires des pays riches (États-Unis, Europe, Japon). Depuis les années 2000, le paysage change : des FMN des \cle{pays émergents} s'imposent dans les classements mondiaux. La Chine en est le meilleur exemple : \textbf{Huawei} (télécommunications, très présente dans le réseau camerounais), \textbf{Transsion} (marques Tecno, Itel, Infinix, leaders du téléphone mobile en Afrique), ou encore des entreprises de BTP comme \textbf{China Harbour} (construction du port en eau profonde de Kribi). La maison-mère de MTN, sud-africaine, rappelle que les FMN ne sont plus l'apanage du seul « Nord ». Cette évolution illustre le basculement du centre de gravité de l'économie mondiale vers l'Asie — un point d'actualité apprécié dans une copie de Bac.

\begin{aretenir}[L'essentiel de la section]
\begin{itemize}
\item Une FMN = une maison-mère + des filiales implantées dans plusieurs pays + une stratégie mondiale ;
\item Ses stratégies : division internationale du travail, délocalisations, IDE, recherche des coûts les plus bas et des marchés ;
\item Elle organise les chaînes de valeur mondiales (smartphone, cacao) et constitue un vecteur majeur de la mondialisation ;
\item Son chiffre d'affaires peut dépasser le PIB d'États comme le Cameroun, d'où un rapport de force parfois déséquilibré ;
\item Apports (emplois, technologies, IDE) mais aussi critiques (évasion fiscale, faible redistribution, dépendance) ;
\item Nouvelle donne : l'essor des FMN chinoises et des pays émergents.
\end{itemize}
\end{aretenir}

\begin{exercice}[Pour t'entraîner]
1) Donne la définition d'une FMN et explique pourquoi une entreprise purement exportatrice n'en est pas une. 2) Construis un croquis simple montrant la chaîne de valeur mondiale du cacao camerounais, du planteur au consommateur européen, en faisant apparaître les FMN à chaque maillon. 3) À l'aide du graphique ci-dessus, rédige un commentaire de cinq lignes comparant la puissance des FMN à celle des États africains.
\end{exercice}

Désormais, tu sais que derrière chaque logo sur un produit se cache une organisation mondiale. Mais les FMN ne sont pas les seuls acteurs de la mondialisation : les États eux-mêmes, et leurs diasporas, jouent aussi un rôle actif. C'est l'objet de la section suivante.

\coursec{L'État et la diaspora}

Après avoir étudié les firmes multinationales, ces acteurs privés qui organisent la production à l'échelle planétaire, tu pourrais penser que la mondialisation échappe totalement aux États. Ce serait une erreur. Imagine la scène suivante : au port de Douala, un conteneur de cacao quitte le quai après avoir payé des droits de douane fixés par l'État camerounais, dans le cadre d'accords signés au sein de la CEMAC et de l'OMC. Quelques semaines plus tard, à Paris, un Camerounais installé en France envoie 200 euros à sa famille de Bafoussam pour financer la scolarité de sa petite sœur. Dans les deux cas, ce ne sont ni des entreprises ni des réseaux de transport qui agissent : ce sont un \cle{État} et des \cle{migrants}. Ce sont précisément les deux acteurs que nous allons étudier maintenant.

\courssub{Le rôle des États dans la libéralisation des échanges}

\begin{definition}[État et libéralisation des échanges]
L'\cle{État} est une institution politique souveraine qui exerce son autorité sur un territoire et une population. Dans la mondialisation, il joue un rôle moteur en \cle{libéralisant} les échanges, c'est-à-dire en supprimant ou en réduisant les obstacles (droits de douane, quotas, réglementations) qui freinent la circulation des marchandises, des capitaux et des services.
\end{definition}

L'intuition est simple : pour que les marchandises circulent librement entre les pays, encore faut-il que quelqu'un ouvre les portes. Or les portes, ce sont les frontières, et les clés sont entre les mains des États. Depuis la fin de la Seconde Guerre mondiale, les États ont volontairement et progressivement \textbf{abaissé leurs barrières douanières}. Sous l'égide du GATT (Accord général sur les tarifs douaniers et le commerce, signé en 1947) puis de l'OMC, les droits de douane moyens mondiaux sont passés d'environ 40 \% de la valeur des marchandises en 1947 à moins de 5 \% aujourd'hui (ordre de grandeur). Cette ouverture a permis la multiplication spectaculaire des échanges mondiaux.

Les États agissent de trois manières principales :

\begin{itemize}
\item \textbf{L'abaissement des barrières douanières} : les États réduisent les taxes à l'importation, ce qui rend les produits étrangers moins chers et stimule le commerce. Le Cameroun, par exemple, applique le tarif extérieur commun de la CEMAC, avec des droits de douane modulés selon les catégories de produits.
\item \textbf{La signature d'accords de libre-échange} : les États s'engagent, à deux (accords bilatéraux) ou à plusieurs (accords multilatéraux), à supprimer les obstacles commerciaux entre eux. Le Cameroun a signé un accord de partenariat économique avec l'Union européenne, qui ouvre progressivement son marché aux produits européens en échange d'un accès privilégié du cacao, de la banane ou du bois camerounais au marché européen.
\item \textbf{Les politiques d'attractivité des IDE} : pour attirer les \cle{investissements directs à l'étranger} (IDE), les États créent des zones franches, offrent des réductions d'impôts, construisent des infrastructures (ports, routes, électricité) et garantissent la stabilité politique. Le Cameroun a ainsi adopté une loi d'incitation à l'investissement privé et mis en place des zones économiques spéciales, comme celle du port en eau profonde de Kribi.
\end{itemize}

\courssub{Les organisations internationales créées ou soutenues par les États}

Pour organiser cette libéralisation, les États ne se contentent pas d'agir seuls : ils se regroupent au sein d'organisations internationales qui fixent les règles du jeu mondial.

\begin{definition}[L'OMC]
L'\cle{OMC} (Organisation mondiale du commerce), créée en 1995 et basée à Genève (Suisse), fixe les règles du commerce international et arbitre les litiges entre États. Elle succède au GATT et compte plus de 160 membres, dont le Cameroun. Son principe fondamental est la non-discrimination : un avantage accordé à un pays membre doit l'être à tous (clause de la nation la plus favorisée).
\end{definition}

Deux autres institutions, issues des accords de Bretton Woods (1944), complètent ce dispositif :

\begin{itemize}
\item Le \cle{FMI} (Fonds monétaire international) veille à la stabilité financière mondiale et vient en aide aux États en difficulté de paiement, en échange de réformes économiques (les fameux « programmes d'ajustement structurel » imposés au Cameroun dans les années 1980-1990).
\item La \cle{Banque mondiale} finance des projets de développement (routes, barrages, éducation, santé) dans les pays en développement.
\end{itemize}

Enfin, les États se regroupent en \cle{unions régionales} pour peser davantage dans la mondialisation. L'\textbf{Union européenne} (UE) est le modèle le plus abouti : marché unique, libre circulation des personnes et des capitaux, monnaie commune (l'euro) pour 20 de ses membres. Pour le Cameroun, deux organisations régionales sont essentielles :

\begin{itemize}
\item la \cle{CEMAC} (Communauté économique et monétaire de l'Afrique centrale), qui regroupe le Cameroun, le Tchad, le Gabon, le Congo, la République centrafricaine et la Guinée équatoriale, avec une monnaie commune, le franc CFA de la zone BEAC ;
\item la \cle{CEEAC} (Communauté économique des États de l'Afrique centrale), plus large (11 membres), qui vise l'intégration régionale de toute l'Afrique centrale.
\end{itemize}

\begin{attention}[La mondialisation ne signifie pas la disparition de l'État]
Erreur fréquente au Baccalauréat : affirmer que la mondialisation « supprime » les États ou les rend inutiles. C'est faux. L'État reste un acteur \textbf{central} : c'est lui qui négocie et signe les accords commerciaux, qui contrôle les frontières (visas, douanes), qui régule les marchés et qui protège certains secteurs jugés stratégiques. Même les États les plus libéraux, comme les États-Unis, imposent parfois des droits de douane élevés pour défendre leur industrie. La mondialisation transforme le rôle de l'État, elle ne l'abolit pas.
\end{attention}

\courssub{Le paradoxe de l'État mondialisé}

Nous touchons ici à un point délicat, souvent valorisé dans les sujets de synthèse : l'État favorise la mondialisation, mais celle-ci réduit en retour sa marge de manœuvre. C'est le \cle{paradoxe de l'État mondialisé}.

D'une part, l'État ouvre ses frontières pour attirer les capitaux et dynamiser son économie. D'autre part, cette ouverture le soumet à de nouvelles contraintes :

\begin{itemize}
\item \textbf{La concurrence fiscale} : pour attirer les firmes multinationales, les États se livrent à une véritable compétition à coups d'avantages fiscaux. Résultat : les recettes fiscales diminuent, et l'État dispose de moins de moyens pour financer ses services publics.
\item \textbf{La pression des marchés financiers} : les capitaux circulent instantanément. Si un État adopte une politique jugée « risquée » par les investisseurs, les capitaux fuient, la monnaie chute, et l'État est contraint de corriger sa politique. Les agences de notation (Standard \& Poor's, Moody's, Fitch) attribuent des notes aux États qui conditionnent leur accès aux emprunts.
\item \textbf{La dépendance aux règles internationales} : en adhérant à l'OMC ou en signant des programmes avec le FMI, l'État accepte de limiter sa propre liberté d'action économique.
\end{itemize}

\begin{aretenir}[Le paradoxe de l'État mondialisé]
L'État est à la fois \textbf{moteur} de la mondialisation (il libéralise les échanges, crée des organisations internationales, attire les IDE) et \textbf{contraint} par elle (concurrence fiscale, pression des marchés, règles de l'OMC). Il ne disparaît pas : il se transforme, passant d'un État « protecteur » à un État « régulateur » et « compétitif ».
\end{aretenir}

\courssub{La diaspora : un acteur humain de la mondialisation}

Changeons maintenant de regard. La mondialisation, ce ne sont pas seulement des conteneurs, des capitaux et des traités : ce sont aussi des \textbf{hommes et des femmes} qui traversent les frontières et maintiennent des liens vivaces avec leur pays d'origine.

\begin{definition}[La diaspora]
La \cle{diaspora} désigne l'ensemble des ressortissants d'un pays vivant durablement à l'étranger et maintenant des liens (familiaux, économiques, culturels, politiques) avec leur pays d'origine. Le mot vient du grec \emph{diaspeirein}, « disperser ».
\end{definition}

L'ampleur des migrations internationales est considérable : on compte environ 280 millions de migrants internationaux dans le monde (ordre de grandeur ONU, début des années 2020), soit près de 3,5 \% de la population mondiale. L'Afrique compte plusieurs dizaines de millions de ressortissants vivant hors de leur pays, dont une partie hors du continent.

La diaspora camerounaise est estimée à plusieurs centaines de milliers de personnes (certaines estimations officieuses dépassent le million, mais les chiffres précis sont difficiles à établir : retiens un ordre de grandeur prudent). Ses destinations principales sont :

\begin{itemize}
\item \textbf{l'Europe} : la France (première destination historique), l'Allemagne, la Belgique, le Royaume-Uni, l'Italie, la Suisse ;
\item \textbf{l'Amérique du Nord} : les États-Unis et le Canada, destinations en forte croissance depuis les années 1990 ;
\item \textbf{les pays africains voisins} : Gabon, Tchad, Nigéria, Afrique du Sud.
\end{itemize}

\begin{popfigure}
\includegraphics[width=\linewidth,height=9.5cm,keepaspectratio]{N11/migrant_stock.png}
\legende{Nombre total d'immigrés par pays de destination en 2020 (tous pays d'origine confondus ; échelle de teintes du jaune clair au bleu foncé, de moins de \num{10000} à plus de 10 millions de personnes). Les plus grands pays d'accueil apparaissent en bleu foncé : États-Unis, Allemagne, Royaume-Uni, Arabie saoudite. La carte ne représente ni les départs du Cameroun ni les remises : elle situe les grands pôles d'accueil des migrants dans le monde, parmi lesquels les pays d'Europe et d'Amérique du Nord où vit une large part de la diaspora camerounaise, selon le cours. \\ Source : Wikimedia Commons, « Migrant stock total, World, 2020 », Our World in Data (données ONU DAES), CC BY 4.0.}
\end{popfigure}

\courssub{Le poids économique des diasporas : les transferts d'argent}

L'apport le plus visible et le plus mesurable de la diaspora est financier : ce sont les \cle{transferts d'argent} des migrants, appelés aussi \cle{remises} (en anglais \emph{remittances}). Chaque mois, des millions de migrants envoient une partie de leurs revenus à leurs familles restées au pays, via des opérateurs comme Western Union, MoneyGram ou les applications de transfert mobile.

\begin{exemplebox}[Des remises plus importantes que l'aide publique]
Les remises des migrants vers l'Afrique subsaharienne dépassent \textbf{45 milliards de dollars par an} (ordre de grandeur Banque mondiale, début des années 2020), soit \textbf{davantage que l'aide publique au développement} (APD) reçue par le continent. Pour le Cameroun, les transferts de la diaspora représentent plusieurs centaines de millions de dollars par an (ordre de grandeur : entre 300 et 500 millions de dollars selon les années, soit environ 1 \% du PIB). Ces chiffres sont probablement sous-estimés, car une partie des transferts emprunte des circuits informels (argent confié à un voyageur, envois en nature).
\end{exemplebox}

\begin{popfigure}
\begin{center}
\begin{tikzpicture}
\begin{axis}[
  ybar, bar width=0.55cm,
  width=0.92\linewidth, height=6.2cm,
  symbolic x coords={2010,2015,2018,2020,2022},
  xtick=data,
  ymin=0, ymax=60,
  ylabel={Milliards de dollars},
  xlabel={Année},
  legend style={at={(0.02,0.97)}, anchor=north west, font=\footnotesize},
  ymajorgrids=true, grid style={popline!40},
  every axis plot/.append style={draw=none},
  nodes near coords, every node near coord/.append style={font=\tiny, rotate=90, anchor=west}
]
\addplot[fill=popblue] coordinates {(2010,30) (2015,35) (2018,46) (2020,42) (2022,53)};
\addplot[fill=poporange] coordinates {(2010,0.2) (2015,0.25) (2018,0.35) (2020,0.3) (2022,0.4)};
\legend{Afrique subsaharienne, Cameroun}
\end{axis}
\end{tikzpicture}
\end{center}
\legende{Évolution des transferts d'argent des migrants vers l'Afrique subsaharienne et vers le Cameroun (ordres de grandeur, en milliards de dollars courants ; source : Banque mondiale, valeurs arrondies et approximatives). On remarque la hausse tendancielle, la résistance des remises pendant la crise sanitaire de 2020, et le poids modeste mais croissant du Cameroun.}
\end{popfigure}

Que financent ces remises ? L'essentiel va à la \textbf{consommation des familles} : alimentation, frais de scolarité, soins de santé, loyers. Une partie sert aussi à des \textbf{investissements} : construction de maisons (les fameuses « maisons de la diaspora » qui transforment les quartiers de Douala, Yaoundé ou Bafoussam), achat de terrains, création de petites entreprises (salons de coiffure, cybercafés, commerces, élevages).

\begin{exoresolu}[Calculer et commenter un taux de variation]
\textbf{Énoncé.} Selon la Banque mondiale, les remises vers l'Afrique subsaharienne sont passées d'environ 30 milliards de dollars en 2010 à environ 53 milliards en 2022. 1) Calcule le taux de variation des remises entre 2010 et 2022. 2) Commente ce résultat en deux phrases, en le comparant à l'évolution de l'aide publique au développement, restée globalement stable sur la période.
\tcblower
\textbf{Solution détaillée.}
1) Le taux de variation se calcule ainsi :
\[ t = \frac{V_{2022} - V_{2010}}{V_{2010}} \times 100 = \frac{53 - 30}{30} \times 100 = \frac{23}{30} \times 100 \approx 76,7 \%. \]
Les remises ont donc augmenté d'environ \textbf{77 \%} en douze ans.
2) \textbf{Commentaire.} Les transferts d'argent des migrants vers l'Afrique subsaharienne ont connu une croissance spectaculaire, progressant de plus des trois quarts entre 2010 et 2022, alors que l'aide publique au développement stagnait. Cela montre que la diaspora est devenue le premier pourvoyeur de financements extérieurs du continent, devant les États donateurs : les migrations internationales constituent ainsi un moteur majeur et croissant de la mondialisation.
\end{exoresolu}

\courssub{Les autres apports de la diaspora}

Réduire la diaspora à un simple « distributeur de billets » serait une grave erreur d'analyse. Ses apports sont multiples :

\begin{itemize}
\item \textbf{Les transferts de compétences} : médecins, ingénieurs, enseignants et cadres camerounais formés ou expérimentés à l'étranger reviennent ponctuellement ou définitivement pour former, opérer, conseiller. On parle de « fuite des cerveaux » (\emph{brain drain}) quand ils partent, mais aussi de « circulation des compétences » quand ils transmettent leur savoir-faire, par exemple lors de missions médicales organisées par des associations de la diaspora.
\item \textbf{Les investissements productifs} : au-delà de l'immobilier, des membres de la diaspora créent des entreprises au Cameroun (agroalimentaire, technologies, hôtellerie), souvent en partenariat avec des proches restés au pays.
\item \textbf{Les réseaux commerciaux} : la diaspora constitue un relais commercial précieux. Les commerçants camerounais de Paris, Bruxelles ou New York écoulent des produits du pays (arachides, épices, tissus, produits artisanaux) et ouvrent des débouchés aux entreprises camerounaises.
\item \textbf{L'influence culturelle} : la musique camerounaise (makossa, bikutsi, afrobeat camerounais), la cuisine (ndolé, eru, poulet DG), la mode et la littérature rayonnent dans le monde grâce aux artistes et aux communautés de la diaspora. Cette « diplomatie culturelle » spontanée améliore l'image du Cameroun et nourrit le tourisme et les échanges.
\end{itemize}

\begin{savaistu}[Le Cameroun mobilise sa diaspora]
Conscient de ce potentiel, l'État camerounais a créé des dispositifs pour mobiliser sa diaspora : des \textbf{journées de la diaspora} sont organisées régulièrement à Yaoundé et dans les ambassades ; des incitations à l'investissement des Camerounais de l'étranger ont été prévues (facilités douanières pour les retours définitifs, possibilité de double nationalité en débat, guichets dédiés dans les représentations diplomatiques). Le gouvernement a même envisagé l'émission d'emprunts obligataires destinés spécialement à la diaspora pour financer des projets nationaux. La diaspora est ainsi officiellement reconnue comme un « partenaire du développement ».
\end{savaistu}

\begin{attention}[Pièges fréquents sur la diaspora]
Trois erreurs à éviter dans tes copies :
\begin{enumerate}
\item Ne confonds pas \cle{diaspora} et \cle{migration} : la diaspora est une communauté installée durablement à l'étranger et organisée en réseaux, pas un simple flux de départs.
\item N'affirme pas que les remises « enrichissent le pays » sans nuance : elles financent surtout la consommation des ménages, peuvent créer une dépendance et ne compensent pas toujours la perte de main-d'œuvre qualifiée (fuite des cerveaux).
\item Ne donne pas de chiffres précis inventés : les statistiques des remises sont des ordres de grandeur (Banque mondiale), car une partie des transferts est informelle. Écris « environ », « plus de », « de l'ordre de ».
\end{enumerate}
\end{attention}

\begin{methode}[Construire un croquis des flux de diaspora]
Pour cartographier la diaspora camerounaise au Baccalauréat, procède en étapes :
\begin{enumerate}
\item Trace un fond de carte très simplifié : trois blocs (Amérique du Nord, Europe, Afrique) suffisent.
\item Place le Cameroun par un point bien visible en Afrique centrale.
\item Dessine les flux migratoires par des flèches pleines dont l'épaisseur est proportionnelle à l'importance de la destination (Europe en premier, puis Amérique du Nord, puis Afrique).
\item Ajoute les remises financières par des flèches pointillées en sens inverse.
\item Soigne la légende : titre précis, signes conventionnels expliqués, flèche du nord, ordres de grandeur si possible.
\end{enumerate}
\end{methode}

\begin{exercice}[Pour t'entraîner]
1) Définis la diaspora et cite les trois principales zones d'installation de la diaspora camerounaise. 2) Explique en un paragraphe rédigé le paradoxe de l'État mondialisé. 3) À l'aide du graphique des remises, montre en cinq lignes que la diaspora est devenue un acteur financier majeur de la mondialisation.
\end{exercice}

En résumé, l'État et la diaspora sont deux acteurs complémentaires de la mondialisation : le premier fixe le cadre juridique et politique des échanges tout en voyant sa souveraineté mise sous tension ; la seconde fait circuler hommes, capitaux, compétences et cultures entre le pays d'origine et les pays d'accueil. Reste à examiner les acteurs de la solidarité internationale — associations et ONG — puis à synthétiser l'ensemble à travers le commerce international.

\coursec{Les associations et les ONG}

Dans la section précédente, nous avons vu que l'État, malgré la mondialisation, conserve un rôle central et que la diaspora camerounaise, par ses transferts d'argent et ses réseaux, relie le pays au reste du monde. Mais il existe d'autres acteurs qui traversent les frontières sans être ni des entreprises, ni des États : ce sont les \cle{associations} et les \cle{ONG}. Quand, en 2017, des milliers de personnes fuient les violences dans les régions du Nord-Ouest et du Sud-Ouest, ce ne sont pas seulement les pouvoirs publics qui leur portent secours : des équipes de Médecins Sans Frontières, de la Croix-Rouge ou d'organisations locales distribuent soins, eau et abris. Qui sont ces acteurs ? Comment agissent-ils ? Et quelle place occupent-ils dans la mondialisation ? C'est l'objet de cette section.

\courssub{Définition et nature des ONG}

\begin{definition}[Qu'est-ce qu'une ONG ?]
Une \cle{ONG} (organisation non gouvernementale) est une organisation \textbf{privée}, \textbf{à but non lucratif} et \textbf{indépendante des États}, qui mène des actions d'intérêt général (humanitaire, développement, environnement, droits humains) à l'échelle nationale ou internationale.
\end{definition}

Décortiquons cette définition, car chaque mot compte au Baccalauréat :

\begin{itemize}
\item \textbf{Privée et non gouvernementale} : une ONG n'est créée ni dirigée par un État. Elle naît de l'initiative de citoyens regroupés autour d'une cause. C'est ce qui la distingue radicalement des organisations internationales.
\item \textbf{À but non lucratif} : contrairement à une firme multinationale, une ONG ne recherche pas le profit. Ses ressources (dons, cotisations, subventions) servent à financer ses actions, non à rémunérer des actionnaires.
\item \textbf{D'intérêt général} : elle défend une cause qui dépasse les intérêts privés de ses membres : soigner, nourrir, protéger la nature, défendre les droits.
\end{itemize}

\begin{attention}[Ne pas confondre ONG et organisations internationales]
C'est l'erreur la plus fréquente dans les copies. Les \textbf{organisations internationales} (ONU, OMC, Banque mondiale, Union africaine) sont créées par des \textbf{États} qui signent des traités et y siègent : on dit qu'elles sont \emph{intergouvernementales}. Les \textbf{ONG} sont créées par des \textbf{personnes privées}. Ainsi, l'UNICEF est une agence de l'ONU (organisation internationale), tandis que Médecins Sans Frontières est une ONG. Dire « l'ONG ONU » est une faute grave de géographie.
\end{attention}

\courssub{La typologie des ONG : quatre grandes familles}

Toutes les ONG ne font pas la même chose. Pour s'y retrouver, on les classe selon leur domaine d'action. Retiens bien un exemple par famille : c'est ce que le correcteur attend.

\begin{itemize}
\item \textbf{Les ONG humanitaires} interviennent en urgence lors de crises : guerres, famines, épidémies, catastrophes naturelles. Exemples : \textbf{Médecins Sans Frontières} (MSF, fondée en France en 1971, prix Nobel de la paix 1999) et la \textbf{Croix-Rouge}.
\item \textbf{Les ONG de développement} agissent sur le long terme pour améliorer les conditions de vie : accès à l'eau potable, éducation, microcrédit, agriculture. Exemple : \textbf{Oxfam}, qui lutte contre la pauvreté et les inégalités.
\item \textbf{Les ONG environnementales} défendent la nature et le climat. Exemples : \textbf{Greenpeace} (actions spectaculaires contre la déforestation ou la pêche industrielle) et le \textbf{WWF} (World Wide Fund for Nature, protection des espèces menacées).
\item \textbf{Les ONG de défense des droits humains} dénoncent les violations des libertés : tortures, emprisonnements arbitraires, discriminations. Exemple : \textbf{Amnesty International} (fondée en 1961, prix Nobel de la paix 1977).
\end{itemize}

\begin{popfigure}
\begin{center}
\begin{tikzpicture}[
  grow=right,
  level 1/.style={sibling distance=2.6cm, level distance=3.6cm},
  level 2/.style={sibling distance=1.15cm, level distance=3.4cm},
  racine/.style={rectangle, rounded corners, draw=popdark, fill=popblue, text=white, font=\bfseries, align=center, inner sep=5pt},
  branche/.style={rectangle, rounded corners, draw=popdark, fill=popblueL, font=\small\bfseries, align=center, inner sep=4pt},
  feuille/.style={rectangle, rounded corners, draw=popmuted, fill=popcream, font=\small, align=center, inner sep=3pt},
  edge from parent/.style={draw=popline, thick}
]
\node[racine, align=center]{ONG\\ (à but non lucratif)}
  child { node[branche, align=center]{Droits\\ humains}
    child { node[feuille, align=center]{Amnesty\\ International} }
  }
  child { node[branche] {Environnement}
    child { node[feuille] {WWF} }
    child { node[feuille] {Greenpeace} }
  }
  child { node[branche] {Développement}
    child { node[feuille] {Oxfam} }
  }
  child { node[branche, align=center]{Humanitaires\\ (urgence)}
    child { node[feuille] {Croix-Rouge} }
    child { node[feuille] {MSF} }
  };
\end{tikzpicture}
\end{center}
\legende{Schéma 1 : typologie des ONG en quatre grandes familles, avec un ou deux exemples par branche. Chaque famille correspond à un domaine d'action : l'urgence, le long terme, la nature, les libertés.}
\end{popfigure}

\begin{aretenir}[L'essentiel sur les ONG]
Les ONG sont des acteurs \textbf{privés, sans but lucratif et indépendants des États}. On distingue quatre familles : humanitaires (MSF, Croix-Rouge), de développement (Oxfam), environnementales (Greenpeace, WWF) et de défense des droits humains (Amnesty International). Elles complètent — et parfois concurrencent — l'action des États.
\end{aretenir}

\courssub{Des acteurs transnationaux de la mondialisation}

Pourquoi étudie-t-on les ONG dans une leçon sur les \textbf{facteurs de la mondialisation} ? Parce que les grandes ONG fonctionnent exactement comme des réseaux planétaires.

\begin{propriete}[Les ONG, acteurs transnationaux]
Les grandes ONG internationales sont présentes sur \textbf{tous les continents} grâce à des réseaux de sections nationales, de bénévoles et de partenaires locaux. Elles mènent des \textbf{campagnes mondiales} coordonnées (pétitions, rapports, médiatisation) qui circulent instantanément grâce aux télécommunications et à Internet. Elles tissent ainsi des liens transfrontaliers indépendants des États et des entreprises : on parle de \cle{société civile mondiale}.
\end{propriete}

Prenons un raisonnement concret. MSF compte des sections dans plus de vingt pays et intervient dans environ 70 pays (ordre de grandeur). Amnesty International revendique plusieurs millions de membres et sympathisants dans plus de 150 pays. Quand Greenpeace lance une campagne contre la déforestation, la même affiche, le même slogan, la même pétition apparaissent simultanément à Paris, à São Paulo et à Yaoundé. Les ONG sont donc des \textbf{vecteurs de circulation} : elles font circuler des idées, des normes (droits humains, protection de l'environnement), des financements et des hommes à l'échelle du globe. À ce titre, elles sont à la fois un \textbf{produit} de la mondialisation (elles profitent des transports et d'Internet) et un de ses \textbf{moteurs} (elles mondialisent des causes).

\courssub{Le rôle des ONG au Cameroun}

Le Cameroun est un terrain d'action important pour de nombreuses ONG, nationales et internationales. Trois domaines se détachent.

\textbf{1. L'action humanitaire dans les zones de crise.} Dans l'Extrême-Nord, les exactions du groupe Boko Haram ont provoqué, depuis 2014, le déplacement de centaines de milliers de personnes (plus de 300 000 déplacés internes selon les estimations de l'OCHA, ordre de grandeur), auxquels s'ajoutent des réfugiés nigérians accueillis notamment dans le camp de \textbf{Minawao} (département du Mayo-Tsanaga, région de Mokolo). Dans les régions du Nord-Ouest et du Sud-Ouest, la crise anglophone a elle aussi déplacé des centaines de milliers de personnes. Des ONG comme MSF, la Croix-Rouge camerounaise ou le Norwegian Refugee Council y distribuent soins, eau, nourriture et abris, souvent en coordination avec les agences de l'ONU (HCR, PAM).

\textbf{2. Les projets de développement rural.} De nombreuses ONG de développement appuient les paysans camerounais : forage de puits, construction de salles de classe et de centres de santé, formation aux techniques agricoles, appui aux groupements de femmes, microcrédits. Ces actions complètent les programmes de l'État dans les zones rurales enclavées.

\textbf{3. La protection de la forêt du bassin du Congo.} Le Cameroun abrite une partie du deuxième massif forestier du monde. Des ONG environnementales y mènent des actions de conservation aux côtés de l'État.

\begin{exemplebox}[Le WWF et les parcs camerounais]
Le \textbf{WWF} intervient depuis des décennies dans le sud-est du Cameroun, notamment autour du parc national de \textbf{Lobéké} et de la réserve de faune du \textbf{Dja} (classée au patrimoine mondial de l'UNESCO en 1987). Ses équipes appuient le \textbf{MINFOF} (ministère des Forêts et de la Faune) : formation et équipement des écogardes contre le braconnage (notamment des éléphants), suivi scientifique de la faune, sensibilisation des populations riveraines et appui à des activités alternatives (agriculture durable, écotourisme). Greenpeace, de son côté, dénonce régulièrement l'exploitation forestière non durable dans le bassin du Congo. Ces exemples montrent qu'une ONG peut \textbf{coopérer} avec l'État tout en restant indépendante de lui.
\end{exemplebox}

\courssub{Les associations altermondialistes : critiquer la mondialisation sans la refuser}

Toutes les associations ne se contentent pas de soigner ou de planter des arbres : certaines veulent transformer la mondialisation elle-même.

\begin{definition}[L'altermondialisme]
L'\cle{altermondialisme} est un mouvement qui \textbf{ne refuse pas la mondialisation} mais en réclame \textbf{une autre forme}, plus juste, plus solidaire et plus respectueuse des hommes et de l'environnement. Il critique la mondialisation libérale actuelle, jugée inégalitaire et dominée par les firmes multinationales et les grandes puissances.
\end{definition}

\begin{attention}[Altermondialiste n'est pas antimondialiste]
Ne rédige jamais « les altermondialistes sont contre la mondialisation » : c'est faux. Le préfixe \emph{alter-} signifie « autre » : ils veulent une \emph{autre} mondialisation. Ils utilisent d'ailleurs eux-mêmes les outils de la mondialisation (Internet, voyages internationaux, réseaux planétaires) pour la critiquer.
\end{attention}

Le rendez-vous symbolique de ce mouvement est le \textbf{Forum social mondial} (FSM), créé en 2001 à Porto Alegre (Brésil), en réponse au Forum économique mondial de Davos qui réunit chaque année chefs d'entreprise et dirigeants politiques. Le FSM rassemble ONG, syndicats, associations de paysans, mouvements de jeunesse autour du slogan « Un autre monde est possible ». Ses revendications typiques : l'annulation de la dette des pays pauvres, un commerce plus équitable, la taxation des transactions financières, la justice climatique. Des forums sociaux régionaux et nationaux, y compris africains, prolongent cette dynamique.

\courssub{Limites et critiques des ONG}

Être rigoureux, c'est aussi savoir nuancer. Les ONG font l'objet de critiques qu'il faut connaître pour construire un devoir équilibré.

\begin{itemize}
\item \textbf{La dépendance financière} : beaucoup d'ONG vivent de subventions publiques ou de bailleurs de fonds étrangers. Cette dépendance peut limiter leur indépendance réelle et orienter leurs priorités vers les causes « visibles » qui attirent les dons.
\item \textbf{La concurrence avec les États} : en se substituant aux pouvoirs publics (santé, éducation, eau), certaines ONG peuvent affaiblir la construction des services publics locaux et créer une dépendance durable. Certains États leur reprochent aussi l'ingérence dans leurs affaires intérieures.
\item \textbf{Une efficacité débattue} : malgré des décennies d'aide au développement, la pauvreté persiste dans de nombreuses régions. On reproche parfois aux ONG des projets de courte durée, mal coordonnés entre eux, ou une méconnaissance des réalités locales.
\item \textbf{Des moyens limités} : à l'échelle mondiale, le poids financier des ONG reste faible comparé à celui des firmes multinationales ou des États.
\end{itemize}

\begin{exoresolu}[Commentaire guidé : les ONG, moteur ou frein de la mondialisation ?]
\textbf{Énoncé.} Un élève écrit : « Les ONG sont des organisations créées par les États pour aider les pays pauvres ; elles freinent la mondialisation. » Relevez et corrigez les erreurs contenues dans cette phrase, puis montrez en quoi les ONG sont un facteur de la mondialisation.

\tcblower
\textbf{Solution détaillée.}

\emph{Étape 1 : correction des erreurs.}
\begin{enumerate}
\item « Créées par les États » : \textbf{faux}. Les ONG sont des organisations \textbf{privées}, nées de l'initiative de citoyens, indépendantes des États. Ce sont les organisations internationales (ONU, OMC) qui sont créées par des États.
\item « Elles freinent la mondialisation » : \textbf{faux}. Même les associations altermondialistes ne refusent pas la mondialisation : elles en réclament une autre forme. Les ONG sont au contraire des acteurs de la mondialisation.
\end{enumerate}

\emph{Étape 2 : les ONG, facteur de la mondialisation.}
\begin{enumerate}
\item Elles forment des \textbf{réseaux transnationaux} présents sur tous les continents (MSF dans environ 70 pays, Amnesty dans plus de 150 pays).
\item Elles font \textbf{circuler} à l'échelle mondiale des idées, des normes (droits humains, environnement), des financements et des personnes, grâce aux télécommunications et aux transports.
\item Elles agissent concrètement dans les territoires : au Cameroun, aide humanitaire dans l'Extrême-Nord et les régions anglophones, projets ruraux, protection des parcs du Dja et de Lobéké avec le WWF.
\end{enumerate}

\emph{Conclusion.} Les ONG sont des acteurs privés, non lucratifs et indépendants des États, qui à la fois profitent de la mondialisation et l'alimentent : elles constituent l'un de ses moteurs, tout en en dénonçant parfois les dérives.
\end{exoresolu}

\begin{exercice}[À toi de jouer]
\begin{enumerate}
\item Recopie et complète : « Une ONG est une organisation \ldots, à but \ldots, indépendante des \ldots »
\item Classe les organisations suivantes selon la famille d'ONG à laquelle elles appartiennent : Greenpeace, MSF, Amnesty International, Oxfam.
\item Explique en cinq lignes la différence entre une ONG et une organisation internationale, en donnant un exemple de chaque.
\item Rédige un paragraphe argumenté (10 lignes) montrant le rôle des ONG au Cameroun, avec au moins deux exemples localisés.
\end{enumerate}
\end{exercice}

\begin{corrige}[Exercice : À toi de jouer]
\begin{enumerate}
\item « Une ONG est une organisation \textbf{privée (non gouvernementale)}, à but \textbf{non lucratif}, indépendante des \textbf{États} ».
\item Greenpeace : ONG environnementale ; MSF : ONG humanitaire ; Amnesty International : ONG de défense des droits humains ; Oxfam : ONG de développement.
\item Une ONG est créée par des personnes privées regroupées autour d'une cause (exemple : MSF) ; une organisation internationale est créée par des États liés par un traité (exemple : l'ONU ou l'OMC). Les secondes sont donc intergouvernementales, les premières non gouvernementales.
\item Éléments attendus : action humanitaire dans l'Extrême-Nord (camp de Minawao, déplacés de la crise Boko Haram) et dans le Nord-Ouest/Sud-Ouest ; projets de développement rural (puits, écoles, microcrédits) ; protection de la forêt du bassin du Congo avec le WWF aux côtés du MINFOF dans les parcs du Dja et de Lobéké. Conclusion nuancée : les ONG complètent l'État mais ne peuvent s'y substituer durablement.
\end{enumerate}
\end{corrige}

\begin{savaistu}[Le prix Nobel de la paix, un habitué des ONG]
Plusieurs ONG ont reçu le prix Nobel de la paix : Amnesty International en 1977, Médecins Sans Frontières en 1999, et la Campagne internationale pour l'abolition des armes nucléaires (ICAN) en 2017. Cela montre la reconnaissance mondiale du rôle de ces acteurs privés dans les grandes questions planétaires — un rôle qui, il y a un siècle, revenait presque exclusivement aux États.
\end{savaistu}

En résumé, après les États et la diaspora, les associations et les ONG apparaissent comme des acteurs à part entière de la mondialisation : privés, solidaires, organisés en réseaux planétaires, ils soignent, développent, protègent et protestent. Mais aucun de ces acteurs ne pèserait pleinement sans les flux de marchandises qui structurent l'économie mondiale : c'est le \textbf{commerce international}, dernière pièce de notre leçon, que nous allons maintenant étudier en synthèse.

\coursec{Le commerce international et synthèse}

Après avoir vu comment les associations et les ONG tissent des réseaux de solidarité à l'échelle de la planète, il nous reste à étudier le plus ancien et le plus visible des moteurs de la mondialisation : le \cle{commerce international}. C'est lui qui fait voyager le cacao de Kumba jusqu'aux chocolateries d'Europe, le pétrole de Kribi jusqu'aux raffineries d'Asie, et qui fait arriver sur les étals du marché Mokolo des téléphones fabriqués en Asie et du riz importé. Cette dernière section permettra aussi de faire la \cle{synthèse} de la leçon : comment les cinq facteurs de la mondialisation s'articulent-ils et se renforcent-ils mutuellement ?

\courssub{Définition et mesure du commerce international}

\begin{definition}[Commerce international]
Le \cle{commerce international} désigne l'ensemble des échanges de biens (marchandises) et de services entre les pays du monde. Pour un pays donné, on distingue les \cle{exportations} (biens et services vendus à l'étranger) et les \cle{importations} (biens et services achetés à l'étranger).
\end{definition}

\begin{definition}[La balance commerciale]
La \cle{balance commerciale} est la différence entre la valeur des exportations et celle des importations d'un pays sur une période donnée (généralement une année) :
\[ \text{Balance commerciale} = \text{Valeur des exportations} - \text{Valeur des importations} \]
Si le résultat est positif, la balance est \cle{excédentaire} (le pays vend plus qu'il n'achète) ; s'il est négatif, elle est \cle{déficitaire} (le pays achète plus qu'il ne vend). On calcule aussi le \cle{taux de couverture} :
\[ \text{Taux de couverture} = \frac{\text{Valeur des exportations}}{\text{Valeur des importations}} \times 100 \]
Un taux inférieur à 100 \% signifie que les exportations ne suffisent pas à payer les importations.
\end{definition}

\begin{exemplebox}[Un exemple simple pour bien comprendre]
Imaginons un petit pays qui exporte pour 10 milliards de dollars de cacao et de pétrole, et qui importe pour 14 milliards de dollars de machines, de carburant raffiné et de riz. Sa balance commerciale est de $10 - 14 = -4$ milliards de dollars : elle est déficitaire, et son taux de couverture est de $\frac{10}{14} \times 100 \approx 71\,\%$. Pour payer la différence, le pays doit puiser dans ses réserves de devises ou s'endetter. C'est exactement la situation que connaissent de nombreux pays africains, dont parfois le Cameroun.
\end{exemplebox}

\begin{attention}[Ne pas confondre]
Ne confonds pas la \emph{balance commerciale} (qui ne concerne que les biens, parfois les services) avec la \emph{balance des paiements} (qui enregistre toutes les opérations avec l'étranger : commerce, mais aussi transferts d'argent de la diaspora, investissements, aides). Au Bac, cette confusion coûte des points.
\end{attention}

\courssub{L'explosion des échanges depuis 1950}

Observe un fait remarquable : depuis la fin de la Seconde Guerre mondiale, le commerce mondial a crû beaucoup plus vite que la production mondiale elle-même. Autrement dit, les pays ne produisent pas seulement davantage : ils échangent une part de plus en plus grande de ce qu'ils produisent.

\begin{popfigure}
\begin{center}
\begin{tikzpicture}
\begin{axis}[
  width=0.95\linewidth, height=7.5cm,
  xlabel={Année}, ylabel={Indice (base 100 en 1950)},
  xmin=1950, xmax=2020, ymin=0, ymax=6200,
  xtick={1950,1960,1970,1980,1990,2000,2010,2020},
  ytick={0,1000,2000,3000,4000,5000,6000},
  legend pos=north west, legend style={font=\small},
  grid=both, grid style={popline!40},
]
\addplot[popcurve, popblue, very thick, domain=1950:2020, samples=60] {100*exp(0.033*(x-1950))};
\addlegendentry{Production mondiale (PIB réel, ordre de grandeur)}
\addplot[popcurve, poporange, very thick, domain=1950:2020, samples=60] {100*exp(0.058*(x-1950))};
\addlegendentry{Commerce mondial de marchandises (volume, ordre de grandeur)}
\end{axis}
\end{tikzpicture}
\end{center}
\legende{Évolution comparée de la production mondiale et du commerce mondial depuis 1950 (indices base 100, ordres de grandeur). Le volume du commerce mondial a été multiplié par plus de 40 depuis 1950 (indice voisin de \num{5 800} en 2020), tandis que la production mondiale était multipliée par environ 10 (indice voisin de \num{1 000}) : le commerce croît environ deux fois plus vite que la production. Source : ordres de grandeur OMC.}
\end{popfigure}

Comment expliquer cet envol ? C'est ici que tous les facteurs étudiés dans cette leçon entrent en jeu : la \emph{containerisation} et la baisse spectaculaire du coût des transports, la libéralisation orchestrée par le GATT puis l'OMC, les firmes multinationales qui découpent leurs chaînes de production entre plusieurs pays (un même smartphone traverse plusieurs frontières avant d'être assemblé), et les télécommunications qui permettent de piloter ces flux en temps réel.

\begin{aretenir}[L'idée-force]
Depuis 1950, le commerce mondial croît plus vite que la production mondiale : l'économie mondiale devient de plus en plus \emph{ouverte} et \emph{intégrée}. C'est le signe le plus mesurable de la mondialisation.
\end{aretenir}

\courssub{Une géographie inégale des échanges}

Le commerce mondial n'est pas également réparti. Trois grandes observations s'imposent.

\textbf{1. La domination historique de la Triade.} L'Amérique du Nord, l'Europe occidentale et l'Asie orientale (Japon, puis Chine et « dragons ») concentrent l'essentiel des échanges mondiaux de marchandises et de services. Les flux les plus intenses relient ces pôles entre eux (axe transatlantique, axe transpacifique, axe Europe--Asie).

\textbf{2. La montée en puissance de la Chine.} Devenue « l'atelier du monde », la Chine est aujourd'hui le premier exportateur mondial de marchandises (depuis la fin des années 2000) et le premier partenaire commercial de la plupart des pays du monde.

\textbf{3. La marginalisation de l'Afrique.}

\begin{exemplebox}[Un chiffre qui doit frapper]
L'Afrique ne pèse qu'environ \cle{2 à 3 \%} du commerce mondial de marchandises, alors qu'elle rassemble environ \cle{17 \%} de la population mondiale. Autrement dit, sa part dans les échanges est cinq à huit fois inférieure à son poids démographique. Ce paradoxe s'explique par la faiblesse de l'industrialisation, la spécialisation dans les produits primaires aux prix fluctuants, le coût élevé des transports et la faiblesse des échanges intra-africains (moins de 20 \% du commerce africain se fait entre pays africains, ordre de grandeur).
\end{exemplebox}

\courssub{Le commerce du Cameroun : un exemple de pays périphérique}

Le Cameroun illustre parfaitement la place des pays en développement dans le commerce mondial : il \emph{exporte des produits primaires bruts ou peu transformés} et \emph{importe des produits manufacturés et des biens de consommation}.

\begin{popfigure}
\begin{center}
\begin{tikzpicture}[scale=1]
% Diagramme circulaire : exportations du Cameroun (ordres de grandeur)
\def\r{2.2}
% Pétrole brut ~40%
\fill[popblue] (0,0) -- (90:\r) arc (90:90-144:\r) -- cycle;
% Cacao ~15%
\fill[poporange] (0,0) -- (90-144:\r) arc (90-144:90-198:\r) -- cycle;
% Bois ~12%
\fill[popgreen] (0,0) -- (90-198:\r) arc (90-198:90-241:\r) -- cycle;
% Coton ~7%
\fill[poppurple] (0,0) -- (90-241:\r) arc (90-241:90-266:\r) -- cycle;
% Bananes ~5%
\fill[popgold] (0,0) -- (90-266:\r) arc (90-266:90-284:\r) -- cycle;
% Autres ~21%
\fill[popmuted!60] (0,0) -- (90-284:\r) arc (90-284:90-360:\r) -- cycle;
% Légendes
\node[anchor=west, font=\small] at (3.0,1.6) {\textcolor{popblue}{\rule{0.4cm}{0.25cm}} Pétrole brut : environ 40 \%};
\node[anchor=west, font=\small] at (3.0,1.0) {\textcolor{poporange}{\rule{0.4cm}{0.25cm}} Cacao : environ 15 \%};
\node[anchor=west, font=\small] at (3.0,0.4) {\textcolor{popgreen}{\rule{0.4cm}{0.25cm}} Bois : environ 12 \%};
\node[anchor=west, font=\small] at (3.0,-0.2) {\textcolor{poppurple}{\rule{0.4cm}{0.25cm}} Coton : environ 7 \%};
\node[anchor=west, font=\small] at (3.0,-0.8) {\textcolor{popgold}{\rule{0.4cm}{0.25cm}} Bananes : environ 5 \%};
\node[anchor=west, font=\small] at (3.0,-1.4) {\textcolor{popmuted}{\rule{0.4cm}{0.25cm}} Autres (aluminium, café, caoutchouc...) : environ 21 \%};
\end{tikzpicture}
\end{center}
\legende{Structure approximative des exportations camerounaises par produit (ordres de grandeur, années 2010--2020 ; les parts varient fortement selon le cours du pétrole). Les produits primaires dominent très largement. Source : ordres de grandeur INS / Douanes camerounaises.}
\end{popfigure}

\textbf{Les exportations} sont dominées par le \cle{pétrole brut} (exploité notamment au large de Rio del Rey et de Kribi, raffiné en partie à Limbe par la SONARA), le \cle{cacao} (le Cameroun figure parmi les premiers producteurs africains, derrière la Côte d'Ivoire et le Ghana), le \cle{bois}, le \cle{coton} (Nord, via la SODECOTON), les \cle{bananes} (plantations du littoral, PHP) et l'aluminium (ALUCAM, à Edéa).

\textbf{Les importations} comprennent les produits manufacturés (machines, véhicules, appareils électriques et électroniques, médicaments), les carburants raffinés (paradoxalement, le Cameroun exporte du pétrole brut mais importe une partie de ses carburants, la capacité de raffinage de la SONARA étant insuffisante), et les céréales (blé, riz), dont la consommation urbaine croît rapidement.

\textbf{Les partenaires} ont profondément changé : longtemps dominés par la France et l'Union européenne, ils sont aujourd'hui marqués par l'ascension de la \cle{Chine}, devenue premier ou deuxième partenaire du Cameroun selon les années (fournisseur de produits manufacturés, client du pétrole et du bois). L'Union européenne (Pays-Bas pour le cacao, France, Espagne, Italie) reste un débouché majeur.

\begin{savaistu}[Le basculement vers la Chine]
La Chine est devenue le premier partenaire commercial de la plupart des pays africains, dont le Cameroun, devançant les anciennes puissances coloniales. Les échanges sino-africains ont été multipliés par plus de 20 entre 2000 et 2020 (ordre de grandeur). Ce basculement illustre la recomposition de la géographie du commerce mondial : le centre de gravité des échanges glisse vers l'Asie.
\end{savaistu}

\begin{exoresolu}[Calculer et commenter une balance commerciale]
\textbf{Énoncé.} Une année donnée, le Cameroun exporte pour environ \num{2 500} milliards de FCFA et importe pour environ \num{3 500} milliards de FCFA (ordres de grandeur). 1) Calcule le solde de la balance commerciale. 2) Calcule le taux de couverture (exportations / importations $\times$ 100). 3) Commente en trois lignes.
\tcblower
\textbf{Solution.}
\begin{align*}
\text{Solde} &= 2500 - 3500 = -1000 \text{ milliards de FCFA} \\
\text{Taux de couverture} &= \frac{2500}{3500} \times 100 \approx 71 \%
\end{align*}
\textbf{Commentaire.} La balance commerciale du Cameroun est déficitaire d'environ \num{1 000} milliards de FCFA : le pays achète à l'étranger bien plus qu'il ne vend. Les exportations ne couvrent qu'environ 71 \% des importations. Ce déficit traduit une économie peu industrialisée, dépendante des produits primaires pour ses ventes et des produits manufacturés importés pour sa consommation et ses équipements.
\end{exoresolu}

\courssub{Synthèse : cinq facteurs qui se renforcent mutuellement}

La mondialisation n'est pas le produit d'un facteur isolé, mais d'un \emph{système} où chaque facteur alimente les autres :

\begin{itemize}
\item Les \cle{transports} (conteneurs, porte-conteneurs géants, avions-cargos) et les \cle{télécommunications} (Internet, câbles sous-marins, satellites) réduisent les coûts et les délais : ils rendent les échanges physiquement possibles à grande échelle.
\item Les \cle{firmes multinationales} organisent la production à l'échelle planétaire : elles utilisent précisément ces transports et ces réseaux pour découper leurs filières entre plusieurs pays.
\item Les \cle{États} libéralisent (OMC, accords de libre-échange, privatisation des ports et des télécoms) et s'appuient sur leurs \cle{diasporas}, dont les transferts d'argent et les réseaux d'affaires facilitent les flux.
\item Les \cle{associations et ONG} tissent des réseaux transnationaux de solidarité, de défense des droits et de régulation, qui accompagnent ou contestent la mondialisation.
\item Le \cle{commerce international} est à la fois la \emph{conséquence} de ces facteurs et leur \emph{moteur} : plus on échange, plus on investit dans les ports, les câbles et les réseaux.
\end{itemize}

\begin{popfigure}
\includegraphics[width=\linewidth,height=9.5cm,keepaspectratio]{N11/dev_emerging.png}
\legende{Fond de carte : les marchés développés (rose : Amérique du Nord, Europe, Japon, Australie, Nouvelle-Zélande), les marchés émergents (bleu : Chine, Inde, Russie, Brésil, Mexique, Afrique du Sud, etc.) et les autres pays (gris, dont la plus grande partie de l'Afrique). \textbf{Le croquis attendu situe :} 1. les flux de marchandises, surtout entre les pôles de la Triade (flèches épaisses) ; 2. l'Afrique, à la marge (flèches fines) ; 3. les capitaux des FMN et les migrations des diasporas reliant Nord et Sud ; 4. les flux d'informations portés par les télécommunications, qui enveloppent l'ensemble. \\ Source : Wikimedia Commons, « Developed and Emerging markets », AlexCovarrubias, domaine public.}
\end{popfigure}

\begin{methode}[Construire un croquis de géographie au Bac]
\begin{enumerate}
\item Trace un \textbf{fond de carte} simple (contours simplifiés, orientation avec la flèche du nord, échelle indicative).
\item Choisis des \textbf{figurés simples et hiérarchisés} : points pour les villes, surfaces colorées pour les espaces, flèches pour les flux.
\item Rends les \textbf{flèches proportionnelles} à l'importance des flux (épaisseur variable).
\item Rédige une \textbf{légende organisée} : classe les figurés du plus important au moins important ; chaque figuré utilisé sur la carte doit figurer dans la légende.
\item Donne un \textbf{titre} précis qui annonce le sujet et l'espace étudié.
\end{enumerate}
\end{methode}

\begin{attention}[L'erreur classique des flèches]
Sur un croquis, une flèche doit être \emph{proportionnelle} à l'importance du flux qu'elle représente. L'erreur la plus fréquente au Bac est de dessiner des flèches toutes identiques entre l'Europe, l'Asie et l'Afrique : cela revient à affirmer, à tort, que ces flux ont la même importance. Une flèche Afrique--Europe doit être nettement plus fine qu'une flèche Asie--Amérique du Nord.
\end{attention}

\courssub{TP guidé : croquis et légendation d'une carte du commerce mondial}

\begin{experience}[TP : « Les facteurs et les flux de la mondialisation »]
\textbf{Matériel} : fond de planisphère simplifié, crayons de couleur, tableau statistique des parts du commerce mondial par grande région (Triade : environ 70 à 80 \% ; Afrique : environ 2 à 3 \% ; ordres de grandeur OMC).

\textbf{Protocole.}
\begin{enumerate}
\item Reporte sur le fond de carte les trois pôles de la Triade (surfaces colorées) et localise l'Afrique et le Cameroun.
\item Trace les principaux flux de marchandises avec des flèches \emph{proportionnelles} : transpacifique, transatlantique, Europe--Asie (épaisses) ; Afrique--Europe et Afrique--Asie (fines).
\item Ajoute en pointillés les flux de capitaux (IDE) et en pointillés fins les migrations des diasporas.
\item Organise la légende en trois rubriques : \emph{espaces}, \emph{flux}, \emph{facteurs} (ports, câbles, sièges de FMN).
\item Termine par le titre, l'orientation et l'échelle indicative.
\end{enumerate}

\textbf{Observations attendues} : les flèches les plus épaisses relient les pôles de la Triade ; l'Afrique n'est reliée au système mondial que par des flux faibles, malgré sa population nombreuse.

\textbf{Interprétation} : la carte confirme que la mondialisation est un processus \emph{inégal} et \emph{hiérarchisé}, dont les flux sont rendus possibles par les transports, les télécommunications, les FMN, les États et les réseaux sociaux (diasporas, ONG).
\end{experience}

\begin{exercice}[Pour t'entraîner]
1) Définis la balance commerciale et explique la différence entre excédent et déficit. 2) À l'aide du graphique de la leçon, montre en cinq lignes que le commerce mondial croît plus vite que la production, et propose deux explications. 3) Construis un croquis intitulé « La place du Cameroun dans le commerce mondial » en respectant les cinq étapes de la méthode.
\end{exercice}

\begin{corrige}[Exercice]
\textbf{1)} La balance commerciale est la différence entre la valeur des exportations et celle des importations d'un pays. Elle est excédentaire si les exportations dépassent les importations, déficitaire dans le cas contraire.
\textbf{2)} Entre 1950 et 2020, l'indice du commerce mondial passe d'environ 100 à près de \num{6 000}, contre environ \num{1 000} pour la production : le commerce croît environ deux fois plus vite. Explications : la baisse des coûts de transport (containerisation) et la libéralisation des échanges (GATT/OMC), auxquelles s'ajoutent les stratégies mondialisées des FMN.
\textbf{3)} Le croquis doit comporter : un fond de planisphère simplifié orienté ; le Cameroun localisé ; des flèches fines vers la Chine et l'UE (exportations de produits primaires) et des flèches d'importation (produits manufacturés) ; une légende organisée et un titre.
\end{corrige}

\begin{aretenir}[L'essentiel de la section]
Le commerce international, mesuré par les exportations, les importations et la balance commerciale, a explosé depuis 1950, plus vite que la production. Sa géographie est inégale : la Triade domine, la Chine monte, l'Afrique est marginalisée (environ 2 à 3 \% des échanges). Le Cameroun exporte des produits primaires (pétrole, cacao, bois, coton, bananes) et importe des produits manufacturés, avec la Chine et l'UE comme principaux partenaires. Les cinq facteurs de la mondialisation — transports et télécommunications, FMN, État et diaspora, associations et ONG, commerce international — forment un système dont les éléments se renforcent mutuellement.
\end{aretenir}

\newpage

\coursec{Activité d'intégration}

\coursec{Les facteurs de la mondialisation}

\courssub{Objectifs de l'activité}

À l'issue de ce travail pratique, tu seras capable de :

\begin{itemize}
\item \cle{observer}, \cle{localiser} et \cle{décrire} les principaux flux mondiaux (maritimes, numériques, financiers, migratoires, humanitaires) à partir d'un dossier documentaire ;
\item \cle{légender} une carte muette du monde en y plaçant les hubs portuaires, les routes maritimes, les câbles sous-marins, les foyers de diaspora et les zones d'intervention des ONG ;
\item \cle{construire un croquis} intitulé « Les facteurs de la mondialisation et la place du Cameroun », avec des figurés proportionnels et une légende organisée en cinq rubriques correspondant aux cinq savoirs du programme ;
\item \cle{rédiger un paragraphe d'interprétation} répondant à une question problématisée : le Cameroun est-il un acteur ou un spectateur de la mondialisation ?
\end{itemize}

Ce TP mobilise l'ensemble des savoir-faire du programme : observer et décrire, représenter et cartographier, inventorier et interpréter des documents.

\courssub{Situation}

En mars 2024, le club de géographie du Lycée bilingue de Yaoundé prépare une exposition intitulée « Le monde entre au Cameroun, le Cameroun entre dans le monde ». Le proviseur a confié au club une mission : produire une grande carte murale montrant comment les \cle{transports}, les \cle{télécommunications}, les \cle{firmes multinationales} (FMN), l'\cle{État}, la \cle{diaspora} et les \cle{associations et ONG} structurent les flux mondiaux, et quelle place le Cameroun y occupe. Le professeur de géographie remet aux élèves le dossier documentaire suivant.

\textbf{Document 1. Les grandes routes maritimes et les hubs portuaires mondiaux.} Environ 80 \% du commerce mondial en volume transite par la mer. Trois axes dominent : la route Asie--Europe par le canal de Suez (environ 12 \% du commerce mondial), la route transpacifique Asie--Amérique du Nord, et la route transatlantique. Les principaux hubs sont Shanghai (plus de 47 millions d'EVP par an, ordre de grandeur 2023), Singapour (environ 39 millions d'EVP), Rotterdam (environ 13 millions d'EVP), Dubaï et Tanger Med (environ 9 millions d'EVP, premier port africain). En Afrique centrale, Douala traite l'essentiel du commerce extérieur camerounais (environ 95 \% des importations du pays) ; le port en eau profonde de Kribi, entré en service en 2015 (première phase), ambitionne de devenir un hub régional.

\textbf{Document 2. Les câbles sous-marins de fibre optique.} Plus de 95 \% du trafic Internet intercontinental passe par des câbles sous-marins. Le Cameroun est connecté par plusieurs câbles dont SAT-3 (côte ouest africaine, depuis 2002), ACE, WACS et le câble SAIL reliant le Brésil à Kribi (mis en service 2018). Les débits ont fait exploser l'usage d'Internet : le taux de pénétration d'Internet au Cameroun est passé d'environ 4 \% en 2010 à plus de 40 \% au début des années 2020 (ordres de grandeur).

\textbf{Document 3. Quelques FMN présentes au Cameroun.}

\begin{center}
\begin{tabularx}{\linewidth}{|l|X|X|}\hline
\rowcolor{popblueL} \textbf{FMN} & \textbf{Secteur} & \textbf{Implantation au Cameroun} \\ \hline
Bolloré / AGL & Logistique, port, rail (Camrail) & Douala, Kribi, axe Douala--Ngaoundéré \\ \hline
Maersk / CMA CGM & Transport maritime & Douala, Kribi \\ \hline
MTN, Orange & Télécommunications & Ensemble du territoire \\ \hline
TotalEnergies & Hydrocarbures & Distribution nationale \\ \hline
SOSUCAM (groupe SOMDIAA) & Agro-industrie sucrière & Mbandjock, Nkoteng \\ \hline
\end{tabularx}
\end{center}

\textbf{Document 4. La diaspora et les transferts d'argent.} La diaspora camerounaise est estimée entre 3 et 5 millions de personnes (ordre de grandeur), installée surtout en France, aux États-Unis, au Canada, en Allemagne, au Nigeria et en Afrique du Sud. Selon la Banque mondiale, les transferts de fonds des migrants vers le Cameroun atteignent environ 300 à 400 millions de dollars par an (ordre de grandeur récent), soit un montant comparable à l'aide publique au développement (APD) reçue par le pays.

\textbf{Document 5. Le commerce extérieur du Cameroun et le rôle de l'État.} Le Cameroun exporte surtout des produits primaires : pétrole brut, cacao, bois, coton, café, bananes, aluminium. Ses principaux clients sont la Chine, les Pays-Bas, l'Inde, l'Italie et la France ; ses principaux fournisseurs sont la Chine, la France, l'Inde et la Belgique. La balance commerciale est structurellement déficitaire (déficit de l'ordre de 1 000 à 2 000 milliards de FCFA certaines années). L'État camerounais encadre cette insertion : membre de l'OMC, de la CEMAC et de la ZLECAf, il négocie des accords de partenariat économique (APE) avec l'UE et pilote le port de Kribi et le corridor ferroviaire.

\textbf{Document 6. Les associations et les ONG : acteurs de la régulation et de la solidarité.} De nombreuses ONG internationales (WWF, CARE International, Croix-Rouge, Médecins Sans Frontières) et associations locales (Réseau de lutte contre la faim, RELUFA ; Association pour la défense des droits de l'homme, ADH) opèrent au Cameroun. Elles interviennent dans l'humanitaire (régions de l'Extrême-Nord, Nord-Ouest, Sud-Ouest), l'environnement (forêt du bassin du Congo, parc de la Lobéké), la gouvernance (transparence extractive, ITIE) et le développement rural. Leurs financements proviennent de donateurs publics (États, UE, ONU) et privés (fondations, dons). Elles constituent un flux d'expertise, de capitaux solidaires et de normes (droits humains, normes environnementales) du Nord vers le Sud, et relaient les revendications locales vers les instances internationales.

\courssub{Consignes}

Par groupes de quatre, et à l'aide du dossier documentaire :

\begin{enumerate}
\item \textbf{Observer et décrire.} Pour chacun des six documents, relève : la nature du flux étudié (marchandises, données, capitaux, personnes, expertise/normes), sa direction principale, et le ou les facteurs de mondialisation qu'il illustre. Présente tes résultats dans un tableau à trois colonnes.
\item \textbf{Localiser.} Sur la carte muette du monde fournie, place et nomme : les hubs portuaires Shanghai, Singapour, Rotterdam, Tanger Med, Douala et Kribi ; les trois grandes routes maritimes ; les câbles sous-marins SAT-3/WACS/ACE (côte ouest Afrique) et SAIL (Brésil--Kribi) ; les foyers de la diaspora camerounaise (France, États-Unis, Canada, Allemagne, Nigeria) ; les sièges de quelques ONG internationales (ex. : Genève, Paris, Washington) et des zones de projets au Cameroun (Extrême-Nord, Sud-Est forestier).
\item \textbf{Calculer.} À partir du document 4, si les transferts de la diaspora s'élèvent à 350 millions de dollars et l'aide publique au développement (APD) à 700 millions de dollars, calcule la part des transferts de la diaspora dans l'ensemble des deux flux. Que conclus-tu sur le poids financier de la diaspora par rapport aux bailleurs institutionnels ?
\item \textbf{Construire un croquis.} Réalise un croquis intitulé « Les facteurs de la mondialisation et la place du Cameroun » : figurés proportionnels (flèches d'épaisseur variable pour les flux), figurés ponctuels pour les hubs et les sièges de FMN/ONG, figurés linéaires pour les câbles et routes, et légende organisée en \textbf{cinq rubriques} correspondant aux cinq savoirs de la leçon : 1. Transports, 2. Télécommunications, 3. Firmes multinationales, 4. État et diaspora, 5. Associations/ONG et commerce international.
\item \textbf{Interpréter.} Rédige un paragraphe de 10 à 15 lignes répondant à la question : \emph{le Cameroun est-il un acteur ou un spectateur de la mondialisation ?} Ton paragraphe doit mobiliser au moins quatre documents et nuancer la réponse.
\end{enumerate}

\courssub{Ressources à mobiliser}

\begin{aretenir}[Ce qu'il faut avoir en tête avant de commencer]
\begin{itemize}
\item Les \cle{facteurs de la mondialisation} : révolution des transports (conteneurisation, hubs portuaires), révolution des télécommunications (câbles sous-marins, Internet), stratégies des \cle{FMN} (délocalisations, réseaux mondiaux), rôle des \cle{États} (libéralisation, OMC, intégration régionale : CEMAC, ZLECAf), poids de la \cle{diaspora} (transferts de fonds) et des \cle{associations et ONG} (solidarité, normes, plaidoyer).
\item Savoir-faire cartographiques : choisir des \cle{figurés} adaptés (ponctuel, linéaire, proportionnel), hiérarchiser l'information, rédiger une \cle{légende} organisée en rubriques, donner un titre, une orientation et une échelle indicative.
\item Calculs utiles : part en pourcentage $= \dfrac{\text{partie}}{\text{total}} \times 100$ ; taux de variation $= \dfrac{V_f - V_i}{V_i} \times 100$.
\end{itemize}
\end{aretenir}

\courssub{Démarche et corrigé commenté}

\begin{exoresolu}[Corrigé commenté du TP]
\textbf{Étape 1 — Tableau d'observation (consigne 1).}

\begin{center}
\begin{tabularx}{\linewidth}{|l|X|X|}\hline
\rowcolor{popblueL} \textbf{Flux} & \textbf{Direction principale} & \textbf{Facteur illustré} \\ \hline
Marchandises (conteneurs) & Asie $\to$ Europe ; Afrique $\to$ Monde & Transports maritimes, conteneurisation \\ \hline
Données numériques & Intercontinental, par câbles sous-marins & Télécommunications \\ \hline
Capitaux, production, technologies (FMN) & Sièges (Nord) $\to$ Implantations mondiales & Firmes multinationales \\ \hline
Transferts d'argent (diaspora) & Europe, Am. du Nord, Afrique $\to$ Cameroun & Diaspora \\ \hline
Exportations primaires / Importations manufacturées & Cameroun $\leftrightarrow$ Chine, UE, Inde & Commerce international, rôle de l'État (OMC, CEMAC, ZLECAf) \\ \hline
Expertise, aide humanitaire, normes (ONG) & Nord (sièges) $\to$ Sud (terrains) ; Plaidoyer Sud $\to$ Nord & Associations et ONG \\ \hline
\end{tabularx}
\end{center}

\tcblower

\textbf{Étape 2 — Calcul de la part des transferts de la diaspora (consigne 3).}
\[
\text{Total des deux flux} = 350 + 700 = 1\,050 \text{ millions USD}
\]
\[
\text{Part diaspora} = \frac{350}{1\,050} \times 100 = \frac{1}{3} \times 100 \approx 33,3\,\%.
\]
Les transferts de la diaspora représentent environ un tiers de l'ensemble des deux flux financiers étudiés. La diaspora s'affirme comme un acteur financier majeur, dont le poids est comparable à celui des bailleurs institutionnels traditionnels, mais avec une logique de proximité et de ciblage familial différente.

\textbf{Étape 3 — Le croquis attendu (consignes 2 et 4).} Le fond de carte ci-dessous (trafic maritime mondial) et sa légende donnent le modèle de ce qui est attendu : planisphère, figurés proportionnels et légende organisée en cinq rubriques officielles.

\begin{popfigure}
\includegraphics[width=\linewidth,height=9.5cm,keepaspectratio]{N11/shipping_routes.jpg}
\legende{Fond de carte pour le croquis : le trafic maritime mondial (plus le bleu est foncé, plus le trafic est dense ; la carte ne comporte aucun nom). \textbf{Le croquis attendu situe et légende en 5 rubriques :} 1. \textbf{Transports} : hubs portuaires mondiaux (Shanghai, Singapour, Rotterdam, Tanger Med) et camerounais (Douala, Kribi), routes maritimes (traits bleus, épaisseur $\propto$ trafic) ; 2. \textbf{Télécommunications} : câbles sous-marins SAT-3/WACS/ACE et SAIL (tirets violets) ; 3. \textbf{Firmes multinationales} : sièges des FMN (carrés violets) ; 4. \textbf{État et diaspora} : transferts de la diaspora vers le Cameroun (flèches roses) ; 5. \textbf{Associations/ONG et commerce} : sièges et zones de projets des ONG (triangles verts), flux d'aide (flèches vertes pointillées). Sur ce fond, on lit déjà que les grands corridors relient l'Atlantique Nord, le Pacifique Nord et l'axe Europe--Suez--Asie, alors que la façade atlantique de l'Afrique est bien moins fréquentée. \\ Source : Wikimedia Commons, « Shipping routes » (densité du trafic maritime commercial), T. Hengl, CC BY-SA 3.0.}
\end{popfigure}

\textbf{Étape 4 — Paragraphe d'interprétation (consigne 5), proposition de rédaction.}

« Le Cameroun est à la fois acteur et spectateur de la mondialisation, mais davantage spectateur qu'acteur dans la division internationale du travail. Spectateur d'abord : il subit des flux organisés depuis l'extérieur — les grandes routes maritimes relient surtout l'Asie, l'Europe et l'Amérique du Nord (Doc. 1), et le port de Douala, qui concentre environ 95 \% des importations, reste un port secondaire à l'échelle mondiale ; il exporte des produits primaires peu transformés (pétrole, cacao, bois) et accuse un déficit commercial chronique (Doc. 5). Les FMN étrangères (Bolloré, Maersk, MTN, TotalEnergies) contrôlent des secteurs clés de son économie (Doc. 3). Acteur cependant : l'État pilote des infrastructures stratégiques (port de Kribi, câble SAIL) et s'engage dans l'intégration régionale (ZLECAf, CEMAC) pour peser sur la libéralisation (Doc. 5, 2) ; la diaspora injecte environ 350 millions de dollars par an — un tiers des flux financiers extérieurs étudiés — devenant un levier de développement endogène (Doc. 4) ; les ONG internationales et locales apportent expertise, normes environnementales et aide humanitaire, reliant les territoires périphériques aux instances globales (Doc. 6). Le Cameroun est donc un acteur émergent qui cherche, par des choix souverains (Kribi, ZLECAf, mobilisation de la diaspora), à transformer une insertion subie en insertion choisie. »
\end{exoresolu}

\begin{attention}[Erreurs fréquentes à éviter dans ce TP]
\begin{itemize}
\item Confondre \cle{flux} (mouvement : marchandises, capitaux, personnes, données, expertise) et \cle{facteur} (ce qui rend le flux possible : transports, FMN, États, diaspora, ONG).
\item Dessiner des flèches d'épaisseur identique pour des flux très inégaux : le figuré proportionnel est justement fait pour hiérarchiser visuellement les intensités.
\item Oublier le titre, l'orientation (flèche Nord), l'échelle indicative ou une rubrique de la légende : au Baccalauréat, un croquis sans légende organisée et complète perd une grande partie des points.
\item Répondre « acteur » ou « spectateur » de façon catégorique : le correcteur attend une réponse \textbf{nuancée}, appuyée sur des chiffres du dossier (parts de marché, montants transferts, déficit commercial, nombre de câbles).
\item Négliger le facteur « Associations et ONG » : c'est un savoir du programme à part entière, il doit apparaître dans le tableau d'observation, sur la carte (sièges/projets) et dans la légende (rubrique 5).
\end{itemize}
\end{attention}

\courssub{Bilan de l'activité}

Ce TP a permis de mettre en évidence trois idées essentielles. D'abord, la mondialisation n'est pas un phénomène abstrait : elle repose sur des \cle{supports matériels} concrets — navires porte-conteneurs, hubs portuaires, câbles sous-marins — que l'on peut localiser et cartographier. Ensuite, les \cle{facteurs} de la mondialisation agissent en système : les FMN organisent les flux que les transports et les télécommunications rendent possibles, tandis que les États, les diasporas et les ONG encadrent, financent, régulent ou contestent ces dynamiques. Enfin, la cartographie révèle une \cle{hiérarchie des territoires} : quelques hubs (Shanghai, Singapour, Rotterdam) concentrent l'essentiel des flux, tandis que des pays comme le Cameroun occupent une position périphérique mais en évolution, entre insertion subie (exportations primaires, déficit commercial) et stratégies d'insertion active (Kribi, ZLECAf, câble SAIL, mobilisation de la diaspora, plaidoyer ONG). Tu disposes désormais d'une méthode complète — observer, localiser, calculer, croquis, interpréter — directement transposable à l'épreuve de géographie du Baccalauréat.

\newpage

\coursec{Méthodes et exercices résolus}

\courssub{Méthodes et exercices résolus}

\begin{methode}[Analyser un tableau statistique : flux commerciaux et IDE]
\textbf{Objectif :} Inventorier, classer et interpréter des données chiffrées (exportations, importations, IDE, transferts de la diaspora) pour mesurer l'intégration du Cameroun dans la mondialisation.
\begin{enumerate}
    \item \textbf{Lire la fiche technique :} Titre, source (INS, BEAC, FMI, Douanes, CNUCED), année, unité (milliards FCFA, \%, part de marché).
    \item \textbf{Identifier les variables :} Lignes (pays partenaires, produits, secteurs) et colonnes (années, flux : export/import, solde, IDE entrants/sortants).
    \item \textbf{Calculer les indicateurs clés :}
    \begin{itemize}
        \item Part d'un partenaire : $\frac{\text{Flux avec partenaire}}{\text{Flux total}} \times 100$.
        \item Taux de variation annuel : $\frac{V_{n} - V_{n-1}}{V_{n-1}} \times 100$.
        \item Taux d'ouverture commerciale : $\frac{\text{Exportations} + \text{Importations}}{\text{PIB}} \times 100$.
        \item Couverture des importations par les exportations : $\frac{\text{Exportations}}{\text{Importations}} \times 100$.
    \end{itemize}
    \item \textbf{Classer et hiérarchiser :} Top 5 partenaires, produits dominants, excédent/déficit par zone (CEMAC, UE, Chine, Inde).
    \item \textbf{Interpréter géographiquement :} Lier les chiffres aux facteurs (FMN extractives, accords APE/UE, corridors, diaspora). Rédiger une synthèse structurée : constats majeurs $\rightarrow$ explications (facteurs) $\rightarrow$ limites/enjeux.
\end{enumerate}
\cle{Réflexe} : Toujours vérifier la cohérence des totaux (somme des parts $\approx 100\%$) et citer la source et l'année dans la conclusion.
\end{methode}

\begin{exoresolu}[Analyse de la structure des échanges du Cameroun (Données 2022-2023 approchées)]
\textbf{Document :} Tableau simplifié des échanges commerciaux du Cameroun (en milliards FCFA, sources : INS, Douanes, BEAC).
\begin{center}
\begin{tabularx}{\linewidth}{|l|X|X|X|X|}\hline
\rowcolor{popblueL}
\textbf{Partenaire / Zone} & \textbf{Exportations 2022} & \textbf{Importations 2022} & \textbf{Exportations 2023} & \textbf{Importations 2023} \\ \hline
Union Européenne (UE) & 1\,200 & 1\,500 & 1\,250 & 1\,550 \\ \hline
Chine & 400 & 1\,100 & 420 & 1\,150 \\ \hline
CEMAC (zone franc) & 300 & 150 & 320 & 160 \\ \hline
Inde & 250 & 300 & 280 & 320 \\ \hline
Autres (USA, Turquie, etc.) & 600 & 800 & 630 & 850 \\ \hline
\textbf{Total} & \textbf{2\,750} & \textbf{3\,850} & \textbf{2\,900} & \textbf{4\,030} \\ \hline
\end{tabularx}
\end{center}
\textbf{PIB 2022 : 27\,000 Mds FCFA ; PIB 2023 : 28\,500 Mds FCFA.}
\textbf{Questions :}
\begin{enumerate}
    \item Calculer le taux d'ouverture commerciale pour 2022 et 2023. Interpréter l'évolution.
    \item Calculer le taux de couverture des importations par les exportations pour 2023.
    \item Déterminer le premier client et le premier fournisseur en 2023. Calculer leurs parts de marché respectives.
    \item En vous appuyant sur les résultats et vos connaissances, expliquez le rôle des \textbf{firmes multinationales (FMN)} et des \textbf{transports} dans cette structure.
\end{enumerate}
\tcblower
\textbf{Solution détaillée}
\begin{enumerate}
    \item \textbf{Taux d'ouverture :}
    \begin{align*}
    2022 &: \frac{2\,750 + 3\,850}{27\,000} \times 100 = \frac{6\,600}{27\,000} \times 100 \approx \textbf{24,4\%} \\
    2023 &: \frac{2\,900 + 4\,030}{28\,500} \times 100 = \frac{6\,930}{28\,500} \times 100 \approx \textbf{24,3\%}
    \end{align*}
    \textit{Interprétation :} Le taux est stable autour de 24\%, niveau modéré (moyenne SSA $\approx$ 50-60\%). L'économie reste peu ouverte, freinée par l'informel, les barrières non tarifaires et la faible transformation locale. La légère baisse s'explique par une croissance du PIB (gaz, BTP) plus rapide que celle du commerce.
    \item \textbf{Taux de couverture 2023 :} $\frac{2\,900}{4\,030} \times 100 \approx \textbf{71,9\%}$.
    \textit{Analyse :} Déficit structurel de la balance commerciale (importations > exportations). Le pays importe des biens d'équipement, produits pétroliers raffinés, médicaments, céréales (riz, blé) que les FMN (MSC/TiL ex-Bolloré, TotalEnergies, CFAO, sociétés chinoises) fournissent via le port de Douala/Kribi.
    \item \textbf{Premier client 2023 :} UE (1\,250 Mds). Part : $\frac{1\,250}{2\,900} \times 100 \approx \textbf{43,1\%}$.
    \textbf{Premier fournisseur 2023 :} UE (1\,550 Mds). Part : $\frac{1\,550}{4\,030} \times 100 \approx \textbf{38,5\%}$.
    \textit{Note :} La Chine est le 2ème fournisseur (28,5\%) mais seulement 4ème client (14,5\%).
    \item \textbf{Rôle des facteurs :}
    \begin{itemize}
        \item \textbf{FMN :} Dominent l'export (pétrole : Perenco, SNH/TotalEnergies ; bois : Rougier, Pallisco ; agro : SOCAPALM, CDC, Plantations du Haut Penja - PHP/Compagnie Fruitière ; aluminium : Alucam/Rio Tinto). Elles orientent les flux vers l'UE (accords APE) et la Chine (bois, coton). Elles importent intrants et machines.
        \item \textbf{Transports :} Le corridor Douala-Ndjamena/Bangui (route/rail) géré par Camrail (MSC/TiL) et le port de Douala, et le nouveau port en eau profonde de Kribi (Kribi Conteneur Terminal - KCT, exploité par CHEC/China Harbour) structurent l'insertion. Le temps de dédouanement (objectif 48h) et le coût du fret conditionnent la compétitivité.
    \end{itemize}
    \textbf{Conclusion :} Le Cameroun reste une économie d'enclave, extravertie, dépendante des FMN et des infrastructures portuaires/corridors pour s'insérer dans la mondialisation, avec un partenaire historique (UE) hégémonique mais un fournisseur émergent (Chine) incontournable.
\end{enumerate}
\end{exoresolu}

\begin{methode}[Construire un croquis de synthèse : Les flux de la mondialisation au Cameroun]
\textbf{Objectif :} Représenter, schématiser et cartographier les facteurs (ports, corridors, FMN, diaspora, ONG) sur une base cartographique simplifiée du Cameroun.
\begin{enumerate}
    \item \textbf{Choisir la base :} Contour simplifié du Cameroun (forme « triangle » pointe au Nord Extrême, base au Sud), méridiens/parallèles principaux (10°E, 15°E ; 2°N, 13°N). Échelle indicative (ex: 1 cm $\approx$ 100 km). Flèche Nord.
    \item \textbf{Structurer la légende AVANT de tracer} (3 à 4 items max par catégorie) :
    \begin{itemize}
        \item \textbf{Portes d'entrée/sortie :} Ports (Douala, Kribi, Limbé, Garoua fluvial), Aéroports (Douala, Yaoundé-Nsimalen, Garoua, Maroua).
        \item \textbf{Corridors / Axes de transport :} Route (N1, N2, N3, N10), Rail (Camrail Douala-Ngaoundéré), Oléoduc (Tchad-Cameroun), Fibre optique (backbone national, câbles sous-marins SAT-3, WACS, ACE, 2Africa).
        \item \textbf{Implantations clés des FMN :} Pétrole (Kribi, Rio del Rey, Logbaba), Bois (Est, Sud), Agro (Sud-Ouest, Littoral, Moungo), Mines (Bipindi-Grand Zambi, Mbalam - projet), Télécoms (Yaoundé, Douala, Bafoussam - Data centers).
        \item \textbf{Flux immatériels :} Flèches « Transferts diaspora » (vers Yaoundé/Douala/Bafoussam), Sièges ONG internationales (Yaoundé, Bertoua, Maroua).
    \end{itemize}
    \item \textbf{Sémiologie graphique rigoureuse :}
    \begin{itemize}
        \item \textbf{Points :} Villes (ronds proportionnels population/fonction), FMN (carrés/losanges couleur secteur : vert=agro, noir=pétrole, bleu=bois, orange=mines, violet=télécoms).
        \item \textbf{Lignes :} Flux (flèches épaisseur $\propto$ volume), Routes (trait continu), Rail (trait pointillé), Oléoduc (trait double), Fibre (trait fin zigzag).
        \item \textbf{Surfaces :} Zones d'influence (hachures), ZAE (carrés grisés).
    \end{itemize}
    \item \textbf{Réaliser le croquis :} Crayon gris $\rightarrow$ Encre/Couleur. Titre complet : « Le Cameroun dans la mondialisation : portes, corridors et acteurs (2024) ». Légende encadrée, organisée, lisible.
\end{enumerate}
\cle{Réflexe} : Un croquis sans légende structurée = 0 point. L'épaisseur des flèches doit être proportionnelle (qualitativement) aux volumes réels.
\end{methode}

\begin{exoresolu}[Construction guidée d'un croquis : « Le Cameroun, carrefour d'Afrique centrale »]
\textbf{Consigne :} À partir de la description ci-dessous, réaliser le croquis demandé en appliquant la méthode. On vous fournit le code TikZ de la base et de la solution pour vérification.
\textbf{Éléments à représenter :}
\begin{itemize}
    \item Villes : Yaoundé (capitale), Douala (pôle économique, port), Kribi (port eau profonde), Ngaoundéré (fin rail), Garoua (port fluvial, aéroport), Maroua (aéroport, ONG humanitaires), Bafoussam (diaspora, agro).
    \item Axes : Route Douala-Yaoundé-Bertoua-Ngaoundéré (N1/N15), Route Douala-Bafoussam-Bamenda (N6), Rail Douala-Ngaoundéré (Camrail), Oléoduc Kribi-Doba (Tchad).
    \item FMN : Alucam (Édéa), SNH/Perenco (Kribi/Logbaba), PHP (Tiko/Moungo), SOCAPALM (Kienké, Dibombari, Mbongo), Camtel/MTN/Orange (Yaoundé, Douala - Data Centers), CHEC (Kribi Port).
    \item Flux : Export bois/minérais vers Kribi/Douala (flèches vertes/noires), Import conteneurs Douala/Kribi $\rightarrow$ Hinterland (Tchad, RCA) (flèches bleues épaisses), Transferts diaspora (flèches pointillées violettes vers Yaoundé/Douala/Bafoussam).
\end{itemize}
\tcblower
\textbf{Solution détaillée (Code TikZ commenté + Analyse)}

\textit{1. Analyse de la structure spatiale :}
Le croquis met en évidence la \textbf{dichotomie côtière/intérieur}. Douala/Kribi sont les « portes » maritimes. Yaoundé est le nœud de commande (État, ONG, Télécoms). Les corridors (Route/Rail/Oléoduc) structurent l'orientation Nord-Sud vers le Tchad/RCA. La diaspora (USA, France, Allemagne) alimente Yaoundé, Douala, Bafoussam (Bamiléké).

\textit{2. Code TikZ de la solution (à compiler pour visualiser) :}
\begin{verbatim}
\begin{tikzpicture}[scale=0.85]
% Base contour simplifié Cameroun
\draw[popdark, thick] (0,0) -- (10,0) -- (5,12) -- cycle; % Triangle approx
% Villes
\node[circle,fill=popblue,minimum size=4pt,inner sep=0] (Yde) at (4.5,7) {};
\node[circle,fill=poporange,minimum size=5pt,inner sep=0] (Dla) at (2.5,3) {};
\node[circle,fill=poporange,minimum size=4pt,inner sep=0] (Kribi) at (2,1.5) {};
\node[circle,fill=popgreen,minimum size=3pt,inner sep=0] (Nga) at (6,10) {};
\node[circle,fill=popgreen,minimum size=3pt,inner sep=0] (Gar) at (7.5,9) {};
\node[circle,fill=popgreen,minimum size=3pt,inner sep=0] (Mar) at (8.5,5) {};
\node[circle,fill=poppurple,minimum size=4pt,inner sep=0] (Baf) at (3.5,5.5) {};
% Labels
\node[above right] at (Yde) {\small Yaoundé};
\node[below left] at (Dla) {\small Douala};
\node[below] at (Kribi) {\small Kribi};
\node[above right] at (Nga) {\small Ngaoundéré};
\node[above] at (Gar) {\small Garoua};
\node[right] at (Mar) {\small Maroua};
\node[left] at (Baf) {\small Bafoussam};
% Axes
\draw[popline, thick] (Dla) -- (Yde) -- (Nga); % Route/Rail N1/N15
\draw[popline, thick, dashed] (Dla) -- (Baf) -- (4,8); % N6 vers Bamenda (approx)
\draw[poporange, thick, double distance=1pt] (Kribi) -- (8,-1); % Oleoduc vers Tchad (hors carte)
\draw[popblue, thin, densely dotted] (Dla) -- (Yde) -- (Nga); % Fibre optique backbone
% FMN (Carrés)
\draw[fill=popgreen!50!black] (3.2,4.5) rectangle ++(0.3,0.3) node[above right,font=\tiny] {PHP};
\draw[fill=black!50] (1.8,1.8) rectangle ++(0.3,0.3) node[above,font=\tiny] {KCT/CHEC};
\draw[fill=poporange] (3.5,4) rectangle ++(0.3,0.3) node[above,font=\tiny] {Alucam};
\draw[fill=poppurple] (Yde) ++(0.3,0.3) rectangle ++(0.3,0.3) node[above right,font=\tiny] {Data Centers};
% Flux Export (Vert)
\draw[popgreen, thick, ->] (4,8) -- (Dla); % Bois Est -> Douala
\draw[popgreen, thick, ->] (4,8) -- (Kribi); % Mines/Bois -> Kribi
% Flux Import (Bleu)
\draw[popblue, very thick, ->] (Dla) -- (Nga); % Conteneurs -> Hinterland
\draw[popblue, very thick, ->] (Kribi) -- (Nga);
% Diaspora (Violet pointillé)
\draw[poppurple, thick, dashed, ->] ($(Yde)+(1,1)$) -- (Yde) node[above,font=\tiny] {Diaspora};
\draw[poppurple, thick, dashed, ->] ($(Dla)+(1,-1)$) -- (Dla);
\draw[poppurple, thick, dashed, ->] ($(Baf)+(-1,1)$) -- (Baf);
% Légende (simplifiée ici, à faire à part en box)
\end{tikzpicture}
\end{verbatim}

\textit{3. Légende type à présenter dans la copie (format tableau) :}
\begin{center}
\begin{tabular}{|l|l|}\hline
\rowcolor{popblueL} \textbf{SYMBOLE} & \textbf{SIGNIFICATION} \\ \hline
\raisebox{-0.5pt}{\tikz{\draw[poporange, thick, double distance=1pt] (0,0)--(1,0);}} & Oléoduc Tchad-Cameroun (Kribi-Doba) \\ \hline
\raisebox{-0.5pt}{\tikz{\draw[popline, thick] (0,0)--(1,0);}} & Corridor routier/ferroviaire principal (Camrail) \\ \hline
\raisebox{-0.5pt}{\tikz{\draw[popblue, very thick, ->] (0,0)--(1,0);}} & Flux imports conteneurisés (Hinterland Tchad/RCA) \\ \hline
\raisebox{-0.5pt}{\tikz{\draw[popgreen, thick, ->] (0,0)--(1,0);}} & Flux exports matières premières (Bois, Mines, Agro) \\ \hline
\raisebox{-0.5pt}{\tikz{\draw[poppurple, thick, dashed, ->] (0,0)--(1,0);}} & Transferts financiers de la diaspora \\ \hline
\tikz{\node[circle,fill=popblue,minimum size=4pt,inner sep=0] at (0.5,0) {};} & Capitale politique/administrative (Yaoundé) \\ \hline
\tikz{\node[circle,fill=poporange,minimum size=5pt,inner sep=0] at (0.5,0) {};} & Métropole économique / Port principal (Douala) \\ \hline
\tikz{\draw[fill=popgreen!50!black] (0,0) rectangle (0.3,0.3);} & FMN Agro-industrie (PHP, CDC, SOCAPALM) \\ \hline
\tikz{\draw[fill=black!50] (0,0) rectangle (0.3,0.3);} & FMN Pétrole/Minier/Portuaire (SNH, Perenco, CHEC) \\ \hline
\tikz{\draw[fill=poppurple] (0,0) rectangle (0.3,0.3);} & FMN Télécoms / Data Centers (MTN, Orange, Camtel) \\ \hline
\end{tabular}
\end{center}
\textbf{Conclusion du croquis :} Le Cameroun affirme son rôle de \cle{porte océane} pour l'Afrique centrale (Tchad, RCA) via Douala/Kribi. L'État (via la SNH, Camrail, Camtel) et les FMN (MSC/TiL, CHEC, PHP, MTN) sont les opérateurs clés. La diaspora et les ONG (Maroua/Bertoua) complètent l'insertion par des flux financiers et humanitaires.
\end{exoresolu}

\begin{methode}[Commenter une photographie géographique : Infrastructure portuaire ou zone industrielle]
\textbf{Objectif :} Observer, localiser et décrire un paysage photographique (port de Kribi, zone industrielle de Douala-Bassa, antenne relais MTN, chantier route) pour en dégager les enjeux de mondialisation.
\begin{enumerate}
    \item \textbf{Identification (Légende) :} Nature du document (photo aérienne/au sol), lieu précis (Kribi Deep Sea Port, zone portuaire de Douala, usine Alucam Édéa), date approximative, photographe/source.
    \item \textbf{Description structurée (Plan par plans) :}
    \begin{itemize}
        \item \textbf{Plan général / Localisation :} Site (côte, embouchure, plateau), relief, hydrographie.
        \item \textbf{Plan moyen / Infrastructures :} Bassins, quais, grues (portiques), terminaux conteneurs/vrac, routes d'accès, rail, stockages (cuves pétrolières, hangars à bois, silos).
        \item \textbf{Plan détaillé / Activité :} Navires (porte-conteneurs, vraquiers, méthaniers), camions, engins de manutention, personnel, signalétique (logos FMN : CHEC, MSC/TiL, Maersk, MSC, SNH).
    \end{itemize}
    \item \textbf{Analyse / Interprétation (Le « Pourquoi » géographique) :}
    \begin{itemize}
        \item \textbf{Fonction :} Porte d'entrée/sortie, interface mer-terre, nœud logistique.
        \item \textbf{Acteurs (Facteurs) :} État (maître d'ouvrage, SNH, Port Autonome), FMN (concessionnaires : KCT/CHEC, Douala Terminal/MSC/TiL, SNH/Perenco), Transporteurs (Camrail, routiers).
        \item \textbf{Flux :} Import (riz, blé, clinker, hydrocarbures, machines), Export (bois, coton, cacao, aluminium, pétrole, GNL).
        \item \textbf{Aménagement / Enjeux :} Profondeur d'eau (16-18m à Kribi vs 7-10m à Douala), extension polder, ZES (Zone Économique Spéciale de Kribi), désenclavement (autoroute Kribi-Lolodorf, rail Mbalam-Kribi projeté).
    \end{itemize}
    \item \textbf{Conclusion synthétique :} Lien avec le cours (libéralisation, attractivité, corridors, concurrence régionale).
\end{enumerate}
\cle{Réflexe} : Ne pas faire un inventaire « catalogue ». Grouper les observations par \textbf{thèmes géographiques} (relief, réseau, acteurs, flux). Utiliser le vocabulaire précis : \emph{chenal d'accès}, \emph{poste à quai}, \emph{arrière-port}, \emph{concession}, \emph{hinterland}.
\end{methode}

\begin{exoresolu}[Commentaire de photographie : Le Port Autonome de Kribi (PAK) - Phase 1]
\textbf{Document :} \textit{Photographie aérienne (drone) du Port de Kribi, prise en 2023. On aperçoit : au premier plan, l'océan Atlantique ; au second plan, deux bassins rectangulaires découpés dans la mangrove/forêt, bordés de quais en béton ; trois portiques de quai (STS cranes) rouges/blancs (logo CHEC/KCT) ; un navire porte-conteneurs (logo Maersk) à quai ; des alignements de conteneurs colorés sur le terre-plein ; une route à 2x2 voies (accès autoroutier) et une voie ferrée (en construction/terrassement) longeant le port ; à l'arrière-plan, la forêt dense et le village de Kribi.}
\textbf{Consigne :} Présentez ce document et montrez en quoi il illustre la place du Cameroun dans la mondialisation.
\tcblower
\textbf{Solution détaillée}
\textbf{1. Présentation du document}
Il s'agit d'une \textbf{photographie aérienne verticale} (prise par drone) du \textbf{Port Autonome de Kribi (PAK)}, infrastructure phare de la \cle{stratégie d'émergence du Cameroun (SND30)}, située dans la région du Sud, à environ 150 km au sud de Douala. Le cliché date de 2023 (phase 1 opérationnelle). Source : Port Autonome de Kribi / CHEC.

\textbf{2. Description analytique}
\begin{itemize}
    \item \textbf{Milieu naturel :} Site littoral bas, sableux, à l'embouchure de la Kienké/Lobé. Forêt tropicale humide en arrière-plan (contrainte d'espace, enjeu environnemental mangroves). Profondeur naturelle exceptionnelle (> 16 m) permettant l'accueil de navires \emph{Post-Panamax} / \emph{New Panamax} (capacité > 10 000 EVP), impossible à Douala (chenal 7,5 m).
    \item \textbf{Infrastructures portuaires (Plan moyen) :} Deux bassins (bassins 1 et 2) protégés par des brise-lames. Quais en béton armé (longueur cumulée $\approx$ 600 m phase 1). Trois portiques de quai \textbf{STS (Ship-to-Shore)} de marque chinoise (ZPMC/CHEC), hauteur 45 m, portée 22 rangées. Zone de stockage (arrière-port) : terre-pleins conteneurs (capacité $\approx$ 30 000 EVP), zone vrac. Accès routier : autoroute Kribi-Lolodorf (2x2 voies, financée par Exim Bank Chine) visible. Emprise ferroviaire (terrassement) pour futur chemin de fer Mbalam-Kribi.
    \item \textbf{Activité (Plan détaillé) :} Un porte-conteneurs \textbf{Maersk} (FMN danoise, 1ère ligne mondiale) à quai, opération de chargement/déchargement active (spreaders sur conteneurs). Camions porte-conteneurs (flotte locale/MSC/TiL) en manœuvre. Absence visible de grues mobiles sur pneus (RTG/RMG automatisés attendus phase 2).
\end{itemize}

\textbf{3. Interprétation : Les facteurs de la mondialisation en action}
\begin{itemize}
    \item \textbf{Transport maritime / Infrastructure :} Kribi incarne la \textbf{modernisation des transports} (facteur 1). C'est une réponse à la saturation de Douala (embouteillages, tirant d'eau limité). Le choix d'un port \emph{« Greenfield »} (site vierge) permet une logistique fluide (autoroute dédiée, rail prévu).
    \item \textbf{FMN / Investissement étranger direct (IED) :} La concession du terminal à conteneurs (KCT) à \textbf{CHEC (China Harbour Engineering Company)} pour 25 ans illustre le rôle des \textbf{FMN chinoises} dans les infrastructures africaines (BRI - Belt and Road Initiative). La présence du navire \textbf{Maersk} montre l'intégration aux \textbf{réseaux maritimes mondiaux} (lignes Europe-Afrique, Asie-Afrique).
    \item \textbf{État / Stratégie :} L'État camerounais (Port Autonome, SNH pour le terminal GNL/Gaz à côté) est le \textbf{régulateur et co-investisseur} (part public-privé). Objectif : capter le trafic de transit Tchad/RCA (hinterland) et devenir \emph{Hub} sous-régional (concurrence Lomé, Libreville, Pointe-Noire).
    \item \textbf{Commerce international :} Les conteneurs (import : biens de consommation, intrants industriels ; export : bois sciés, coton, aluminium, futur fer Mbalam) matérialisent la \textbf{libéralisation des échanges}.
\end{itemize}

\textbf{4. Conclusion}
Cette photo résume la \textbf{réorientation stratégique} du Cameroun : d'un port unique (Douala) contraint, vers un système portuaire dual (Douala-Kribi) attractif pour les FMN logistiques (MSC/TiL, CHEC, Maersk, CMA CGM) et miniers (Sundance/Australia, Sinosteel/Chine pour Mbalam). Le défi reste le \textbf{désenclavement ferroviaire} (Mbalam-Kribi, Ngaoundéré-Kribi) et la \textbf{compétitivité des coûts} (facture portuaire, temps d'attente) face aux ports voisins.
\end{exoresolu}

\begin{methode}[Étudier un cas d'acteur non étatique : FMN, Diaspora ou ONG]
\textbf{Objectif :} Inventorier, classer et interpréter des documents (rapports annuels, articles de presse, données CNUCED, enquêtes) sur un acteur spécifique pour expliquer son rôle dans la mondialisation camerounaise.
\begin{enumerate}
    \item \textbf{Cadrer l'acteur :} Nom, nationalité, secteur (pétrole, bois, agro, télécoms, banque, logistique, humanitaire), taille (CA, effectifs, filiales).
    \item \textbf{Analyser la stratégie spatiale :} Implantations au Cameroun (usines, plantations, bureaux, antennes), liens avec la maison-mère, flux intra-firme (transferts de technologie, rapatriement bénéfices, import intrants).
    \item \textbf{Identifier les interactions avec l'État et le territoire :} Convention d'établissement, fiscalité (régime d'incitation, convention de stabilité), emploi (nationaux/expatriés), RSE (routes, écoles, santé), contentieux (droit du travail, environnement - ex: PHP/Greenpeace, Herakles Farms).
    \item \textbf{Évaluer l'apport à la mondialisation :} Part dans exportations/PIB/IDE, insertion dans les chaînes de valeur mondiales (CVM), transfert de compétences, dépendance (prix mondiaux, décisions extérieures).
    \item \textbf{Pour la Diaspora/ONG :} Cartographier les flux financiers (transferts formels/informels - Western Union, Mobile Money, Wave, virements), zones d'origine/destination, usage (consommation, investissement foncier/immobilier, PME). Pour ONG : zones d'intervention (Extrême-Nord, Est, Nord-Ouest/Sud-Ouest), secteurs (santé, éducation, protection, sécurité alimentaire), financement (donateurs publics/privés internationaux).
\end{enumerate}
\cle{Réflexe} : Distinguer \textbf{IDE (Investissement Direct Étranger)} - contrôle durable (>10\% capital) - et \textbf{Investissement de Portefeuille}. Citer des noms propres : \textbf{MSC/TiL (Logistique, ex-Bolloré), PHP/Compagnie Fruitière (Banane), SOCAPALM (Palmier), Alucam/Rio Tinto (Aluminium), Perenco/SNH (Pétrole), MTN/Orange/Camtel (Télécoms), CFAO/Toyota (Distribution), UBA/Ecobank/Afriland (Banque), OMS/UNICEF/MSF/CRS/Caritas (ONG)}.
\end{methode}

\begin{exoresolu}[Étude de cas : La diaspora camerounaise, premier investisseur du pays ?]
\textbf{Documents :}
\begin{itemize}
    \item \textbf{Doc 1 (Banque Mondiale / BEAC) :} Transferts formels de la diaspora vers le Cameroun : $\approx$ 300 Mds FCFA/an (moyenne 2020-2023). Estimation informel : $\times$ 2 à 3.
    \item \textbf{Doc 2 (Enquête INS/Diaspora 2021) :} 60\% des transferts vers le Littoral et le Centre (Yaoundé/Douala). 15\% vers l'Ouest (Bafoussam/Dschang). Usage : 55\% consommation courante, 25\% immobilier/foncier, 10\% éducation/santé, 5\% investissement productif (PME, agro).
    \item \textbf{Doc 3 (Ministère Relations Extérieures) :} Création du \emph{Programme d'Appui aux Investissements de la Diaspora (PAID)} et de la \emph{Caravane de la Diaspora}. Cible : 1 000 projets/an.
\end{itemize}
\textbf{Questions :}
\begin{enumerate}
    \item En utilisant les documents, montrez que la diaspora est un acteur majeur de la mondialisation camerounaise.
    \item Expliquez pourquoi l'État camerounais cherche à « capter » et « orienter » ces flux (Doc 3).
    \item Quels sont les freins à la transformation de ces transferts en investissements productifs structurants ?
\end{enumerate}
\tcblower
\textbf{Solution détaillée}
\begin{enumerate}
    \item \textbf{Acteur majeur : Preuves chiffrées.}
    \begin{itemize}
        \item \textbf{Volume financier :} $\approx$ 300 Mds FCFA/an (formel) $\rightarrow$ potentiel 600-900 Mds FCFA (avec informel). \textbf{Comparaison cruciale :} IDE entrants annuels (flux) au Cameroun $\approx$ 200-300 Mds FCFA (CNUCED). Aide Publique au Développement (APD) $\approx$ 200-250 Mds FCFA. \textbf{La diaspora est le 1er pourvoyeur de devises net du pays}, devant les exportations de cacao/bois/coton et l'IDE.
        \item \textbf{Flux transnationaux :} Réseaux États-Unis, France, Allemagne, Belgique, Canada, Pays du Golfe $\rightarrow$ Cameroun. Utilisation des opérateurs mondiaux (Western Union, MoneyGram) et fintechs (Wave, Orange Money, MTN MoMo, Wise) $\rightarrow$ \textbf{Télécommunications} (facteur clé).
        \item \textbf{Impact territorial :} Concentration Littoral/Centre/Ouest (Doc 2). Dynamise l'immobilier (Yaoundé/Douala/Bafoussam), le commerce de détail, le BTP local.
    \end{itemize}
    \item \textbf{Volonté étatique (Doc 3) :}
    \begin{itemize}
        \item \textbf{Perte de change :} Une part importante reste en numéraire ou finance de l'importation (consommation) $\rightarrow$ fuite de devises.
        \item \textbf{Levier de développement :} L'État veut transformer la \emph{rente migratoire} en \emph{capital productif} (création d'emplois, industrialisation, agro-business). Le PAID vise à baisser le risque (garanties, accompagnement technique, fiscale).
        \item \textbf{Souveraineté / Diplomatie :} Reconnaitre la diaspora comme \emph{« 11ème région »} renforce le \emph{soft power} et le lobbying international.
    \end{itemize}
    \item \textbf{Freins structurels :}
    \begin{itemize}
        \item \textbf{Climat des affaires :} Rank Doing Business (faible), corruption, insécurité juridique (conflits fonciers), justice lente.
        \item \textbf{Confiance :} Méfiance vis-à-vis de l'administration, gestion à distance difficile, arnaques familiales.
        \item \textbf{Structure des transferts :} Majorité « \emph{sociale} » (survie familles) et non « \emph{investissement} » (création valeur ajoutée). Manque d'instruments financiers adaptés (obligations diaspora, fonds d'investissement dédiés, capital-risque).
        \item \textbf{Double imposition / Fiscalité :} Complexité, absence d'incitations fortes (exonération IR sur investissement productif ?).
    \end{itemize}
    \textbf{Synthèse :} La diaspora est le \textbf{moteur financier invisible} de la mondialisation camerounaise. Son passage de \emph{filet de sécurité social} à \emph{moteur d'investissement} nécessite une réforme profonde du climat des affaires et une ingénierie financière innovante (Diaspora Bonds).
\end{enumerate}
\end{exoresolu}

\begin{methode}[Calculer et interpréter des indicateurs de libéralisation : Taux d'ouverture, Part de marché, IDE par habitant]
\textbf{Objectif :} Maîtriser les formules de base pour quantifier l'intégration économique et les comparer dans l'espace (Cameroun vs CEMAC vs SSA vs Monde) et dans le temps.
\begin{enumerate}
    \item \textbf{Taux d'ouverture commerciale (TOC) :} $\frac{X + M}{PIB} \times 100$.
    \begin{itemize}
        \item $X$ = Exportations (FOB), $M$ = Importations (CIF ou FOB selon source, préciser).
        \item Interprétation : $< 30\%$ = économie relativement fermée ; $30-60\%$ = modérément ouverte ; $> 60\%$ = économie très ouverte (ex: Singapour, Hong Kong, petits États insulaires). Cameroun $\approx$ 25\%.
    \end{itemize}
    \item \textbf{Taux de couverture (TC) :} $\frac{X}{M} \times 100$.
    \begin{itemize}
        \item $TC > 100\%$ = Excédent commercial. $TC < 100\%$ = Déficit. Cameroun structurellement $< 100\%$ (sauf années pétrolières hautes).
    \end{itemize}
    \item \textbf{Part des exportations primaires / manufacturées :} $\frac{\text{Export Primaires}}{\text{Export Total}} \times 100$. Indicateur de \textbf{spécialisation} (piège de la rente). Cameroun $> 80\%$ primaires (pétrole, bois, cacao, coton, aluminium).
    \item \textbf{Stock d'IDE par habitant :} $\frac{\text{Stock IDE entrant (Mds \$)}}{\text{Population (M hab.)}}$. Permet comparaison inter-pays. Cameroun $\approx$ 150-200 \$/hab (faible vs Gabon, Guinée Éq., Ghana).
    \item \textbf{Indice de concentration géographique (Hirschman-Herfindahl simplifié) :} Somme des carrés des parts de marché des 5 premiers partenaires. $> 0,18$ = forte concentration. Cameroun concentré sur UE + Chine.
\end{enumerate}
\cle{Formules} : TOC, TC, Part $\%$, Taux variation $\%$. \textbf{Unités :} Vérifier homogénéité (Mds FCFA vs Mds \$ ; PIB current vs constant). 1 \$ $\approx$ 600-650 FCFA (taux change fixe EUR/FCFA = 655,957).
\end{methode}

\begin{exoresolu}[Calculs et comparaison : Ouverture du Cameroun vs CEMAC (Données 2022 approx.)]
\textbf{Données 2022 (Sources : FMI, BEAC, CNUCED, estimations) :}
\begin{center}
\begin{tabularx}{\linewidth}{|l|X|X|X|X|}\hline
\rowcolor{popblueL}
\textbf{Pays} & \textbf{PIB (Mds \$)} & \textbf{Export (Mds \$)} & \textbf{Import (Mds \$)} & \textbf{Stock IDE entrant (Mds \$)} \\ \hline
Cameroun & 44 & 5,5 & 7,5 & 8,5 \\ \hline
Gabon & 20 & 10 & 3,5 & 6 \\ \hline
Tchad & 13 & 4 & 2,5 & 3 \\ \hline
Guinée Éq. & 12 & 8 & 2 & 4 \\ \hline
RCA & 2,5 & 0,2 & 0,5 & 0,1 \\ \hline
Congo & 12 & 5 & 3 & 5 \\ \hline
\end{tabularx}
\end{center}
Population 2022 : Cameroun 28 M, Gabon 2,4 M, Tchad 18 M, Guinée Éq. 1,7 M, RCA 5,5 M, Congo 5,8 M.
\textbf{Questions :}
\begin{enumerate}
    \item Calculer le TOC et le TC pour chaque pays. Classer les pays du plus ouvert au moins ouvert.
    \item Calculer le Stock IDE par habitant pour chaque pays. Commenter la position du Cameroun.
    \item Le Cameroun a un TOC faible mais un PIB total CEMAC le plus élevé. Expliquez ce paradoxe apparent.
\end{enumerate}
\tcblower
\textbf{Solution détaillée}
\begin{enumerate}
    \item \textbf{Calculs (Tableau de synthèse) :}
    \begin{center}
    \begin{tabularx}{\linewidth}{|l|X|X|X|}\hline
    \rowcolor{popblueL}
    \textbf{Pays} & \textbf{TOC (\%)} & \textbf{TC (\%)} & \textbf{IDE/hab (\$)} \\ \hline
    Cameroun & $\frac{5,5+7,5}{44}\times100 = \textbf{29,5}$ & $\frac{5,5}{7,5}\times100 = \textbf{73,3}$ & $\frac{8,5}{28} = \textbf{304}$ \\ \hline
    Gabon & $\frac{10+3,5}{20}\times100 = \textbf{67,5}$ & $\frac{10}{3,5}\times100 = \textbf{285,7}$ & $\frac{6}{2,4} = \textbf{2\,500}$ \\ \hline
    Tchad & $\frac{4+2,5}{13}\times100 = \textbf{50,0}$ & $\frac{4}{2,5}\times100 = \textbf{160,0}$ & $\frac{3}{18} = \textbf{167}$ \\ \hline
    Guinée Éq. & $\frac{8+2}{12}\times100 = \textbf{83,3}$ & $\frac{8}{2}\times100 = \textbf{400,0}$ & $\frac{4}{1,7} = \textbf{2\,353}$ \\ \hline
    RCA & $\frac{0,2+0,5}{2,5}\times100 = \textbf{28,0}$ & $\frac{0,2}{0,5}\times100 = \textbf{40,0}$ & $\frac{0,1}{5,5} = \textbf{18}$ \\ \hline
    Congo & $\frac{5+3}{12}\times100 = \textbf{66,7}$ & $\frac{5}{3}\times100 = \textbf{166,7}$ & $\frac{5}{5,8} = \textbf{862}$ \\ \hline
    \end{tabularx}
    \end{center}
    \textbf{Classement TOC (Ouverture) :} 1. Guinée Éq. (83\%), 2. Gabon (67,5\%), 3. Congo (66,7\%), 4. Tchad (50\%), 5. Cameroun (29,5\%), 6. RCA (28\%).
    \textbf{Analyse TC :} Seul le Cameroun, la RCA et (légèrement) le Tchad ont un TC $< 100\%$ (déficit). Les pays pétroliers (Gabon, G.Éq., Congo) ont des excédents massifs.
    \item \textbf{IDE/hab :} Le Cameroun (304 \$) est \textbf{dernier des pays pétroliers} et très loin du Gabon (2\,500 \$) ou Guinée Éq. (2\,353 \$). Il devance seulement le Tchad (167 \$) et la RCA (18 \$).
    \textit{Interprétation :} Malgré un marché intérieur vaste (28 M hab), le Cameroun attire peu d'IDE \emph{par tête}. L'IDE va vers l'extractif (pétrole/mines) dans les petits pays (rente facile) ou vers les services/telecoms (MTN/Orange) au Cameroun mais dilué par la population.
    \item \textbf{Paradoxe PIB vs Ouverture :}
    \begin{itemize}
        \item \textbf{Économie diversifiée (relative) :} Le PIB camerounais repose sur l'agriculture (15-20\%), l'industrie (25-30\% dont agro-alimentaire, BTP, aluminium), les services (45-50\%). La demande intérieure (consommation 28 M hab, BTP infrastructures) absorbe une grande partie de la production. Le commerce extérieur ne représente qu'une fraction de l'activité.
        \item \textbf{Pays pétroliers :} PIB = Rente pétrolière (exportée à 90\%) + peu d'activité locale $\rightarrow$ TOC mécaniquement très élevé.
        \item \textbf{Barrières :} Douanes lentes, corruption, infrastructure déficiente (énergie, routes), informel massif (réexport vers Nigeria non enregistré), barrières non tarifaires CEMAC.
    \end{itemize}
    \textbf{Conclusion :} Le Cameroun est la \textbf{première puissance économique de la CEMAC par le volume} (PIB, population, diversification), mais la \textbf{moins intégrée commercialement} (TOC le plus bas hors RCA en crise). Son défi : transformer la taille du marché en attractivité pour l'IDE productif (ZES, climat affaires) et réduire le déficit structurel via la transformation locale (bois, cacao, coton).
\end{enumerate}
\end{exoresolu}

\begin{methode}[Construire une carte mentale / schéma synthèse : Les 5 moteurs de la mondialisation au Cameroun]
\textbf{Objectif :} Schématiser les liens logiques entre les 5 facteurs du programme pour une révision efficace ou une introduction de dissertation.
\begin{enumerate}
    \item \textbf{Nœud central :} « \textbf{Facteurs de la mondialisation au Cameroun} ».
    \item \textbf{5 Branches principales (Couleurs distinctes) :}
    \begin{enumerate}
        \item \textbf{🟦 Transports \& Télécoms} (Physique \& Virtuel) : Ports (Douala, Kribi, Limbé, Garoua), Corridors (Route, Rail Camrail, Oléoduc), Aéroports, Câbles sous-marins (SAT-3, WACS, ACE, 2Africa), Backbone fibre, Data Centers (Yde, Dla), Mobile Money.
        \item \textbf{🟧 Firmes Multinationales (FMN)} : Extractives (SNH/Perenco/TotalEnergies, Alucam/Rio Tinto, Sundance/Sinosteel), Agro-industrie (PHP/Compagnie Fruitière, CDC, SOCAPALM, Safacam), Bois (Rougier, Pallisco, SEFAC), Services/Logistique (MSC/TiL/Camrail/Douala Terminal/KCT/CHEC, Maersk/MSC/CMA CGM), Télécoms (MTN, Orange, Camtel/Viettel), Banque (Afriland, UBA, Ecobank, BICEC/BNP).
        \item \textbf{🟩 État \& Diaspora} : État (SND30, ZES Kribi, Code Investissement, Douanes, Régulation ARMP/ART), Diaspora (Transferts 300-900 Mds FCFA, PAID, Caravane, Investissement immobilier/PME, Lobbying).
        \item \textbf{🟪 Associations \& ONG} : Internationales (ONU : OMS, UNICEF, PNUD, HCR, PAM ; ONG : MSF, CRS, Care, Plan International, WWF, Greenpeace), Nationales (RELUFA, Réseau Lutte contre la Faim, ACDIC, Dynamique Citoyenne). Rôle : Service public substitutif (Nord/NO/SO), Plaidoyer (gouvernance, environnement, droits humains), Capacitation.
        \item \textbf{🟥 Commerce International} : Accords (OMC, APE UE-CEMAC, ZLECAf, AGOA), Structure (Export primaires 80\%+, Import machines/produits finis), Partenaires (UE 40\%, Chine 25\%, CEMAC 10\%), Déficit chronique, Enjeu transformation locale (2ème/3ème transformation bois, cacao, coton).
    \end{enumerate}
    \item \textbf{Liens transversaux (Flèches entre branches) :}
    \begin{itemize}
        \item FMN $\rightarrow$ Transports (investissent/usent ports, rail, fibre).
        \item État $\rightarrow$ Transports (finance infrastructures, concède à FMN).
        \item Diaspora $\rightarrow$ Commerce (import véhicules, matériel, financement PME export).
        \item ONG $\rightarrow$ État (partenariat santé/éducation, contrôle transparence extractive - ITIE).
        \item Télécoms $\rightarrow$ Diaspora/Commerce (Mobile Money, e-commerce, transfert fonds).
    \end{itemize}
    \item \textbf{Mise en page :} A3 paysage, centre aéré, mots-clés courts, pictogrammes, flèches de relation, légende des couleurs.
\end{enumerate}
\cle{Réflexe} : Une carte mentale n'est pas un cours recopié. C'est une \textbf{hiérarchisation visuelle}. Au Bac, elle peut servir de plan de dissertation ou d'illustration en géographie.
\end{methode}

\begin{exoresolu}[Réalisation d'une carte mentale synthétique : « Le Cameroun face à la mondialisation »]
\textbf{Consigne :} Reproduisez la carte mentale décrite dans la méthode ci-dessus sur une feuille A3 (ou format numérique). Vous devez faire apparaître : les 5 facteurs, 3 exemples précis par facteur (noms propres, lieux, chiffres), et 3 flèches d'interaction entre facteurs.
\tcblower
\textbf{Solution détaillée (Structure textuelle de la carte mentale à reproduire graphiquement)}

\textbf{CENTRE :} \fbox{\textbf{FACTEURS DE LA MONDIALISATION AU CAMEROUN (2024)}}

\textbf{BRANCHE 1 (BLEU) : TRANSPORTS \& TÉLÉCOMS}
\begin{itemize}
    \item \textbf{Ports :} Douala (1er, 90\% trafic, saturation), Kribi (Eau profonde 16m, KCT/CHEC, GNL SNH/Perenco), Limbé (Pétrole/SONARA), Garoua (Fluvial/Bénoué).
    \item \textbf{Corridors :} Douala-Ngaoundéré (Camrail/MSC/TiL + Route N1/N15), Douala-Bertoua-Ngaoundéré (Route), Kribi-Lolodorf (Autoroute 2x2 voies), Projet Rail Mbalam-Kribi (Fer/Sinosteel).
    \item \textbf{Numérique :} Câbles SAT-3/WACS/ACE/2Africa (atterrissement Kribi/Limbé/Douala), Backbone national 10 000 km fibre (Camtel), Data Centers (Yaoundé, Douala - MTN/Orange/Camtel), Mobile Money (MoMo, Orange Money $\approx$ 10 M utilisateurs).
\end{itemize}

\textbf{BRANCHE 2 (ORANGE) : FIRMES MULTINATIONALES (FMN)}
\begin{itemize}
    \item \textbf{Extractif :} SNH/Perenco/TotalEnergies (Pétrole/Gaz - Rio del Rey, Logbaba, Kribi), Alucam/Rio Tinto (Aluminium - Édéa), Sundance/Sinosteel (Fer - Mbalam/Nabeba, *projet*).
    \item \textbf{Agro/Bois :} PHP/Compagnie Fruitière (Banane - Tiko, Moungo, 300k t/an), CDC (Palmier/Hévéa - Sud-Ouest), SOCAPALM/MSC/TiL (Palmier - Kienké, Dibombari, Mbongo), Rougier/Pallisco/SEFAC (Bois - Est/Sud, 1er export bois Afrique).
    \item \textbf{Services/Logistique :} MSC/TiL (Camrail, Douala Terminal, Logistique), CHEC (KCT Kribi), Maersk/MSC/CMA CGM (Armateurs), MTN/Orange/Camtel (Télécoms - 90\% parts marché), UBA/Ecobank/Afriland/BICEC (Banque).
\end{itemize}

\textbf{BRANCHE 3 (VERT) : ÉTAT \& DIASPORA}
\begin{itemize}
    \item \textbf{État :} SND30 (Stratégie Nationale Développement), ZES Kribi (exonérations 10 ans), Code Investissement 2013, Douanes (CAMCIS, guichet unique), ARMP (Marchés publics), ART (Régulation Télécoms), ITIE (Transparence extractive).
    \item \textbf{Diaspora :} 3-4 M personnes (France, USA, Allemagne, Canada, Pays Golfe). Transferts : 300 Mds FCFA (formel) $\rightarrow$ 900 Mds (estimé total). 1er pourvoyeur devises. PAID (Programme Appui Investissements Diaspora), Caravane Diaspora. Investissement : Immobilier, PME Agro/Tech, Santé.
\end{itemize}

\textbf{BRANCHE 4 (VIOLET) : ASSOCIATIONS \& ONG}
\begin{itemize}
    \item \textbf{Internationales (Humanitaire/Dév) :} OMS/UNICEF/PNUD/HCR/PAM (ONU), MSF/Médecins du Monde (Santé urgences NO/SO/Extrême-Nord), CRS/Caritas/Plan International (Éducation/Protection), WWF/Greenpeace/Global Witness (Environnement/Bois/Forêt).
    \item \textbf{Nationales :} RELUFA (Lutte faim), ACDIC (Développement communautaire), Réseau Citoyen (Gouvernance), SYNATRUC (Transport), Conférences Épiscopales (Médiation).
    \item \textbf{Rôle :} Suppléance État (Nord/NO/SO), Plaidoyer (ZLECAf, Dette, Climat, Droits Humains), Expertise technique.
\end{itemize}

\textbf{BRANCHE 5 (ROUGE) : COMMERCE INTERNATIONAL}
\begin{itemize}
    \item \textbf{Cadre :} OMC, APE UE-CEMAC (2014/2016, libéralisation 80\% lignes tarifaires), ZLECAf (ratifiée 2021, démarrage opérationnel), AGOA (USA - éligible textile/agro).
    \item \textbf{Structure :} Export : Pétrole (40\%), Bois (20\%), Cacao (10\%), Aluminium, Coton, Banane. Import : Produits pétroliers raffinés, Machines, Céréales (Riz/Blé), Médicaments, Véhicules.
    \item \textbf{Partenaires :} UE (Client 43\%, Fournisseur 38\%), Chine (Fournisseur 28\%, Client 14\%), CEMAC (Client 11\%, Fournisseur 4\%), Inde, USA.
    \item \textbf{Enjeu :} Déficit structurel ($\approx$ -1000 Mds FCFA/an), Transformation locale (Décrets 2017/2023 interdiction export grumes/cacao brut/coton brut).
\end{itemize}

\textbf{INTERACTIONS (FLÈCHES TRANSVERSALES - Exemples) :}
\begin{enumerate}
    \item \textcolor{poporange}{FMN (MSC/TiL/CHEC)} $\xrightarrow{\text{Concession/Investissement}}$ \textcolor{popblue}{Transports (Camrail, KCT Kribi)} $\xrightarrow{\text{Désenclavement}}$ \textcolor{popred}{Commerce (Transit Tchad/RCA)}.
    \item \textcolor{popgreen}{État (ZES Kribi, Code Invest)} $\xrightarrow{\text{Cadre légal/Fiscal}}$ \textcolor{poporange}{FMN (Sinosteel, PHP, MTN)} $\xrightarrow{\text{Emploi/Transfert tech}}$ \textcolor{popgreen}{État (Recettes fiscales)}.
    \item \textcolor{popgreen}{Diaspora (Transferts MoMo/Wave)} $\xrightarrow{\text{Télécoms (MTN/Orange/Camtel)}}$ \textcolor{popblue}{Télécoms (Data Centers, Fibre)} $\xrightarrow{\text{Inclusion financière}}$ \textcolor{popred}{Commerce (E-commerce, PME)}.
\end{enumerate}
\end{exoresolu}

\begin{attention}[Les 5 erreurs qui coûtent des points au Bac]
\begin{enumerate}
    \item \textbf{Confondre « Mondialisation » et « Libéralisation » :} La mondialisation est le \textbf{processus global} d'interdépendance (flux multiples). La libéralisation des échanges est une \textbf{politique délibérée} (baisse tarifs, suppression barrières) qui \emph{accélère} la mondialisation. Ne pas écrire « La mondialisation a libéralisé les échanges » mais « La libéralisation \emph{renforce} la mondialisation ».
    \item \textbf{Oublier le « Cameroun » dans les exemples :} Le programme impose l'\textbf{ancrage territorial}. Citer « les FMN exportent des matières premières » sans nommer \textbf{SNH, PHP, Alucam, Rougier, SOCAPALM, MSC/TiL, MTN, CHEC} fait perdre la moitié des points d'illustration. De même pour les ports (Douala/Kribi), corridors (Camrail, N1), ONG (MSF au Nord, WWF dans l'Est).
    \item \textbf{Traiter la Diaspora uniquement comme « envoi d'argent » :} C'est un \textbf{acteur géographique à part entière} : réseaux transnationaux, transfert de compétences (retour pays, télétravail), investissement foncier/immobilier (reconfiguration urbaine Yaoundé/Douala/Bafoussam), lobbying politique (vote extérieur, pression diplomatique). Il faut analyser la \textbf{spatialité} des flux (Origine $\rightarrow$ Destination, Usage).
    \item \textbf{Ne pas hiérarchiser les facteurs dans la synthèse :} Tous les facteurs ne pèsent pas le même poids. \textbf{Ordre de grandeur :} 1. Transports/Télécoms (infrastructure physique sine qua non), 2. FMN (opérateurs principaux flux/IDE), 3. Commerce (cadre règlementaire/résultat), 4. État (régulateur/stratège), 5. Diaspora/ONG (acteurs complémentaires, flux financiers/sociaux). Une conclusion qui met tout sur le même plan manque de rigueur.
    \item \textbf{Erreurs cartographiques / Croquis :} \textbf{(a)} Légende absente ou illisible (pas de cadre, symboles non reproduits). \textbf{(b)} Flèches de flux toutes de même épaisseur (ne pas distinguer flux massifs conteneurs Douala$\rightarrow$Ngaoundéré vs flux diaspora ponctuels). \textbf{(c)} Oublier l'orientation (Nord) et l'échelle indicative. \textbf{(d)} Placer Kribi au Nord de Douala ou Ngaoundéré au Sud de Yaoundé (respecter la géométrie de base du triangle Cameroun).
\end{enumerate}
\end{attention}

\newpage

\coursec{Exercices}

\courssub{Je vérifie mes connaissances}

\begin{exercice}[QCM : Les moteurs de la mondialisation]
Pour chaque question, une seule réponse est exacte. Recopie le numéro et la lettre correspondante.
\begin{enumerate}
    \item Lequel des éléments suivants n'est \textbf{pas} un facteur technique de la mondialisation ?
    \begin{enumerate}
        \item Le conteneur normalisé (ISO)
        \item Le câble sous-marin à fibre optique
        \item La convention de l'OMC sur les subventions
        \item L'avion gros-porteur (ex. Boeing 747, Airbus A380)
    \end{enumerate}
    \item Une firme multinationale (FMN) se caractérise principalement par :
    \begin{enumerate}
        \item Son siège social situé dans un pays en développement
        \item La détention d'actifs productifs dans au moins deux pays
        \item Son statut d'organisme non gouvernemental
        \item Son financement exclusif par la diaspora
    \end{enumerate}
    \item Les transferts de fonds de la diaspora camerounaise vers le pays (environ \num{300} milliards de FCFA/an ces dernières années) relèvent de :
    \begin{enumerate}
        \item L'aide publique au développement (APD)
        \item L'investissement direct étranger (IDE)
        \item Les envois de fonds privés (remittances)
        \item Le commerce international de services
    \end{enumerate}
    \item L'Organisation mondiale du commerce (OMC) a pour rôle principal de :
    \begin{enumerate}
        \item Fixer le prix du baril de pétrole
        \item Réguler les flux migratoires internationaux
        \item Négocier les accords de libre-échange et régler les différends commerciaux
        \item Contrôler les budgets des États membres
    \end{enumerate}
\end{enumerate}
\end{exercice}

\begin{exercice}[Vrai ou Faux ? Justifie chaque réponse]
Indique si l'affirmation est vraie (V) ou fausse (F) et justifie brièvement en mobilisant tes connaissances du cours.
\begin{enumerate}
    \item La révolution des télécommunications (Internet, téléphonie mobile) a réduit les coûts de transaction de l'information à quasi zéro, accélérant la fragmentation internationale des processus de production.
    \item Les firmes multinationales n'investissent dans les pays du Sud que pour extraire des matières premières (pétrole, minerais) sans créer d'emplois locaux qualifiés.
    \item L'État camerounais, à travers la loi sur les investissements privés (loi n°2013/004) et la création de zones industrielles (ex. Kribi, Pola), joue un rôle actif dans l'attractivité du territoire pour la mondialisation.
    \item Les ONG humanitaires (ex. Croix-Rouge, MSF) et les associations de développement (ex. ONG locales d'appui aux groupements paysans) sont des acteurs non étatiques qui tissent des réseaux transnationaux, participant ainsi à la mondialisation par le bas.
\end{enumerate}
\end{exercice}

\begin{exercice}[Définitions et concepts clés]
Donne une définition précise (2 à 3 lignes) pour chacun des termes suivants en géographie de la mondialisation :
\begin{enumerate}
    \item \cle{Mondialisation} (distinction avec \cle{internationalisation})
    \item \cle{Investissement direct étranger (IDE)} (critère de propriété \SI{10}{\%} du capital)
    \item \cle{Chaîne de valeur mondiale (CVM)} ou \cle{fragmentation productive}
    \item \cle{Diaspora} (au sens géographique et économique)
    \item \cle{Libéralisation des échanges} (rôle de l'OMC, baisse des droits de douane)
\end{enumerate}
\end{exercice}

\begin{exercice}[Calculs directs : Spécialisation commerciale du Cameroun]
Le tableau ci-dessous présente les exportations camerounaises par produit principal (données arrondies, source : INS / Douanes, exercices récents).
\begin{center}
\begin{tabular}{|l|c|c|}\hline
\textbf{Produit} & \textbf{Valeur (milliards FCFA)} & \textbf{Part (\%)} \\ \hline
Pétrole brut & 1\,800 & \\ \hline
Cacao (fèves + dérivés) & 450 & \\ \hline
Bois (grumes, sciages) & 300 & \\ \hline
Coton (fibre) & 120 & \\ \hline
Aluminium & 100 & \\ \hline
Autres (banane, café, caoutchouc, etc.) & 230 & \\ \hline
\textbf{Total} & \textbf{3\,000} & \textbf{100} \\ \hline
\end{tabular}
\end{center}
\begin{enumerate}
    \item Calcule la part (\% ) de chaque produit dans le total des exportations (complète le tableau).
    \item Calcule le \cle{taux de couverture} des exportations par le pétrole seul.
    \item Commenter en une phrase la structure de ces exportations (spécialisation, vulnérabilité).
\end{enumerate}
\end{exercice}

\courssub{Je m'entraîne}

\begin{exercice}[Analyse de document : La connectivité maritime mondiale]
\textbf{Document :} Extrait simplifié du trafic conteneurisé des 5 premiers ports mondiaux (2023, millions d'EVP\footnote{EVP = Équivalent Vingt Pieds, unité de mesure standard d'un conteneur.}).
\begin{center}
\begin{tabular}{|l|c|c|}\hline
\textbf{Port} & \textbf{Pays / Région} & \textbf{Trafic (millions EVP)} \\ \hline
Shanghai & Chine (Asie de l'Est) & 49,1 \\ \hline
Singapour & Singapour (Asie du Sud-Est) & 39,0 \\ \hline
Ningbo-Zhoushan & Chine (Asie de l'Est) & 35,3 \\ \hline
Shenzhen & Chine (Asie de l'Est) & 29,9 \\ \hline
Busan & Corée du Sud (Asie de l'Est) & 22,8 \\ \hline
\end{tabular}
\end{center}
\textit{Premier port africain : Tanger Med (Maroc) \textasciitilde 8,6 millions EVP ; Douala (Cameroun) \textasciitilde 0,5 million EVP.}
\begin{enumerate}
    \item Localise ces ports sur un planisphère mental : quelle région du monde domine le trafic conteneurisé mondial ?
    \item Calcule la part cumulée de la Chine (Shanghai + Ningbo + Shenzhen) dans le trafic de ce top 5.
    \item Explique pourquoi l'Afrique (hors Tanger Med) reste marginale dans les flux conteneurisés mondiaux (au moins 3 facteurs : infrastructure, volume d'échanges, position dans les CVM).
\end{enumerate}
\end{exercice}

\begin{exercice}[Étude de cas : Stratégies d'implantation de MTN Cameroon]
\textbf{Contexte :} MTN Cameroon (filiale du groupe sud-africain MTN Group, leader télécoms en Afrique) est le premier opérateur mobile au Cameroun (\textasciitilde 18 millions d'abonnés, \textasciitilde 50\% de part de marché, chiffre d'affaires \textasciitilde 400 milliards FCFA/an).
\begin{enumerate}
    \item Identifie la nationalité du groupe mère et le secteur d'activité (biens ou services ?).
    \item Cite trois stratégies concrètes d'implantation locale observables (ex. : recrutement, RSE, infrastructure, adaptation de l'offre \emph{Mobile Money}).
    \item Montre, par deux exemples chiffrés ou factuels, comment cette FMN intègre le Cameroun aux flux mondiaux (financiers, technologiques, informationnels).
\end{enumerate}
\end{exercice}

\begin{exercice}[Rôle de l'État et de la diaspora : Le corridor Kribi -- Mbalam]
\textbf{Document :} Le port en eau profonde de Kribi (phase 1 opérationnelle depuis 2014, géré par \emph{Kribi Conteneur Terminal} -- KCT, consortium international) et le projet ferroviaire Kribi--Mbalam (fer de Mbalam, opérateur \emph{Sundance Resources} / \emph{Australia} puis consortium chinois).
\begin{enumerate}
    \item En quoi l'État camerounais (via le Port Autonome de Kribi - PAK et les conventions de concession) agit-il comme \textbf{acteur d'intégration} à la mondialisation ?
    \item La diaspora camerounaise finance-t-elle ce type d'infrastructure ? Précise son rôle réel (transferts vers ménages, investissement immobilier/PME, lobbying diplomatique).
    \item Pourquoi le choix du partenaire gestionnaire (KCT : \emph{ICTSI} Philippines / \emph{Maersk} Danemark / \emph{Bolloré Logistics} France) illustre-t-il la mondialisation des services portuaires ?
\end{enumerate}
\end{exercice}

\begin{exercice}[Croquis commenté : Les flux de la mondialisation et la place de l'Afrique]
À partir de la description ci-dessous, réalise un \textbf{croquis schématique} (format A4 paysage) intitulé \og La place de l'Afrique dans les flux de la mondialisation (années 2020) \fg.
\textbf{Données à représenter :}
\begin{itemize}
    \item \textbf{Triade} (Amérique du Nord, Europe de l'Ouest, Asie de l'Est) : flux massifs bidirectionnels (flèches épaisses), pôles de commandement (FMN, Bourses, R\&D).
    \item \textbf{Pays émergents} (Chine, Inde, Brésil, ASEAN) : flux croissants vers Triade et vers Afrique (matières premières), flux d'IDE vers Afrique (infrastructures).
    \item \textbf{Afrique} : exportations de matières premières (pétrole, minerais, cacao, bois) vers Triade et Chine (flèches moyennes sortantes) ; importations de produits manufacturés, céréales, hydrocarbures raffinés (flèches moyennes entrantes) ; IDE entrants (infrastructures, télécoms) ; envois de fonds diaspora (flèches pointillées entrantes).
    \item \textbf{Cameroun} : positionner Douala/Kribi (ports), Yaoundé (capitale), pipeline Tchad--Kribi, axe routier Douala--N'Djamena.
\end{itemize}
\textbf{Exigences :} Légende organisée (Acteurs, Flux, Territoires), figurés de surface/linéaires/ponctuels respectés, flèche Nord, échelle approximative.
\end{exercice}

\courssub{Je me prépare au Bac}

\begin{exercice}[Sujet type Bac : Dissertation géographique]
\textbf{Sujet :} \og Les firmes multinationales sont-elles les principaux acteurs de la mondialisation ? \fg
\textbf{Consignes :}
\begin{enumerate}
    \item Analyse le sujet (mots-clés : \cle{principaux}, \cle{acteurs}, \cle{mondialisation} ; délimitation spatio-temporelle).
    \item Propose une \textbf{problématique} pertinente.
    \item Construis un \textbf{plan détaillé} (2 ou 3 parties, 2 sous-parties chacune) avec arguments, exemples camerounais et mondiaux (FMN : TotalEnergies, MTN, Dangote, Huawei ; États : Chine, USA, Cameroun ; Diaspora, ONG, OMC).
    \item Rédige l'\textbf{introduction} complète (accroche, définition, sujet, problématique, annonce du plan).
\end{enumerate}
\end{exercice}

\begin{exercice}[Sujet type Bac : Étude de documents -- Géographie du commerce international]
\textbf{Document 1 :} Carte schématique des principales routes maritimes mondiales (description textuelle).
\og Les flux maritimes structurent l'espace mondial : \textbf{axe Pacifique Nord} (Chine/Japon/Corée -- Côte Ouest USA) ; \textbf{axe Atlantique Nord} (Europe -- Côte Est USA) ; \textbf{axe Europe -- Asie} (via Suez -- Mer Rouge -- Océan Indien -- Malacca) ; \textbf{axe Amérique du Sud -- Chine} (soja, minerais via Cap Horn ou Panama). L'Afrique est traversée par l'axe Europe--Asie (côte Ouest : Golfe de Guinée ; côte Est : Cap) mais génère peu de flux propres.\fg

\textbf{Document 2 :} Part des régions dans les exportations mondiales de marchandises (2023, OMC).
\begin{center}
\begin{tabular}{|l|c|}\hline
\textbf{Région} & \textbf{Part (\%)} \\ \hline
Asie (dont Chine \textasciitilde 14\%) & 38 \\ \hline
Europe (UE 27 \textasciitilde 15\%) & 36 \\ \hline
Amérique du Nord & 14 \\ \hline
Afrique & 2,3 \\ \hline
Amérique latine & 5 \\ \hline
\end{tabular}
\end{center}

\textbf{Questions :}
\begin{enumerate}
    \item \textbf{Document 1 :} Décris la géographie des grandes routes maritimes. Pourquoi le détroit de Malacca et le canal de Suez sont-ils des \cle{points de passage obligés} (goulets d'étranglement) ?
    \item \textbf{Document 2 :} Calcule la part cumulée de la Triade (Asie développée + Europe + Am. Nord) et commente la place de l'Afrique.
    \item \textbf{Synthèse :} En t'appuyant sur les documents et tes connaissances, explique pourquoi le commerce international reste fortement \cle{hiérarchisé} et \cle{régionalisé} (blocs régionaux : ALENA/USMCA, UE, ASEAN, ZLECAf) malgré la mondialisation.
\end{enumerate}
\end{exercice}

\begin{exercice}[Sujet type Bac : TP Noté -- Dossier « Le Cameroun dans la mondialisation »]
Tu disposes du dossier suivant (données simplifiées pour l'examen) :
\begin{itemize}
    \item \textbf{Exportations 2023 (milliards FCFA) :} Pétrole 1\,800 ; Cacao 450 ; Bois 300 ; Coton 120 ; Aluminium 100 ; Autres 230. Total 3\,000.
    \item \textbf{Importations 2023 (milliards FCFA) :} Produits pétroliers raffinés 800 ; Machines/Électrique 600 ; Céréales/Riz/Blé 400 ; Véhicules 300 ; Médicaments/Pharma 250 ; Autres 1\,150. Total 3\,500.
    \item \textbf{IDE entrants (stock 2022) :} \textasciitilde 3\,500 milliards FCFA (secteurs : hydrocarbures, télécoms, banques, ciment, agro-industrie). Principaux pays investisseurs : France, Chine, Maroc, Afrique du Sud, USA.
    \item \textbf{Transferts diaspora (2023) :} \textasciitilde 320 milliards FCFA (source : Banque Mondiale).
    \item \textbf{Infrastructures :} Port de Kribi (eau profonde), Pipeline Tchad--Kribi, Autoroute Yaoundé--Douala (en cours), Réseau 4G/5G (MTN, Orange, Camtel), Câble sous-marin SAT-3/WASC + ACE + WIOCC.
\end{itemize}

\textbf{Travail demandé :}
\begin{enumerate}
    \item \textbf{Traitement statistique (6 pts) :} Construis un \textbf{diagramme en barres groupées} comparant Exportations et Importations par grande catégorie (Hydrocarbures, Produits agricoles/bois, Produits manufacturés, Autres). Utilise une échelle adaptée. Titre, axes, unités, source.
    \item \textbf{Analyse de la structure commerciale (6 pts) :} Commente le déséquilibre de la balance commerciale (déficit \textasciitilde 500 Mds FCFA). Caractérise la \cle{spécialisation} du Cameroun (type de produits, élasticité-prix, vulnérabilité). Cite le concept de \cle{termes de l'échange}.
    \item \textbf{Acteurs et intégration (8 pts) :} Montre comment la combinaison \textbf{FMN (ex. TotalEnergies/Perenco, MTN, Cimencam/LafargeHolcim) + État (politique d'attractivité, ZLECAf, infrastructures) + Diaspora} structure l'insertion du Cameroun dans la mondialisation. Distingue les apports (capitaux, techno, emplois, devises, infrastructures) et les limites (rapatriement bénéfices, enclaves, dépendance technologique, endettement).
\end{enumerate}
\end{exercice}

\newpage

\coursec{Corrigés des exercices}

\coursec{Leçon 11 : Les facteurs de la mondialisation — Corrigés détaillés des exercices et sujets types}

\begin{corrige}[Exercice 1 — QCM : Les moteurs de la mondialisation]
\textbf{Réponses : 1 -- c ; 2 -- b ; 3 -- c ; 4 -- c.}

\textbf{Justifications :}
\begin{enumerate}
    \item \textbf{c.} Le conteneur ISO, la fibre optique et l'avion gros-porteur sont des \cle{innovations techniques} majeures qui réduisent drastiquement les coûts et les délais de transport (maritime, aérien) et de communication (internet). La convention de l'OMC sur les subventions relève du cadre \cle{institutionnel et juridique} (libéralisation des échanges), non du facteur technique.
    \item \textbf{b.} Selon la définition de la CNUCED, une FMN est une entreprise qui détient ou contrôle des \cle{actifs productifs} (usines, filiales, plantations, mines) dans au moins deux pays. Son siège social est le plus souvent situé dans un pays développé ou émergent ; c'est un acteur \emph{privé à but lucratif} (ni ONG, ni structure financée par la diaspora).
    \item \textbf{c.} Les envois de fonds des migrants (\emph{remittances}) sont des transferts \cle{privés} de personne à personne (ménage à ménage). Ce n'est ni de l'aide publique au développement (État à État), ni un investissement direct étranger (IDE, pas de prise de contrôle d'entreprise), ni un échange commercial de services.
    \item \textbf{c.} L'OMC (créée en 1995, succédant au GATT 1947) a pour missions la négociation de la baisse des barrières douanières et le règlement des différends entre États membres via son Organe de règlement des différends (ORD). Elle ne fixe pas les prix du pétrole (rôle de l'OPEP et des marchés), ne gère pas les flux migratoires, ni les budgets nationaux.
\end{enumerate}
\end{corrige}

\begin{corrige}[Exercice 2 — Vrai ou Faux ?]
\begin{enumerate}
    \item \textbf{Vrai.} Internet, la téléphonie mobile et les câbles sous-marins (SAT-3/WASC, ACE pour le Cameroun) ont fait chuter le coût de transmission de l'information. Les FMN peuvent ainsi piloter en temps réel des usines dispersées sur plusieurs continents : c'est la \cle{fragmentation internationale de la production} (ou nouvelles divisions internationales du travail), base des chaînes de valeur mondiales.
    \item \textbf{Faux.} L'extractivisme existe (pétrole, minerais), mais les FMN investissent massivement dans les télécoms (MTN, Orange), la banque, le ciment (LafargeHolcim, Dangote Cement), l'agro-industrie (SOCAPALM, PHP) et créent des emplois locaux, y compris qualifiés (ingénieurs, cadres, techniciens). L'affirmation est trop générale ; la critique économique pertinente porte sur le \cle{rapatriement des bénéfices} (dividendes), la faible intégration locale (effet d'enclave) et l'insuffisance des transferts de technologie.
    \item \textbf{Vrai.} La loi n°2013/004 incitative à l'investissement privé, la création de l'APIPA (Agence de Promotion des Investissements), la construction du port en eau profonde de Kribi et l'aménagement de zones industrielles (Bomono, Edea, Kribi) montrent que l'\cle{État aménage l'attractivité} du territoire. Il reste un acteur majeur de la mondialisation, contrairement à l'idée reçue d'un État « effacé » ou impuissant.
    \item \textbf{Vrai.} Les ONG internationales (MSF, Croix-Rouge, WWF) et locales (appui aux groupements de producteurs de cacao, coton, café) forment des \cle{réseaux transnationaux} d'idées, de normes (environnementales, sociales) et de financements. C'est la \og mondialisation par le bas \fg{}, complémentaire de la mondialisation économique pilotée par les FMN.
\end{enumerate}
\end{corrige}

\begin{corrige}[Exercice 3 — Définitions et concepts clés]
\begin{enumerate}
    \item \textbf{Mondialisation :} processus d'intensification et d'accélération des flux (biens, capitaux, informations, personnes) qui met les territoires du monde en interdépendance et crée un \emph{système} mondial hiérarchisé. L'\cle{internationalisation} désigne seulement le développement des échanges \emph{entre} États qui restent les cadres politiques et juridiques fermés ; la mondialisation suppose l'émergence d'acteurs transnationaux (FMN, ONG, réseaux financiers, diasporas) qui opèrent \emph{au-delà} et \emph{à travers} les frontières.
    \item \textbf{IDE (Investissement Direct Étranger) :} investissement réalisé par une entreprise d'un pays dans une entreprise d'un autre pays, avec prise de participation d'au moins \SI{10}{\%} du capital social (seuil OCDE/CNUCED). Ce seuil traduit une \cle{volonté de contrôle} durable sur la gestion (siège au conseil d'administration), et non un simple placement financier de portefeuille (spéculatif, à court terme).
    \item \textbf{Chaîne de valeur mondiale (CVM) :} découpage du processus de production d'un bien en étapes distinctes (R\&D/conception, fabrication de composants, assemblage final, logistique, marketing, services après-vente) localisées dans plusieurs pays selon leurs avantages comparatifs (coût du travail, compétences, ressources, fiscalité), sous le pilotage stratégique d'une FMN \cle{entreprise pilote}. Exemple : smartphone conçu aux USA (Apple), puces de Taïwan (TSMC), écrans de Corée (Samsung/LG), assemblé en Chine/Inde (Foxconn), vendu mondialement.
    \item \textbf{Diaspora :} ensemble des personnes originaires d'un pays vivant durablement à l'étranger et maintenant des liens forts (identitaires, économiques, politiques) avec leur pays d'origine. Économiquement, elle est un acteur majeur par ses \cle{envois de fonds} (\textasciitilde 300 à 350 milliards FCFA/an pour le Cameroun, source Banque mondiale), ses investissements directs (immobilier, PME, agriculture), et ses transferts de compétences (retour d'expertise, télétravail, lobbying).
    \item \textbf{Libéralisation des échanges :} réduction progressive, négociée et multilatérale des obstacles au commerce international (droits de douane, quotas, barrières non tarifaires) engagée depuis le GATT (1947) puis l'OMC (1995), afin de favoriser la libre circulation des biens et services. Elle se traduit aussi par des accords régionaux préférentiels (UE, USMCA, ASEAN, \cle{ZLECAf} pour l'Afrique entrée en vigueur 2021).
\end{enumerate}
\end{corrige}

\begin{corrige}[Exercice 4 — Calculs : Spécialisation commerciale du Cameroun]
\textbf{1. Parts calculées} (part = valeur / total $\times$ 100, total exportations = \num{3000} milliards FCFA) :
\begin{center}
\begin{tabular}{|l|c|c|}\hline
\textbf{Produit} & \textbf{Valeur (Mds FCFA)} & \textbf{Part (\%)\ } \\ \hline
Pétrole brut & 1\,800 & $\frac{1800}{3000}\times100 = \SI{60,0}{\%}$ \\ \hline
Cacao & 450 & $\frac{450}{3000}\times100 = \SI{15,0}{\%}$ \\ \hline
Bois & 300 & $\frac{300}{3000}\times100 = \SI{10,0}{\%}$ \\ \hline
Coton & 120 & $\frac{120}{3000}\times100 = \SI{4,0}{\%}$ \\ \hline
Aluminium & 100 & $\frac{100}{3000}\times100 \approx \SI{3,3}{\%}$ \\ \hline
Autres & 230 & $\frac{230}{3000}\times100 \approx \SI{7,7}{\%}$ \\ \hline
\textbf{Total} & \textbf{3\,000} & \textbf{100,0} \\ \hline
\end{tabular}
\end{center}
Vérification : $60,0 + 15,0 + 10,0 + 4,0 + 3,3 + 7,7 = 100,0$. Correct.

\textbf{2. Taux de couverture par le pétrole seul :} le pétrole représente \SI{60,0}{\%} des exportations totales ; il « couvre » à lui seul 60\,\% de la valeur totale exportée. Rapporté aux importations (\num{3500} Mds FCFA), il couvrirait $\frac{1800}{3500}\times100 \approx \SI{51,4}{\%}$ des importations.

\textbf{3. Commentaire :} les exportations camerounaises sont très \cle{concentrées} (pétrole + cacao + bois = \SI{85}{\%}) sur des matières premières brutes ou peu transformées (premières transformations : bois sciage, fèves de cacao, coton fibre). Cette spécialisation expose le pays à la \cle{volatilité des cours mondiaux} (vulnérabilité extérieure : choc pétrolier 2014-2016, chute cacao 2017) et rapporte peu de valeur ajoutée locale (emplois peu qualifiés, faible fiscalité sur les volumes).

\textbf{Points de vigilance :} ne pas oublier de vérifier que la somme des parts fait exactement \SI{100}{\%} ; arrondir à un chiffre après la virgule (\SI{3,3}{\%}, \SI{7,7}{\%}) et le signaler explicitement ; distinguer \emph{taux de couverture des importations par les exportations} (ici \SI{85,7}{\%}) et \emph{part d'un produit dans les exportations}.
\end{corrige}

\begin{corrige}[Exercice 5 — Analyse de document : La connectivité maritime mondiale]
\textbf{1.} Les cinq premiers ports mondiaux par trafic conteneurisé (EVP) sont tous situés en \cle{Asie de l'Est et du Sud-Est} (façade pacifique asiatique : Shanghai, Singapour, Ningbo-Zhoushan, Shenzhen, Guangzhou/Qingdao). C'est cette région, \og atelier du monde \fg{}, qui domine les échanges manufacturés conteneurisés.

\textbf{2. Part cumulée de la Chine (hors Hong Kong) dans le top 5 :}
\begin{align*}
\text{Trafic ports chinois (Shanghai, Ningbo, Shenzhen)} &= 49,1 + 35,3 + 29,9 = 114,3 \text{ millions EVP} \\
\text{Total top 5} &= 49,1 + 39,0 + 35,3 + 29,9 + 22,8 = 176,1 \text{ millions EVP} \\
\text{Part} &= \frac{114,3}{176,1} \times 100 \approx \SI{64,9}{\%} \approx \SI{65}{\%}.
\end{align*}
La Chine concentre à elle seule près des deux tiers du trafic des cinq premiers ports mondiaux, confirmant sa place de premier exportateur mondial de biens manufacturés.

\textbf{3. Marginalité de l'Afrique (hors Tanger Med) :}
\begin{itemize}
    \item \textbf{Infrastructures portuaires :} ports peu profonds (tirant d'eau insuffisant pour les porte-conteneurs \emph{Post-Panamax} / \emph{Ultra Large}), congestionnés (Douala : \textasciitilde 0,5 million EVP/an, soit \SI{1}{\%} de Shanghai), arrière-pays mal déservis (routes dégradées, voies ferrées insuffisantes ou inexistantes), coûts de manutention et délais de dédouanement élevés.
    \item \textbf{Faible volume d'échanges conteneurisés :} l'Afrique ne pèse qu'environ \SI{2}{\%} à \SI{3}{\%} du commerce mondial de marchandises ; elle exporte majoritairement des vracs liquides (pétrole) et solides (minerais, bois en grumes) qui n'utilisent pas de conteneurs.
    \item \textbf{Position dans les CVM :} l'Afrique est positionnée en \emph{amont} des chaînes de valeur mondiales (fournisseur de matières premières) et importe les produits manufacturés conteneurisés (machines, véhicules, pharmaceutique, électronique). Elle n'accueille pas les grandes plateformes de transbordement (hubs) mondiales, sauf Tanger Med (Maroc), connecté aux flux Europe-Asie-Amériques.
\end{itemize}
\textbf{Point de vigilance :} bien distinguer trafic \emph{conteneurisé} (mesuré en EVP, indicateur de l'intégration aux CVM manufacturières) et trafic global en tonnage (vracs) ; ne pas conclure que l'Afrique est « hors de la mondialisation » : elle y est insérée, mais de façon \emph{subordonnée} (spécialisation primaire, dépendance aux importations manufacturées).
\end{corrige}

\begin{corrige}[Exercice 6 — Étude de cas : MTN Cameroon]
\textbf{1.} Le groupe mère, MTN Group, est de nationalité \cle{sud-africaine} (siège social à Johannesburg, coté à la Bourse de Johannesburg). Il opère dans le secteur des \cle{services} (télécommunications : téléphonie mobile voix/data, Internet fixe/mobile, services financiers numériques \emph{Mobile Money}).

\textbf{2. Stratégies d'implantation locale (trois exemples concrets) :}
\begin{itemize}
    \item \textbf{Recrutement et formation locaux :} l'immense majorité des \textasciitilde 1\,500 salariés directs et des cadres est camerounaise ; partenariats avec l'ENSP (Yaoundé), l'ENSET (Douala), l'Université de Ngaoundéré pour stages et insertion.
    \item \textbf{Adaptation de l'offre au marché (innovation frugale) :} lancement pionnier du \emph{Mobile Money} (MoMo, 2010) répondant à la faible bancarisation (\textasciitilde 15\,\% de la population) ; forfaits voix/data à très petits montants (quotidiens/hebdomadaires) accessibles aux faibles revenus ; distribution par un réseau dense de revendeurs informels (jeunes, femmes).
    \item \textbf{Investissement en infrastructure et RSE :} déploiement de milliers de sites 3G/4G (et tests 5G) et de fibre optique (backbone national) ; actions RSE via la Fondation MTN (santé : dons aux hôpitaux ; éducation : laboratoires informatiques dans les lycées ; mécénat culturel).
\end{itemize}

\textbf{3. Intégration du Cameroun aux flux mondiaux par MTN :}
\begin{itemize}
    \item \textbf{Flux financiers :} chiffre d'affaires \textasciitilde 400 milliards FCFA/an ; bénéfices en partie rapatriés vers l'Afrique du Sud (dividendes, redevances de marque, frais de gestion) ; \emph{Mobile Money} connecte les transferts de la diaspora (Western Union, MoneyGram, Wave, Ria partenaires) aux ménages camerounais (rural/urbain) en temps réel.
    \item \textbf{Flux technologiques et informationnels :} MTN raccorde le Cameroun aux câbles sous-marins internationaux (SAT-3/WASC, ACE, WACS) dont elle est actionnaire/consortiée ; déploie la 4G/5G avec des équipementiers mondiaux (Huawei Chine, Ericsson Suède, Nokia Finlande) ; \textasciitilde 18 millions d'abonnés (part de marché \textasciitilde 50\,\%) accèdent à l'information mondiale, aux réseaux sociaux, au e-commerce, à l'e-gouvernement en temps réel.
\end{itemize}
\textbf{Point de vigilance :} une FMN de services crée des emplois et des infrastructures, mais rapatrie des dividendes (sortie nette de devises), dépend d'équipements et licences importés (dépendance technologique) et peut pratiquer l'optimisation fiscale : nuancer le bilan (apports \emph{et} limites).
\end{corrige}

\begin{corrige}[Exercice 7 — État et diaspora : Le corridor Kribi -- Mbalam]
\textbf{1. État acteur d'intégration :} l'État camerounais décide, planifie et co-finance (avec des partenaires financiers, ex. Eximbank of China, fonds propres, PPP) le port en eau profonde de Kribi (PAK, mis en service 2014, phase 2 en cours), crée l'entité gestionnaire (Port Autonome de Kribi) et signe des \cle{conventions de concession} (BOT -- Build Operate Transfer) avec des opérateurs portuaires internationaux pour le terminal à conteneurs (KCT) et le terminal minéralier. Il aménage ainsi les conditions d'accueil des capitaux étrangers (infrastructures, cadre juridique, guichet unique) et branche le pays (et l'hinterland tchadien/centrafricain via le pipeline et la route/rail projetés) sur les routes maritimes mondiales : l'État \emph{organise} la mondialisation au lieu de la subir passivement.

\textbf{2. Rôle réel de la diaspora :} la diaspora ne finance \textbf{pas} directement ce type de mégaprojet d'infrastructure lourde (coûts de plusieurs centaines de milliards FCFA, réservés aux États, banques de développement, FMN minières/portuaires). Son rôle économique est ailleurs : \cle{envois de fonds} aux ménages (\textasciitilde 300-350 milliards FCFA/an : consommation courante, santé, éducation, cérémonies), investissements dans l'\cle{immobilier} (construction villas, appartements à Yaoundé/Douala) et les \cle{PME} (commerce, transport, agro-pastoral), transferts de compétences (expertise médicale, ingénierie, enseignement, diaspora numérique) et lobbying politique en faveur du pays.

\textbf{3. Mondialisation des services portuaires :} le consortium KCT (Kribi Conteneur Terminal) associe ICTSI (Philippines, 1er opérateur portuaire privé mondial), des partenaires liés à Maersk (Danemark, 2e armateur mondial) et Bolloré Logistics / Africa Global Logistics (France, logistique intégrée) : la gestion opérationnelle d'un port stratégique camerounais est assurée par des FMN de trois continents. Cela montre que les \cle{services} (logistique portuaire, manutention, transit) sont eux-mêmes mondialisés et que les savoir-faire managériaux et techniques circulent à l'échelle planétaire.

\textbf{Point de vigilance :} ne pas attribuer à la diaspora un rôle de financeur d'infrastructures lourdes (ports, barrages, autoroutes) ; c'est l'erreur classique du sujet de dissertation/étude de documents. La diaspora finance le \emph{bas de bilan} des ménages et des PME, pas le \emph{haut de bilan} de l'État.
\end{corrige}

\begin{corrige}[Exercice 8 — Croquis commenté : L'Afrique dans les flux de la mondialisation]
\textbf{Méthode de construction (barème indicatif sur 10 : titre 1, légende organisée 2, figurés respectés 3, contenu géographique 3, orientation/échelle 1).}
\begin{enumerate}
    \item \textbf{Fond de carte :} esquisse simplifiée des continents (blocs polygonaux) ; placer la Triade (Amérique du Nord, Europe occidentale, Asie de l'Est) en haut/hémisphère Nord, l'Afrique au centre, les pays émergents (Chine, Inde, Brésil, Afrique du Sud) repérés.
    \item \textbf{Figurés linéaires (flux) -- épaisseur proportionnelle :}
    \begin{itemize}
        \item Flèches \emph{très épaisses}, bidirectionnelles, entre les trois pôles de la Triade (flux manufacturés, financiers, technologiques intenses).
        \item Flèches \emph{moyennes} sortantes d'Afrique vers Europe, Chine, USA (matières premières : pétrole, minerais, bois, cacao, coton).
        \item Flèches \emph{moyennes} entrantes vers Afrique depuis Triade et Chine (produits manufacturés, machines, céréales, médicaments, véhicules).
        \item Flèches \emph{pointillées} entrantes vers Afrique (envois de fonds de la diaspora depuis Europe/USA/Monde arabe).
        \item Flèches IDE (investissements) : Chine/émergents $\rightarrow$ Afrique (infrastructures, mines, zones industrielles) ; Triade $\rightarrow$ Afrique (pétrole, services, télécoms).
        \item Axe routier/ferroviaire Douala -- Yaoundé -- Ngaoundéré -- N'Djamena (Tchad) et pipeline Tchad (Doba) -- Kribi en traits fins nommés.
    \end{itemize}
    \item \textbf{Figurés ponctuels :} carrés rouges = pôles de commandement mondial (New York, Londres, Tokyo/Shanghai, Paris, Francfort) ; cercles bleus = grands ports mondiaux (Shanghai, Singapour, Rotterdam, Los Angeles) et ports africains (Douala, Kribi, Tanger Med, Durban, Lagos, Mombasa) ; étoile jaune = capitales politiques (Yaoundé, Pretoria, Abuja, Le Caire, Addis-Abeba).
    \item \textbf{Figurés de surface :} teinte foncée pour la Triade (cœur), teinte intermédiaire pour les émergents (BRICS+), trame claire pour l'Afrique (périphérie active mais subordonnée).
    \item \textbf{Finitions obligatoires :} titre précis (\og L'Afrique dans les flux de la mondialisation : une insertion subordonnée \fg), légende structurée en 3 rubriques (Acteurs / Flux / Territoires), flèche du Nord, échelle indicative (ex. 1 cm = 5000 km), sources (CNUCED, OMC, FMI, Banque mondiale).
\end{enumerate}
\textbf{Points de vigilance :} l'épaisseur des flèches \emph{doit} être proportionnelle à l'importance relative des flux (visuellement lisible) ; toute flèche non définie dans la légende est sanctionnée ; ne pas surcharger (5 à 6 types de figurés maximum) ; respecter la sémiologie graphique (surface pour territoires, linéaire pour flux, ponctuel pour nœuds).
\end{corrige}

\begin{corrige}[Exercice 9 — Sujet type Bac : Dissertation — « Les firmes multinationales sont-elles les principaux acteurs de la mondialisation ? »]
\textbf{1. Analyse du sujet.} Mots-clés : \og principaux \fg{} (invite à \emph{hiérarchiser}, nuancer, discuter), \og acteurs \fg{} (qui fait la mondialisation ? agents, stratégies), \og mondialisation \fg{} (processus contemporain, depuis années 1980, à l'échelle planétaire, multidimensionnel). Le sujet appelle une réponse \emph{nuancée} : les FMN sont centrales mais pas exclusives.

\textbf{2. Problématique :} Si les firmes multinationales organisent l'essentiel des flux productifs, financiers et technologiques mondiaux, dans quelle mesure d'autres acteurs -- États, organisations internationales, diasporas, ONG, réseaux de villes -- concourent-ils, l'encadrent-ils, voire le contestent-ils, dans la fabrique quotidienne de la mondialisation ?

\textbf{3. Plan détaillé (dialectique, 3 parties équilibrées) :}
\begin{itemize}
    \item \textbf{I. Les FMN, acteurs centraux et moteurs de la mondialisation économique.}
    \begin{itemize}
        \item[a)] Puissance économique sans précédent (chiffre d'affaires > PIB de nombreux États) et organisation des \cle{chaînes de valeur mondiales} (CVM) : elles décident où produire, quoi, comment, pour quel marché (ex. : TotalEnergies, Apple, Samsung ; au Cameroun : MTN, Orange, LafargeHolcim, Dangote Cement, SOCAPALM, PHP, TotalEnergies Marketing).
        \item[b)] Vecteurs principaux de diffusion des capitaux (\cle{IDE} : stock mondial \textasciitilde 40\,000 Mds \$), des technologies (brevets, R\&D), des normes de gestion (ISO, RSE) et de consommation (marques globales).
    \end{itemize}
    \item \textbf{II. Mais d'autres acteurs structurent, encadrent et orientent la mondialisation.}
    \begin{itemize}
        \item[a)] Les \cle{États} restent décisifs : politiques d'attractivité (loi n°2013/004, APIPA, zones franches, port de Kribi au Cameroun), puissance régulatrice (USA, Chine -- BRI/Nouvelles Routes de la Soie, UE), souveraineté monétaire/fiscale ; les \cle{organisations internationales} (OMC -- règles commerce, FMI/Banque mondiale -- conditionnalités, ONU -- ODD) fixent le cadre normatif mondial.
        \item[b)] Les \cle{acteurs non étatiques} : diasporas (envois de fonds \textasciitilde 300-350 Mds FCFA/an vers le Cameroun, investissements, lobbying), ONG et associations (mondialisation par le bas : normes environnementales, droits humains, commerce équitable, santé mondiale), réseaux de villes mondiales (commandement financier, innovation).
    \end{itemize}
    \item \textbf{III. Une mondialisation coproduite, inégale et contestée.}
    \begin{itemize}
        \item[a)] Interdépendance structurée : les FMN ont besoin des États (infrastructures, sécurité juridique, main-d'œuvre formée, stabilité) et les États ont besoin des FMN (capitaux, emplois, technologies, recettes fiscales, accès aux marchés) : relation \emph{partenariale mais asymétrique}.
        \item[b)] Limites et contestations : rapatriement massif des bénéfices (érosion base fiscale, optimisation), effets d'enclave (peu de liens locaux), dépendance technologique, altermondialisme (critique OMC, paradis fiscaux), relocalisations/réshoring (sécurisation chaînes d'approvisionnement post-Covid) : la mondialisation est un processus \emph{hiérarchisé} (Nord/Sud, centre/périphérie) et perfectible, non un destin figé.
    \end{itemize}
\end{itemize}

\textbf{4. Introduction rédigée (modèle) :}
\og En 2014, le port en eau profonde de Kribi, construit avec des capitaux chinois (Eximbank) et géré par un consortium associant un opérateur philippin (ICTSI), un armateur danois (Maersk) et un logisticien français (Bolloré), entre en service. Un seul équipement camerounais y associe ainsi des entreprises de trois continents. Cette scène illustre la mondialisation, processus d'intensification et d'accélération des flux de biens, de capitaux, d'informations et de personnes qui met les territoires du monde en interdépendance. Au cœur de ce processus, les firmes multinationales (FMN), entreprises détenant des actifs productifs dans au moins deux pays, apparaissent souvent comme les acteurs dominants. Mais sont-elles réellement les \emph{principaux} acteurs de la mondialisation ? Nous verrons d'abord que les FMN occupent une place centrale dans l'organisation des flux productifs et financiers mondiaux, puis que les États, les organisations internationales et les acteurs non étatiques structurent tout autant ce processus, enfin que la mondialisation résulte d'une coproduction inégale et contestée entre ces acteurs. \fg

\textbf{Barème indicatif (sur 20) :} analyse du sujet et problématique 4 pts ; solidité du plan et pertinence des arguments 8 pts ; exemples précis et variés (Cameroun + monde) 4 pts ; qualité rédactionnelle (vocabulaire, connecteurs, structure) 4 pts.
\textbf{Points de vigilance :} définir \emph{tous} les termes du sujet dans l'introduction ; éviter le catalogue d'exemples sans argumentation ; chaque partie doit contenir \emph{au moins} un exemple camerounais et un exemple mondial ; soigner la conclusion (ouverture sur ZLECAf, transition écologique, souveraineté numérique).
\end{corrige}

\begin{corrige}[Exercice 10 — Sujet type Bac : Étude de documents (Routes maritimes + Parts commerce mondial)]
\textbf{1. Géographie des routes maritimes :} les grandes routes relient les pôles de la Triade : Pacifique Nord (Asie de l'Est -- côte ouest USA), Atlantique Nord (Europe -- côte est USA), Europe -- Asie via Suez et Malacca, et route Amérique du Sud -- Chine (soja, minerais). Malacca et Suez sont des \cle{goulets d'étranglement} (passages étroits, quasi obligés entre deux mers) : les détourner (Cap de Bonne-Espérance, détroit de Magellan) allongerait fortement temps (10-15 jours) et coûts (carburant, équipage, assurance) ; leur blocage (ex. \emph{Ever Given}, Suez mars 2021) perturbe instantanément tout le commerce mondial (hausse fret, pénuries).

\textbf{2. Part de la Triade dans les exportations mondiales (Document 2) :}
\[ 38\,\% (\text{UE}) + 36\,\% (\text{Asie de l'Est, dont Chine}) + 14\,\% (\text{Amérique du Nord}) = \SI{88}{\%}. \]
(En retirant la Chine \textasciitilde \SI{14}{\%} de l'Asie, la Triade « classique » UE+USA+Japon pèse environ \SI{74}{\%} ; avec la Chine, les trois grandes façades maritimes concentrent près de 9 dixièmes des exportations mondiales.) L'Afrique, avec \SI{2,3}{\%}, reste marginale malgré son poids démographique (\textasciitilde 18\,\% population mondiale) et ses ressources : elle exporte peu et surtout des produits bruts à faible valeur ajoutée.

\textbf{3. Synthèse — Un commerce international hiérarchisé et régionalisé :} le commerce international reste dominé par trois pôles (Asie de l'Est, Europe, Amérique du Nord) qui concentrent richesses, FMN, grands ports, places financières et technologies : d'où une \cle{hiérarchie} Nord/Sud persistante. Par ailleurs, l'essentiel des échanges se fait \emph{à l'intérieur} de blocs régionaux (UE : \textasciitilde 60\,\% échanges intra-zone ; USMCA ; ASEAN) : la proximité, les accords préférentiels et les chaînes de valeur régionales favorisent la \cle{régionalisation} (mondialisation « par blocs »). L'Afrique tente de rattraper ce retard avec la \cle{ZLECAf} (Zone de Libre-Échange Continentale Africaine, 2021), qui vise à créer un marché unique continental de 1,3 milliard de consommateurs. La mondialisation n'efface donc ni les inégalités de développement ni les territoires : elle les recompose en réseaux inégaux.

\textbf{Barème indicatif (sur 20) :} Q1 (description + explication goulets) : 6 pts ; Q2 (calcul exact écrit + commentaire hiérarchie) : 6 pts ; Q3 (mobilisation des deux documents + connaissances personnelles, notions de hiérarchie et régionalisation, vocabulaire précis) : 8 pts.
\textbf{Points de vigilance :} citer explicitement les documents (« le document 1 montre que... », « le document 2 indique que... ») ; le calcul doit être écrit explicitement ; ne pas confondre part du commerce mondial (relative) et volume absolu (croissance africaine rapide mais partie faible) ; définir \emph{goulet d'étranglement} et \emph{régionalisation}.
\end{corrige}

\begin{corrige}[Exercice 11 — Sujet type Bac : TP Noté « Le Cameroun dans la mondialisation »]
\textbf{1. Diagramme en barres groupées (6 pts).} Regroupement des catégories en milliards FCFA (données INS/Douanes 2023 arrondies) :
\begin{center}
\begin{tabularx}{\linewidth}{|l|X|X|}\hline
\textbf{Catégorie} & \textbf{Exportations (Mds FCFA)} & \textbf{Importations (Mds FCFA)} \\ \hline
Hydrocarbures & 1\,800 (pétrole brut) & 800 (produits raffinés) \\ \hline
Produits agricoles et bois & $450+300+120 = 870$ (cacao, bois, coton) & 400 (céréales : riz, blé, maïs) \\ \hline
Produits manufacturés & 100 (aluminium, première transformation) & $600+300+250 = 1\,150$ (machines, véhicules, pharma, chimie) \\ \hline
Autres & 230 & 1\,150 \\ \hline
\textbf{Total} & \textbf{3\,000} & \textbf{3\,500} \\ \hline
\end{tabularx}
\end{center}
\textbf{Construction attendue :} deux barres juxtaposées par catégorie (couleur distincte : ex. bleu export, rouge import), échelle adaptée (ex. 1 cm = 200 Mds FCFA), axes gradués et légendés (abscisses : catégories ; ordonnées : milliards FCFA), titre précis (\og Structure comparée des exportations et importations du Cameroun, 2023 \fg), source (INS / Douanes, données simplifiées).

\textbf{2. Analyse de la structure commerciale (6 pts).} La balance commerciale est structurellement déficitaire :
\[ \text{Solde} = 3\,000 - 3\,500 = -\SI{500}{\text{ milliards FCFA}} \quad ; \quad \text{Taux de couverture} = \frac{3\,000}{3\,500}\times100 \approx \SI{85,7}{\%}. \]
Le Cameroun est spécialisé dans l'exportation de \cle{matières premières brutes} (pétrole \SI{60}{\%}, cacao, bois, coton) dont la demande est peu élastique aux prix et dont les cours fluctuent fortement (volatilité) : d'où une \cle{vulnérabilité} extérieure majeure (chocs asymétriques). Il importe massivement des produits manufacturés (machines, véhicules, pharmaceutique) et même des produits pétroliers raffinés (essence, gasoil), faute de capacité de raffinage locale suffisante (SONARA limitée, projet de nouvelle raffinerie en attente). Les \cle{termes de l'échange} (rapport entre l'indice des prix des exportations et l'indice des prix des importations) tendent à se dégrader structurellement pour les pays exportateurs de produits primaires (thèse Prebisch-Singer) : il faut exporter toujours plus de volume de matières premières pour acheter la même quantité de produits manufacturés.

\textbf{3. Acteurs et intégration (8 pts).}
\begin{itemize}
    \item \textbf{FMN :} TotalEnergies et Perenco (hydrocarbures, \textasciitilde 80\,\% production), MTN et Orange (télécoms, \textasciitilde 95\,\% parts de marché), Cimencam/LafargeHolcim et Dangote Cement (ciment, oligopole), SOCAPALM/PHP (palmier à huile/hévéa), Guinness/SABC (brasserie). Apports : capitaux (stock d'IDE \textasciitilde \SI{3500}{\text{ milliards FCFA}}), technologies, emplois directs/indirects (\textasciitilde 100\,000), devises, formation. Limites : rapatriement bénéfices (\textasciitilde 1000-1500 Mds FCFA/an), \cle{enclaves} peu connectées au tissu local (sous-traitance faible), dépendance technologique, optimisation fiscale (prix de transfert).
    \item \textbf{État :} politique d'attractivité (loi n°2013/004, APIPA, zones franches), infrastructures structurantes (port Kribi, pipeline Tchad-Kribi, autoroute Yaoundé-Douala, barrages Lom Pangar/Mekin/Nachtigal, câbles SAT-3/ACE/WACS), adhésion ZLECAf, facilitation douanière (guichet unique GUCE). Limites : endettement extérieur croissant (notamment Chine \textasciitilde 3000 Mds FCFA), gouvernance (corruption, lenteurs administratives), délais chantiers, insuffisance formation technique adaptée.
    \item \textbf{Diaspora :} \textasciitilde \SI{320}{\text{ milliards FCFA}} de transferts formels en 2023 (Banque mondiale), 2e source de devises après hydrocarbures. Soutien consommation, santé, éducation, immobilier, PME locales. Limite : ces fonds financent peu l'industrie exportatrice à forte valeur ajoutée ; coûts d'envoi encore élevés (\textasciitilde 7-8\,\% frais).
\end{itemize}
\textbf{Conclusion du TP :} l'insertion du Cameroun dans la mondialisation est réelle (flux financiers, humains, numériques, portuaires) mais \cle{subordonnée} : exportateur de brut, importateur de transformé, dépendant des décisions d'investissement des FMN et des conditions de prêt des créanciers bilatéraux/multilatéraux. La transformation locale (raffinerie, usines cacao/bois/coton, ZLECAf) et la montée en gamme des ressources humaines sont les pistes prioritaires pour une intégration plus souveraine.

\textbf{Barème indicatif :} diagramme 6 pts (titre, axes, échelle, exactitude valeurs, propreté, source) ; analyse structurelle 6 pts (calcul déficit/taux couverture écrit, spécialisation, termes de l'échange définis et appliqués) ; acteurs 8 pts (triptyque FMN/État/Diaspora, apports \emph{et} limites pour chacun, exemples précis, vocabulaire de programme).
\textbf{Points de vigilance :} écrire explicitement le calcul du taux de couverture ; définir les termes de l'échange avant de les utiliser ; un TP noté exige une copie soignée : légende complète du graphique, traits à la règle, écriture lisible, respect de l'échelle.
\end{corrige}

\newpage

\coursec{Fiche bilan}

\begin{popfigure}
\centering
\begin{tikzpicture}[
  mindmap,
  concept color=popblue,
  level 1 concept/.append style={level distance=4.5cm, sibling angle=60},
  level 2 concept/.append style={level distance=3cm, sibling angle=45},
  every node/.style={font=\small},
  text=popdark,
  branch1/.style={concept color=poporange, grow=-90},
  branch2/.style={concept color=popgreen, grow=-30},
  branch3/.style={concept color=poppurple, grow=30},
  branch4/.style={concept color=poppink, grow=90},
  branch5/.style={concept color=popgold, grow=150},
  branch6/.style={concept color=popblue, grow=210},
]
\node[concept, minimum size=2.8cm, font=\bfseries\large] {Facteurs de la mondialisation}
  child[branch1] { node[concept, minimum size=2.2cm] {Transports \& Télécoms}
    child { node[concept, minimum size=1.6cm] {Conteneurisation} }
    child { node[concept, minimum size=1.6cm] {Câbles sous-marins (WACS, ACE)} }
    child { node[concept, minimum size=1.6cm] {Internet, mobile money} }
  }
  child[branch2] { node[concept, minimum size=2.2cm] {Firmes Multinationales (FMN)}
    child { node[concept, minimum size=1.6cm] {IDE \& filiales (ex. Bolloré, MTN, Dangote)} }
    child { node[concept, minimum size=1.6cm] {Chaînes de valeur mondiales} }
    child { node[concept, minimum size=1.6cm] {Transfert de technologie} }
  }
  child[branch3] { node[concept, minimum size=2.2cm] {État \& Diaspora}
    child { node[concept, minimum size=1.6cm] {Politiques d'ouverture (ZLECAf)} }
    child { node[concept, minimum size=1.6cm] {Transferts de fonds (> 300 Mds FCFA/an)} }
    child { node[concept, minimum size=1.6cm] {Compétences \& investissements} }
  }
  child[branch4] { node[concept, minimum size=2.2cm] {Associations \& ONG}
    child { node[concept, minimum size=1.6cm] {Plaidoyer (climat, droits humains)} }
    child { node[concept, minimum size=1.6cm] {Services (santé, éducation)} }
    child { node[concept, minimum size=1.6cm] {Réseaux transnationaux} }
  }
  child[branch5] { node[concept, minimum size=2.2cm] {Commerce international}
    child { node[concept, minimum size=1.6cm] {OMC, APE, ZLECAf} }
    child { node[concept, minimum size=1.6cm] {Flux : matières 1ères vs manufacturés} }
    child { node[concept, minimum size=1.6cm] {Port de Douala-Kribi, corridor CEMAC} }
  }
  child[branch6] { node[concept, minimum size=2.2cm] {Finance \& Institutions}
    child { node[concept, minimum size=1.6cm] {FMI, BM, BAD, BEAC} }
    child { node[concept, minimum size=1.6cm] {Marchés financiers, change} }
    child { node[concept, minimum size=1.6cm] {Conditionnalités, PPTE} }
  };
\end{tikzpicture}
\legende{Carte mentale synthétique des 6 moteurs majeurs de la mondialisation (adaptée au contexte camerounais).}
\end{popfigure}

\begin{bilan}
\courssub{Définitions clés à connaître par cœur}
\begin{itemize}
  \item \cle{Mondialisation} : Processus d'intensification des flux (biens, capitaux, informations, personnes) à l'échelle planétaire, réduisant les distances et les frontières.
  \item \cle{Firme Multinationale (FMN)} : Entreprise qui exerce un contrôle sur des actifs productifs (filiales) dans au moins deux pays (ex. \textbf{Bolloré Logistics}, \textbf{MTN}, \textbf{Société des Brasseries du Cameroun}).
  \item \cle{IDE (Investissement Direct à l'Étranger)} : Investissement impliquant une relation durable et un contrôle significatif (généralement $\ge 10\%$ du capital) d'une entreprise dans un pays par une entité d'un autre pays.
  \item \cle{Diaspora} : Ensemble des migrants originaires d'un même pays résidant à l'étranger, maintenant des liens forts (transferts, réseaux) avec le pays d'origine.
  \item \cle{ONG (Organisation Non Gouvernementale)} : Association à but non lucratif, indépendante des États, agissant à l'échelle internationale (humanitaire, développement, environnement).
  \item \cle{ZLECAf (Zone de Libre-Échange Continentale Africaine)} : Accord de l'Union Africaine (opérationnel depuis 2021) créant un marché unique africain (1,3 Md hab., PIB $\approx$ 3 400 Mds \$).
\end{itemize}

\courssub{Rôles et mécanismes des 5 facteurs (programme officiel)}
\begin{center}
\begin{tabularx}{\linewidth}{|l|X|X|}\hline
\rowcolor{popblueL}
\textbf{Facteur} & \textbf{Mécanisme principal} & \textbf{Impact / Exemple Cameroun} \\ \hline
\textbf{1. Transports \& Télécoms} & Réduction coûts/délais (conteneur, câbles WACS/ACE, 4G/5G, Mobile Money) & Port de \textbf{Kribi} (eau profonde), corridor Douala-Ndjamena, inclusion financière (Orange Money, MTN MoMo). \\ \hline
\textbf{2. Firmes Multinationales} & IDE, sous-traitance, chaînes de valeur mondiales (CVM), transfert techno. & \textbf{Dangote Cement} (Figuil), \textbf{Neo Industry} (transformation cacao), \textbf{Azure Energy} (solaire). \\ \hline
\textbf{3. État \& Diaspora} & Politique commerciale (tarifs, ZLECAf), transferts financiers, compétences. & Transferts diaspora $\approx$ \textbf{300--400 Mds FCFA/an} (1er flux entrant > APD), loi incitative 2013. \\ \hline
\textbf{4. Associations \& ONG} & Plaidoyer normes mondiales, services de base, veille citoyenne. & \textbf{Greenpeace} (forêt), \textbf{Plan International} (éducation), \textbf{Croix-Rouge} (crises NOSO/Extrême-Nord). \\ \hline
\textbf{5. Commerce international} & Règles OMC, accords régionaux (APE UE-Cameroun, ZLECAf), spécialisation. & Export : bois, cacao, pétrole, coton, aluminium. Import : machines, hydrocarbures, céréales, médicaments. \\ \hline
\end{tabularx}
\end{center}

\courssub{Méthodes express (Bac)}
\begin{enumerate}
  \item \textbf{Commenter une carte de flux / croquis de mondialisation} : 1) Titre, échelle, légende. 2) Identifier les pôles (ports, capitales, zones industrielles). 3) Décrire les flux (épaisseur = volume, sens = direction). 4) Analyser la place du Cameroun (insertion périphérique vs hub sous-régional). 5) Conclure sur dépendance/émergence.
  \item \textbf{Construire un croquis « Le Cameroun dans la mondialisation »} : Fond de carte (contour, 10 régions, villes Douala/Yaoundé/Kribi/Maroua/Garoua). Symboles : ports (ancres), aéroports (avions), axes routiers/fer (traits), flux export/import (flèches colorées), FMN (logos/carrés), câbles (traits pointillés mer). Légende organisée par thèmes.
  \item \textbf{Étude de cas : Une FMN au Cameroun} : Identifier (nom, siège, secteur), Localiser (usines, bureaux), Analyser stratégie (recherche matière 1ère, marché, coûts), Évaluer impacts (emplois, recettes fiscales, transfert tech, environnement, lien local).
\end{enumerate}
\end{bilan}

\begin{aretenir}[Checklist avant le Bac]
$\square$ Je sais définir \textbf{mondialisation, FMN, IDE, diaspora, ONG, ZLECAf} avec un exemple camerounais précis.\par
$\square$ Je sais expliquer \textbf{comment les transports (conteneur, port Kribi) et télécoms (câbles, Mobile Money) réduisent les distances}.\par
$\square$ Je sais analyser le rôle des \textbf{FMN} (Bolloré, MTN, Dangote, Neo Industry) : IDE, création d'emplois, chaîne de valeur, rapatriement bénéfices.\par
$\square$ Je connais le poids de la \textbf{diaspora} (transferts $\approx$ 300--400 Mds FCFA/an, 1ère source de devises) et les politiques de l'État (loi 2013, ZLECAf).\par
$\square$ Je sais citer des \textbf{ONG} actives au Cameroun et leurs domaines (environnement, santé, éducation, droits humains).\par
$\square$ Je maîtrise la structure du \textbf{commerce extérieur} : partenaires (UE, Chine, CEMAC), produits (export matières 1ères / import manufacturés), balance souvent déficitaire.\par
$\square$ Je sais \textbf{légender et analyser une carte de flux} (ports, corridors, câbles sous-marins, IDE, flux migratoires).\par
$\square$ Je sais \textbf{construire un croquis synthétique} « Le Cameroun dans la mondialisation » avec une légende hiérarchisée (3 à 4 items).\par
$\square$ Je sais discuter la \textbf{place du Cameroun} : « pays charnière » Afrique centrale, hub logistique (Kribi), mais insertion inégale (dépendance matières 1ères, déficit infrastructurel).\par
$\square$ Je connais les \textbf{enjeux de la ZLECAf} pour le Cameroun : élargir le marché, transformer localement, réduire la dépendance extra-africaine.\par
$\square$ Je distingue \textbf{mondialisation subie (dépendance) et mondialisation choisie (stratégie émergence, ZLECAf, transformation locale)}.
\end{aretenir}

\findecours

\end{document}
